\documentclass[11pt]{article}
\usepackage[margin=2.6cm]{geometry}
\usepackage{booktabs}
\usepackage{graphicx}
\usepackage{tikz}
\usetikzlibrary{positioning,arrows.meta,fit,backgrounds,calc}
\usepackage{siunitx}
\sisetup{group-digits=integer}
\usepackage{amsmath}
\usepackage{hyperref}
\usepackage{seqsplit}
\graphicspath{{figures/}}
\setlength{\emergencystretch}{3em}
\newcommand{\runid}[1]{\texttt{\small\seqsplit{#1}}}
\newcommand{\dsid}[1]{\texttt{\small #1}}
\newcommand{\best}[1]{\textbf{#1}}
\newcommand{\Rp}{\ensuremath{R'}}
\newcommand{\Cp}{\ensuremath{C'}}

\title{Next-Track Session Recommendation: Which Trained Models Win, and Why}
\author{}
\date{\today}

\begin{document}
\maketitle

\begin{abstract}
This document compares the \emph{trained} models built for the next-track task
in this lab and argues why the successful ones succeed; the markdown reports in
\texttt{experiments/} are the primary record and this synthesis transcribes
rather than recomputes their numbers. The task is next-item \emph{retrieval}:
given a listening \emph{session} (an ordered sequence of tracks), rank the
\num{19402}-item catalog and recover the next \emph{distinct} track the user
actually played. The board is scored on one leak-free chronological split
(\dsid{seq-20260715-131139}: \num{7154} sessions / \num{19402} items,
\num{5723} train / \num{1431} test) and is denominated in \textbf{Recall@10},
the record's comparability currency; the server's \emph{declared} primary metric
is \textbf{holisticness@10}, a four-factor list-quality composite that has had
\emph{three} incompatible code generations and that orders the models almost
inversely, and this document keeps both boards distinct rather than silently
choosing one. Six trained predictor families are in scope
(Section~\ref{sec:zoo}); pure table lookups over the dataset appear only as
reference bars and as blend \emph{legs}, never as \emph{compared} models. The
one fact that explains most of the results is that \SI{55}{\percent} of test
targets are \emph{cold} --- never seen in any training session --- which caps
every seen-items-only mechanism at Recall@10 \num{0.4493} and makes retrieval in
a \emph{learned content geometry}, not per-item memory, the only channel to the
majority of the test set. On that reading the record has one champion --- the
promoted \Rp+M+\Cp{} blend at Recall@10 \num{0.2117} (\num{303} of \num{1431}),
paired $\Delta$ over the unprojected three-leg crown $+\num{0.0259}$
$[+0.012,+0.041]$, CI \emph{entirely above zero}, seed-robust across three seeds
and hardened on a second, colder split at $+\num{0.0426}$
$[+\num{0.0280},+\num{0.0573}]$ --- a single-model bar of \num{0.12299}
(\num{176} of \num{1431}) that no topology change has beaten \emph{with a CI
above zero}, five refuted design axes whose mechanisms are now \emph{measured}
rather than guessed (dual-tower view fusion, a ReLU pre-encoder, an MLP fusion
head, four learned combiners, and depth/width/direction), and an eval-time
re-ranking policy --- maximal marginal relevance, MMR --- that is the first
honest lever on the crown and from which \textbf{nothing has been promoted}. The
honest frame for the nulls is statistical: at $n=\num{1431}$ the paired
half-width is ${\sim}\num{0.0124}$ on Recall@10 against a whole-family spread of
\num{0.028}, so large harms were reliably detected while \emph{moderate wins
were never detectable at all}.
\end{abstract}

\section{Introduction: one split, one number, and why the cold half decides everything}
\label{sec:intro}

\emph{Rationale.} Every model in this document solves the same problem. A
listening \emph{session} is an ordered sequence of tracks; the model reads a
prefix and must produce a ranking of the entire \num{19402}-item catalog such
that the track the user actually played next lands near the top. Two separate
conventions shape that, and the phrase ``next-distinct'' is easy to misread as
covering both. First, the \textbf{target} is the next \emph{distinct} track: an
immediate repeat of something already in the prefix is skipped over when the
label is constructed, so the task is not repeat detection in disguise --- this
is the operation the name refers to. Second, and independently, the
\textbf{candidate set} excludes every item already present in the prefix, which
is the parity convention the Markov leg is built on; a model is therefore never
credited for re-proposing a track the session already contains.

\emph{The one artifact --- and the three places it is not.} The board is scored
on \dsid{seq-20260715-131139}, split \emph{chronologically by session start} at
2024-08-24T14:28:02Z into \num{5723} train and \num{1431} test sessions, with
every item represented by a frozen \num{192}-dimensional latent
(\dsid{spotify\_tracks\_song\_pca192}). Three groups of numbers below are
deliberately \emph{not} on that artifact and are named at every use: the
\texttt{seq-ann} pooled twin and one stacker arm (\dsid{seq-20260715-030509},
PCA-128), the second-split hardening of the champion
(\dsid{seq-20260718-211238}), and the 2026-07-13/14 topology and first-blend
campaigns, which predate it (\dsid{seq-20260713-014226}, built on the earlier
\emph{AE-64} item space --- a \num{64}-dimensional autoencoder embedding of the
same track features, which a compression bake-off later replaced with PCA-192
because the linear projection retrieved better at every dimension tested). Their \emph{verdicts} are paired $\Delta$s against
a same-split anchor and do travel; their \emph{levels} do not, and must never be
read beside the canonical-split column.

\emph{The hard fact.} A test target is \textbf{cold} if its item index never
appears in \emph{any} of the \num{5723} training sessions. On this split
\num{788} of \num{1431} targets are cold --- a rate of \num{0.550664}, which
independently reproduces the manifest's declared \num{0.5507}. The consequence
is a hard structural ceiling: any mechanism that can only retrieve items it
observed during training --- a co-occurrence table, a first-order Markov bigram,
a free per-item embedding table --- is capped at Recall@10 $\num{643}/\num{1431}
= \num{0.449336}$ no matter how good it is. Everything above that line has to
come from somewhere else, and on this corpus the only ``somewhere else'' is
\emph{content geometry}: a track that never appeared in training still has a
feature vector, hence a position in the item space, hence a chance of being
retrieved.

\emph{The second hard fact.} The supervision budget is thin. Under per-step
teacher forcing the \num{5723} training sessions yield \num{73632}
(prefix, next-item) pairs against a \num{19402}-item retrieval catalog, of which
only \num{15816} items are ever observed in training; the \emph{median} observed
item appears exactly \emph{once}, the mean \num{5.017} times, and only
\num{3294} items (\SI{17}{\percent} of the catalog) reach five occurrences.
These figures were recomputed here from the dataset artifact and agree with the
\textsc{data budget} block in the header of
\texttt{predictors/seq\_embed.py}. Read together with the cold rate they explain
the shape of the entire record: every win on this board came from improving the
\emph{shared} latent geometry, and every attempt to spend parameters per item,
per leg or per extra tower failed.

\emph{Interpretation.} The rest of this document is one argument in that frame.
Section~\ref{sec:machinery} defines the machinery; Section~\ref{sec:metrics}
handles the awkward fact that this lab now has two primary metrics;
Sections~\ref{sec:zoo}--\ref{sec:board} lay out the trained models and their
scores; Sections~\ref{sec:fusion}--\ref{sec:mmr} explain the four things that
work; Section~\ref{sec:failures} explains the five that do not, with a measured
mechanism for each; Section~\ref{sec:power} reframes those nulls as upper
bounds; and Section~\ref{sec:gains} states what the record implies about where
gains on this problem actually come from.

{\rightskip=0pt plus 2em \emph{Sources: \texttt{PROJECT-FACTS.md}, and \texttt{2026-\allowbreak 07-\allowbreak 15c-\allowbreak nexttrack-\allowbreak phase2-\allowbreak model-\allowbreak sweep.md}.}\par}

\clearpage
\section{Background: the machinery, defined once}
\label{sec:machinery}

% ===========================================================================
%  topo-style.tex --- the shared TikZ style block for the next-track topology
%  diagrams AND for the evaluation-protocol diagram in the body of
%  model-comparison.tex.  Extracted from topologies.tex so that the protocol
%  figure, which appears BEFORE % ===========================================================================
%  topologies.tex --- TikZ network-topology diagrams for the next-track
%  MODEL-COMPARISON document.  Not a standalone document: \input it from the
%  body of the main .tex, which must be compiled from docs/:
%
%      % ===========================================================================
%  topologies.tex --- TikZ network-topology diagrams for the next-track
%  MODEL-COMPARISON document.  Not a standalone document: \input it from the
%  body of the main .tex, which must be compiled from docs/:
%
%      \input{figures/topologies}
%
%  (\graphicspath{{figures/}} governs \includegraphics only, NOT \input, so the
%   figures/ prefix is spelled out above.)
%
%  ------------------------------------------------------------------------
%  REQUIRED PREAMBLE LINES in the main document --- add these two lines after
%  \usepackage{graphicx}, before \begin{document}:
%
%      \usepackage{tikz}
%      \usetikzlibrary{positioning,arrows.meta,fit,backgrounds,calc}
%
%  Nothing else is required: xcolor arrives with tikz, and every colour and
%  every style used below is defined in this file.  \dsid, \Rp and \Cp are
%  \providecommand-guarded here so the file also compiles inside a bare
%  wrapper; when the main document defines them, its definitions win.
%
%  PACKAGES DELIBERATELY *NOT* REQUIRED: amssymb is absent from the house
%  preamble, so nothing below uses \mathbb.  Every picture is laid out with
%  RELATIVE positioning (right=/above=/below= of ...), never with hard
%  coordinates, so no node can collide with its neighbour if the font or the
%  wording changes; the per-figure notes and the legend are anchored to
%  `current bounding box', so they always land underneath the diagram.
%  ------------------------------------------------------------------------
%
%  LABELS DEFINED IN THIS FILE (in file order)
%      fig:topo-gru       seq-nexttrack --- the baseline; ALSO CARRIES THE LEGEND
%      fig:topo-ann       seq-ann       --- pooled feed-forward twin
%      fig:topo-dualgru   seq-dualgru   --- two towers + fusion head + pre-MLP variant
%      fig:topo-embed     seq-embed     --- learned, weight-tied item table (DESIGN ONLY)
%      fig:topo-blend     seq-blend     --- the champion, three legs
%      fig:topo-stacker   seq-stacker   --- learned XGBoost combiner
%
%  VISUAL LANGUAGE (the primary axis is the LINE, not the colour, so it
%  survives greyscale printing; stated for the reader in the caption of
%  fig:topo-gru and re-pointed-at by every other caption)
%      SOLID  border  = TRAINED component (parameters fitted by gradient descent)
%      DASHED border  = FROZEN or PARAMETER-FREE component
%      blue   = the on-record default path
%      purple dashed  = non-parametric table lookup over the dataset
%      dashed blue    = parameter-free, but operating inside a LEARNED geometry
%      green  = variant favoured by the record
%      red    = variant refuted by the record
% ===========================================================================

\input{figures/topo-style}



% ===========================================================================
% 1. seq-nexttrack --- the baseline (and the legend)
% ===========================================================================
\begin{figure}[ht]
  \centering
  \begin{tikzpicture}[topo]
    % ---- the frozen item-latent table, read TWICE ------------------------
    \node[frozen, minimum width=20mm, minimum height=12mm] (tab)
      {frozen item\\latents $Z$\\[1pt]{\scriptsize$19{,}402\times192$}};
    % ---- the prefix, in order (one column per step) ----------------------
    \node[frozen, minimum width=12mm, right=7mm of tab]  (x1) {$x_1$};
    \node[frozen, minimum width=12mm, right=4mm of x1]   (x2) {$x_2$};
    \node[tag, right=2.5mm of x2]                        (xd) {$\cdots$};
    \node[frozen, minimum width=12mm, right=2.5mm of xd] (x3) {$x_t$};
    \begin{scope}[on background layer]
      \node[grp, fit=(x1)(x3)] (pfx) {};
    \end{scope}
    \draw[data] (tab.east) -- (pfx.west);
    \node[tag, below=1.2mm of pfx] (pfxlab)
      {prefix latents --- \emph{order preserved}};
    % ---- the recurrence --------------------------------------------------
    \node[trained, minimum width=12mm, above=7mm of x1] (g1) {GRU};
    \node[trained, minimum width=12mm, above=7mm of x2] (g2) {GRU};
    \node[tag] (gd) at (g1 -| xd) {$\cdots$};
    \node[trained, minimum width=12mm, above=7mm of x3] (g3) {GRU};
    \draw[flow] (x1) -- (g1);
    \draw[flow] (x2) -- (g2);
    \draw[flow] (x3) -- (g3);
    \draw[flow] (g1) -- node[edg, above, inner sep=0.7pt] {$h_t$} (g2);
    \draw[flow] (g2) -- (gd);
    \draw[flow] (gd) -- (g3);
    % ---- head and retrieval ---------------------------------------------
    \node[trained, right=6mm of g3] (head)
      {linear head\\{\scriptsize$256\to192$}};
    \node[op, right=7mm of head] (rk)
      {rank the catalogue by $\cos(\hat z, z_v)$\\{\scriptsize prefix excluded $\to$ top-10}};
    \draw[flow] (g3) -- (head);
    \draw[flow] (head) -- node[edg, above, inner sep=1pt] {$\hat z$} (rk);
    % ---- the SAME frozen table is the retrieval target -------------------
    \coordinate (c1) at ([yshift=-3mm]tab.south |- current bounding box.south);
    \coordinate (c2) at (rk.south |- c1);
    \draw[data] (tab.south) -- (c1);
    \draw[data] (c1) -- node[tag, above]
      {the \emph{same} frozen table is the retrieval target} (c2);
    \draw[data] (c2) -- (rk.south);
    % ---- supervision note ------------------------------------------------
    \node[note] (nt) at ([yshift=-3.5mm]current bounding box.south west)
      {\textbf{Supervision --- per-step teacher forcing.} The head is shared
       across positions and must predict the next item's latent at
       \emph{every} prefix position, so one session of length $L$ yields
       $L-1$ examples: $73{,}632$ (prefix, next-item) pairs on the train
       split. The loss is masked in-batch InfoNCE at $\tau=0.07$ --- each
       step's true next latent must outscore the other steps' targets in the
       same minibatch, which supplies negatives for free.};
    % ---- LEGEND (two rows so it cannot widen the picture) ----------------
    \node[trained, swatch, anchor=north west] (l1)
      at ([yshift=-3mm]current bounding box.south west) {};
    \node[tag, right=1.1mm of l1] (t1) {trained (solid)};
    \node[frozen, swatch, right=3.4mm of t1] (l2) {};
    \node[tag, right=1.1mm of l2] (t2) {frozen / parameter-free (dashed)};
    \node[lookup, swatch, right=3.4mm of t2] (l3) {};
    \node[tag, right=1.1mm of l3] (t3) {table lookup};
    \node[learnspace, swatch, anchor=north west] (l4)
      at ([yshift=-1.6mm]l1.south west) {};
    \node[tag, right=1.1mm of l4] (t4) {parameter-free, learned space};
    \node[won, swatch, right=3.4mm of t4] (l5) {};
    \node[tag, right=1.1mm of l5] (t5) {favoured variant};
    \node[refuted, swatch, right=3.4mm of t5] (l6) {};
    \node[tag, right=1.1mm of l6] (t6) {refuted variant};
  \end{tikzpicture}
  \caption{\texttt{seq-nexttrack}, the single-model baseline every other
    topology in this document is measured against: a frozen item space, a
    recurrence over the \emph{ordered} prefix, a linear head, and cosine
    retrieval back into the \emph{same} frozen space. Nothing about the item
    geometry is learned here --- the $19{,}402\times192$ latent table (offline
    PCA-192 on \dsid{seq-20260715-131139}) is read twice, once to embed the
    prefix and once as the retrieval target, and only the GRU and its linear
    head carry parameters (\num{394944} at $h=256$). \textbf{The visual
    language used throughout this document is fixed here.} A \emph{solid}
    border marks a trained component and a \emph{dashed} border a frozen or
    parameter-free one; that distinction rides on the line style, not the
    colour, so it survives greyscale. Colour then carries the verdict ---
    blue: the on-record default path; purple dashed: a non-parametric table
    lookup over the dataset; dashed blue: parameter-free but operating inside
    a \emph{learned} geometry; green: a variant the record favours; red: a
    variant the record refutes. On the canonical split this configuration
    scores Recall@10 \num{0.12299} (\num{176} of \num{1431}), reproduced
    bit-identically eleven times.}
  \label{fig:topo-gru}
\end{figure}


% ===========================================================================
% 2. seq-ann --- the pooled feed-forward twin
% ===========================================================================
\begin{figure}[ht]
  \centering
  \begin{tikzpicture}[topo]
    \node[frozen, minimum width=20mm, minimum height=12mm] (tab)
      {frozen item\\latents $Z$\\[1pt]{\scriptsize$19{,}402\times192$}};
    \node[frozen, minimum width=12mm, right=7mm of tab]  (x1) {$x_1$};
    \node[frozen, minimum width=12mm, right=4mm of x1]   (x2) {$x_2$};
    \node[tag, right=2.5mm of x2]                        (xd) {$\cdots$};
    \node[frozen, minimum width=12mm, right=2.5mm of xd] (x3) {$x_t$};
    \begin{scope}[on background layer]
      \node[grp, fit=(x1)(x3)] (pfx) {};
    \end{scope}
    \draw[data] (tab.east) -- (pfx.west);
    % ---- the two parameter-free summaries of the same prefix -------------
    \node[op, right=8mm of pfx, yshift=10mm] (mp)
      {causal mean-pool\\$\frac{1}{t}\sum_{i\le t}x_i$\\[1pt]{\scriptsize\emph{permutation-invariant}}};
    \node[op, right=8mm of pfx, yshift=-10mm] (lt)
      {last item\\$x_t$\\[1pt]{\scriptsize the only recency cue}};
    \draw[flow] (pfx.east) -- (mp.west);
    \draw[flow] (pfx.east) -- (lt.west);
    % ---- the MLP ---------------------------------------------------------
    \node[trained, right=5mm of mp, yshift=-10mm] (mlp)
      {MLP\\$384\to256\to256\to192$\\[1pt]{\scriptsize ReLU $+$ dropout}};
    \draw[flow] (mp.east) -- (mlp.north west);
    \draw[flow] (lt.east) -- (mlp.south west);
    \node[edg] at ([xshift=-4.2mm]mlp.west) {$\|$};
    % ---- retrieval, placed below to keep the picture inside the text ------
    \node[op, below=9mm of mlp] (rk)
      {$\hat z\to$ cosine rank\\{\scriptsize prefix excluded $\to$ top-10}};
    \draw[flow] (mlp) -- (rk);
    \coordinate (c1) at ([yshift=-3mm]tab.south |- current bounding box.south);
    \coordinate (c2) at (rk.south |- c1);
    \draw[data] (tab.south) -- (c1);
    \draw[data] (c1) -- node[tag, above]
      {the same frozen table is the retrieval target} (c2);
    \draw[data] (c2) -- (rk.south);
    \node[note] at ([yshift=-3.5mm]current bounding box.south west)
      {\textbf{The same information, with the order discarded.} Both summaries
       are strictly causal, so the objective, the padded batching and the
       early stopping are reused \emph{verbatim} from
       Figure~\ref{fig:topo-gru} --- but a mean weights the track from twenty
       steps ago exactly like the one before last, and cannot represent ``A
       then B'' differently from ``B then A''. Only the special-cased $x_t$ tap
       carries any notion of recency, and the concatenation ($\|$) is this
       model's entire account of sequence.};
  \end{tikzpicture}
  \caption{\texttt{seq-ann}, the GRU's non-recurrent twin: the recurrence of
    Figure~\ref{fig:topo-gru} is replaced by two parameter-free summaries of
    the same prefix --- a causal mean-pool concatenated with the last item's
    latent --- feeding a plain three-layer ReLU MLP. Everything else is
    literally the same code (\texttt{seq\_ann.py} imports
    \texttt{make\_batches} and \texttt{step\_loss} from
    \texttt{seq\_nexttrack}), which is what makes this pair a clean test of
    the recurrence itself rather than of a training recipe. Visual convention
    as in Figure~\ref{fig:topo-gru}: the two pooling operators are dashed
    because they hold no parameters, and only the MLP is trained. The
    defensible comparison is same-space: on \dsid{seq-20260715-030509}
    (PCA-128) the pooled twin scores Recall@10 \num{0.0901} (\num{129} of
    \num{1431}) against the GRU's \num{0.1146} (\num{164}), paired $\Delta$
    $-\num{0.0245}$ $[-0.039, -0.011]$, a CI \emph{entirely below zero}.
    \texttt{seq-ann} has no run at all on the canonical PCA-192 split, so it
    never belongs in the same column as the \num{0.12299} bar.}
  \label{fig:topo-ann}
\end{figure}


% ===========================================================================
% 3. seq-dualgru --- two towers, a fusion head, and the pre-MLP variant
% ===========================================================================
\begin{figure}[ht]
  \centering
  \begin{tikzpicture}[topo]
    \node[frozen, minimum width=19mm, minimum height=11mm] (pf)
      {prefix\\latents\\$x_1\dots x_t$};
    % ---- the three configurable causal views ----------------------------
    \node[op, right=7mm of pf] (v2) {\texttt{delta}\\$x_t-x_{t-1}$};
    \node[op, above=2.6mm of v2] (v1) {\texttt{latent}\\$x_t$};
    \node[op, below=2.6mm of v2] (v3) {\texttt{cummean}\\$\frac{1}{t}\sum_{i\le t}x_i$};
    \begin{scope}[on background layer]
      \node[grp, fit=(v1)(v3)] (vg) {};
    \end{scope}
    \draw[data] (pf.east) -- (vg.west);
    % ---- two independent towers -----------------------------------------
    \node[trained, right=8mm of vg, yshift=9mm] (ta)
      {tower A on \texttt{view\_a}\\GRU $h_a{=}256$\\[1pt]{\scriptsize own weights}};
    \node[trained, right=8mm of vg, yshift=-9mm] (tb)
      {tower B on \texttt{view\_b}\\GRU $h_b{=}256$\\[1pt]{\scriptsize own weights}};
    \draw[flow] (vg.east) -- (ta.west);
    \draw[flow] (vg.east) -- (tb.west);
    % ---- concat + the two fusion depths ---------------------------------
    \node[op, right=7mm of ta, yshift=-9mm] (cc)
      {concat\\$h_a\|h_b$\\[1pt]{\scriptsize$512$}};
    \draw[flow] (ta.east) -- (cc.north west);
    \draw[flow] (tb.east) -- (cc.south west);
    \node[won, right=7mm of cc, yshift=9mm] (f0)
      {\texttt{fusion\_layers}${}=0$\\$512\to192$, linear};
    \node[refuted, right=7mm of cc, yshift=-9mm] (f1)
      {\texttt{fusion\_layers}${}=1$\\$512\to256\to192$\\[1pt]{\scriptsize$+$ ReLU $+$ dropout}};
    \draw[flow] (cc.north east) -- (f0.west);
    \draw[flow] (cc.south east) -- (f1.west);
    % ---- the refuted per-step pre-MLP variant ---------------------------
    \node[refuted, below=12mm of tb] (pb)
      {\emph{variant} --- per-step pre-MLP\\on a tower:
       Linear$(192{\to}256)$\\$+$ ReLU $+$ Dropout};
    \draw[variant] (vg.south) |- (pb.west);
    \draw[variant] (pb.north) -- (tb.south);
    \node[note] at ([yshift=-3.5mm]current bounding box.south west)
      {\textbf{Both heads emit the same $192$-d $\hat z$ and retrieve exactly
       as in Figure~\ref{fig:topo-gru}.} The towers share nothing --- separate
       modules, separate parameters --- and their per-step hidden states are
       concatenated, so the combiner sees the towers' \emph{representations}
       rather than their scores and gradients flow back into both. There is no
       bidirectional mode by construction: a backward pass would let step $t$
       read item $t+1$ and copy its own teacher-forcing target. All three views
       are strictly causal for the same reason, and so is the pre-MLP, which is
       applied \emph{timestep-wise}; it is drawn solid because it is trained,
       red because it is refuted.};
  \end{tikzpicture}
  \caption{\texttt{seq-dualgru}: two recurrent towers with entirely separate
    weights read the prefix in parallel and a head fuses their hidden states
    --- the strongest available form of ``train the combination end-to-end so
    the parts co-adapt''. Each tower's input is a configurable causal view
    (\texttt{view\_a} / \texttt{view\_b}); the default is
    \texttt{latent}/\texttt{latent}, i.e.\ two independent GRUs over the same
    stream. Two of the boxes are coloured by the record rather than by the
    design. Removing the MLP fusion head (\texttt{fusion\_layers}${}=0$,
    green) \emph{beats} keeping it in three independent measurements, most
    recently by $+\num{0.02655}$ $[+0.01398, +0.03913]$ CI$>$0 --- and with
    some $9$--$10\%$ \emph{fewer} parameters, so the linear readout is the
    family default; and the
    per-step pre-MLP (red) loses the decisive single-tower ablation, Recall@10
    \num{0.10901} (\num{156} of \num{1431}) against the baseline's \num{176},
    paired $\Delta$ $-\num{0.01398}$ $[-0.02657, -0.00140]$ CI$<$0. Visual
    convention as in Figure~\ref{fig:topo-gru}. The best dual arm reaches
    \num{0.11880} and merely \emph{ties} the single-GRU bar, as does a
    $0.098\%$ parameter-matched width control --- so capacity is not the
    explanation.}
  \label{fig:topo-dualgru}
\end{figure}


% ===========================================================================
% 4. seq-embed --- a learned, weight-tied item table (DESIGN ONLY)
% ===========================================================================
\begin{figure}[ht]
  \centering
  \begin{tikzpicture}[topo]
    \node[frozen, minimum width=20mm, minimum height=11mm] (ids)
      {session prefix\\item \emph{ids}\\$i_1\dots i_t$};
    \node[trained, right=10mm of ids, minimum height=11mm] (E)
      {learned item\\table $E$\\[1pt]{\scriptsize$19{,}402\times192$}};
    \node[trained, right=6mm of E] (gru)
      {GRU $h{=}256$\\{\scriptsize$+$ dropout}};
    \node[trained, right=6mm of gru] (hd)
      {linear head\\{\scriptsize$256\to192$}};
    \node[op, right=6mm of hd] (rk)
      {$\cos(\hat z, E_v)$ over all\\$19{,}402$ rows $\to$ top-10};
    \draw[flow] (ids) -- node[edg, above, inner sep=1pt] {gather} (E);
    \draw[flow] (E) -- (gru);
    \draw[flow] (gru) -- (hd);
    \draw[flow] (hd) -- node[edg, above, inner sep=1pt] {$\hat z$} (rk);
    % ---- the three modes -------------------------------------------------
    \node[op, text width=64mm, align=flush left, anchor=south west] (md)
      at ([xshift=4mm,yshift=8mm]E.north east)
      {\texttt{embed\_mode} --- what the table $E$ actually is:\\
       \textbf{\texttt{frozen}} $E=Z$, nothing trained (the reproduction gate)\\
       \textbf{\texttt{residual}} $E=Z+\Delta E$: frozen content vector $+$ a
       trained correction, so a cold item still has a position\\
       \textbf{\texttt{free}} $E$ fully trained, optionally randomly initialised};
    \draw[flow] (md.west) -- (E.north east);
    % ---- weight tying: the same table is the retrieval target -----------
    \coordinate (c1) at ([yshift=-8mm]E.south);
    \coordinate (c2) at (rk.south |- c1);
    \draw[tie] (E.south) -- (c1);
    \draw[tie] (c1) -- node[tag, above, text=topogood]
      {weight tying --- the \emph{same} table is\\the input embedding
       \emph{and} the retrieval target} (c2);
    \draw[tie] (c2) -- (rk.south);
    % ---- the gradient loop that makes the geometry trainable ------------
    \coordinate (t1) at ([yshift=5mm]rk.north |- current bounding box.north);
    \coordinate (t2) at (E.north |- t1);
    \draw[grad] (rk.north) -- (t1);
    \draw[grad] (t1) -- node[tag, above, text=topobad]
      {$\nabla$ --- one InfoNCE step moves the query \emph{and} every target:
       the geometry itself is trainable} (t2);
    \draw[grad] (t2) -- (E.north);
    \node[note] at ([yshift=-3.5mm]current bounding box.south west)
      {\textbf{Why the loop is the whole point.} In
       Figures~\ref{fig:topo-gru}--\ref{fig:topo-dualgru} the retrieval space
       is fixed and only the map into it is learned; here the space is a
       parameter, shaped by next-track adjacency rather than inherited from a
       reconstruction objective. That also creates a collapse hazard absent
       everywhere else: with a learned table a plain cosine objective has a
       trivial optimum at ``send every item to one point'', so only a
       contrastive loss is admissible and \texttt{loss=cosine} is rejected
       outright for the trained modes.};
  \end{tikzpicture}
  \caption{\texttt{seq-embed}, the one axis the record has not tested:
    \textbf{a design, not a result}. No run of this predictor exists on any
    dataset, it is absent from the server's registered predictor list, and
    nothing here is a measured score. The item table $E$ is used twice, as the
    GRU's input embedding and as the retrieval target (weight tying, the
    GRU4Rec/SASRec formulation), so a single gradient step moves the query and
    all the candidates together --- the return arrow along the top is what
    makes the geometry trainable, and it is exactly what every other topology
    in this document lacks. The three \texttt{embed\_mode} settings span the
    space from a reproduction gate to a fully free table. What \emph{is}
    measured is the data budget that bounds them: \num{73632} training steps
    to fit \num{15816} item vectors, a median per-item training frequency of
    \num{1}, only \num{3294} items occurring five or more times, and
    \num{788} of \num{1431} test targets never appearing in any training
    session --- so a \texttt{free} table cannot retrieve the cold
    \SI{55}{\percent} at all and is capped at Recall@10 \num{0.4493}, whereas
    \texttt{residual} keeps the content vector and with it the cold reach.
    Visual convention as in Figure~\ref{fig:topo-gru}.}
  \label{fig:topo-embed}
\end{figure}


% ===========================================================================
% 5. seq-blend --- the champion
% ===========================================================================
\begin{figure}[ht]
  \centering
  \begin{tikzpicture}[topo]
    \node[frozen, minimum width=18mm, minimum height=11mm] (pf)
      {session\\prefix\\$i_1\dots i_t$};
    % ---- the learned metric map: shared by two legs, bypassed by one -----
    \node[trained, right=4mm of pf] (W)
      {learned map $W$\\$192\times192$\\[1pt]{\scriptsize no bias, identity init}\\{\scriptsize
       InfoNCE $\tau{=}0.07$}\\{\scriptsize $36{,}864$ params}};
    % ---- the three legs -------------------------------------------------
    \node[learnspace, right=3mm of W] (C)
      {\Cp\ content-$k$NN in $W$-space\\$\max_{u}\cos(Wz_v,Wz_u)$\\[1pt]{\scriptsize
       no parameters of its own}};
    \node[trained, above=7mm of C] (R)
      {\Rp\ projected GRU\\$h{=}256$, $394{,}944$ params};
    \node[lookup, below=7mm of C] (M)
      {$M$ first-order Markov counts\\$P(v\mid u)$, $u=$ last prefix item\\[1pt]{\scriptsize
       Dirichlet trust $n_u/(n_u{+}8)$,}\\{\scriptsize artist/genre back-off; $0$ params}};
    \draw[flow] (pf.east) -- (W.west);
    % ---- the frozen item table W acts ON: the same Z as Figure~\ref{fig:topo-gru} --
    \node[frozen, anchor=south east, minimum width=18mm] (Z)
      at ([yshift=6mm]W.north west)
      {frozen item latents $Z$\\{\scriptsize$19{,}402\times192$}};
    \draw[data] (Z.south east) -- node[edg, left, inner sep=1.2pt] {$z_i$} (W.north west);
    \draw[proj] (W.north east) -- (R.west);
    \draw[proj] (W.east) -- (C.west);
    \draw[data] (pf.south) |- node[tag, above, pos=0.55]
      {item \emph{ids} only ---\\the $M$ leg bypasses $Z$ and $W$} (M.west);
    % ---- the fixed combiner ---------------------------------------------
    \node[op, right=3mm of C] (cb)
      {z-normalise each leg\\over the candidate set,\\then average\\[1pt]{\scriptsize
       equal thirds ---}\\{\scriptsize\emph{no learned weights}}};
    \draw[flow] (R.east) -- node[edg, above, sloped, pos=0.55] {$s_{\Rp}$} (cb.north west);
    \draw[flow] (C.east) -- node[edg, above, inner sep=1pt] {$s_{\Cp}$} (cb.west);
    \draw[flow] (M.east) -- node[edg, above, sloped, pos=0.30] {$s_M$} (cb.south west);
    \node[op, right=2mm of cb] (rk) {rank\\$\to$ top-10};
    \draw[flow] (cb) -- (rk);
    \node[note] at ([yshift=-3.5mm]current bounding box.south west)
      {\textbf{Which legs are learned.} $W$ and the \Rp\ recurrence are
       trained (\num{431808} parameters between them); \Cp\ has no parameters
       of its own but retrieves inside $W$'s learned geometry, which is why it
       is drawn dashed-blue; $M$ is a table of training-set transition counts,
       drawn dashed-purple. The combiner is a fixed rule: each leg's
       full-vocabulary score vector is z-scored over the admissible candidates
       (every item minus those already in the prefix) and the three are
       averaged with hard-coded equal weights. The z-scoring is the
       load-bearing step --- it puts three incomparable score scales onto one
       per-session scale while \emph{preserving} relative magnitude, so a
       confidently-peaked leg can outvote two diffuse ones. Rank-only fusion,
       which discards that magnitude, loses by $-\num{0.029}$ CI$<$0.};
  \end{tikzpicture}
  \caption{\texttt{seq-blend} with \texttt{projection=true}, the champion:
    three full-vocabulary score vectors, z-normalised over the candidate set
    and averaged in fixed equal thirds. $W$ acts on the \emph{frozen item
    latents} $Z$ of Figure~\ref{fig:topo-gru} --- it warps the retrieval geometry,
    it does not read the prefix ids directly --- and it is shared by two of the
    three legs and bypassed by the third; that is the whole mechanism. $W$ is fit by in-batch InfoNCE on consecutive within-train
    pairs, which rebuilds the retrieval metric from ``reconstruct the item
    matrix'' into ``what follows what''; both legs that retrieve by cosine in
    that space improve --- the parameter-free content-$k$NN leg more than
    doubles, \num{0.0685} to \num{0.1516}, and the recurrent leg rises from
    \num{0.1230} to $\approx$\num{0.172} --- while the item-ID lookup is
    untouched at \num{0.106918}. The two projected-leg standalone figures come
    from the off-server driver of the 2026-07-18 campaign and have no server
    run of their own; only the $M$ leg's number is a canonical-split run.
    Together the three legs reach Recall@10 \num{0.21174} (\num{303} of
    \num{1431}), paired $\Delta$ over the unprojected three-leg blend
    $+\num{0.0259}$ $[+0.012, +0.041]$, a CI entirely above zero and robust
    across three seeds. Visual convention as in Figure~\ref{fig:topo-gru}.}
  \label{fig:topo-blend}
\end{figure}


% ===========================================================================
% 6. seq-stacker --- the learned combiner
% ===========================================================================
\begin{figure}[ht]
  \centering
  \begin{tikzpicture}[topo]
    \node[frozen, minimum width=18mm, minimum height=11mm] (pf)
      {session\\prefix\\$i_1\dots i_t$};
    \node[trained, right=6mm of pf] (R)
      {$R$ GRU $h{=}256$\\[1pt]{\scriptsize trained separately}};
    \node[lookup, above=6mm of R] (M) {$M$ Markov counts};
    \node[trained, below=6mm of R] (A)
      {$A$ pooled MLP\\[1pt]{\scriptsize trained separately}};
    \draw[flow] (pf.east) -- (M.west);
    \draw[flow] (pf.east) -- (R.west);
    \draw[flow] (pf.east) -- (A.west);
    % ---- the per-candidate feature row ----------------------------------
    \node[op, right=6mm of R] (fx)
      {per-candidate feature row\\$[\,s_M,\ s_R,\ s_A,\ \log(1{+}\text{plays}),$\\$\text{recency},\
       \text{artist match},\ \text{genre match}\,]$};
    \draw[flow] (M.east) -- (fx.north west);
    \draw[flow] (R.east) -- (fx.west);
    \draw[flow] (A.east) -- (fx.south west);
    \node[refuted, right=5mm of fx] (xgb)
      {XGBoost meta-learner\\$200$ rounds, depth $6$\\[1pt]{\scriptsize\texttt{binary:logistic}}};
    \draw[flow] (fx) -- (xgb);
    \node[op, below=9mm of xgb] (rk)
      {score all $19{,}402$ items\\$\to$ top-10};
    \draw[flow] (xgb) -- (rk);
    % ---- the two failure mechanisms, marked on the diagram --------------
    \node[refuted, text width=56mm, align=flush left, anchor=north west] (tr)
      at ([yshift=-10mm]fx.south west)
      {\textbf{Training rows}: $1$ positive $+$ $50$ popularity-sampled
       negatives per train session, i.e.\ \num{5723} positives, against the
       \num{73632} teacher-forced steps its own $R$ leg trains on
       ($12.9\times$ less signal) --- and a decision surface calibrated on
       $51$ candidates ($0.26\%$ of the catalogue) is then asked to rank all
       \num{19402}.};
    \draw[grad] (tr.north east) -- (xgb.south west);
    \node[note] at ([yshift=-3.5mm]current bounding box.south west)
      {\textbf{Why the shape matters.} A $200$-round depth-$6$ booster is
       piecewise \emph{constant}, so ranking \num{19402} candidates with it
       produces large tied blocks that the harness breaks by item index. The
       fixed z-blend of Figure~\ref{fig:topo-blend} is instead a strictly
       monotone continuous function of the leg scores: it preserves each leg's
       internal ordering exactly and can only trade the legs off against one
       another.};
  \end{tikzpicture}
  \caption{\texttt{seq-stacker}, the learned counterpart of
    Figure~\ref{fig:topo-blend}'s fixed rule: the same separately-trained base
    legs, but their per-candidate scores become columns of a feature row
    alongside four cheap candidate features, and a gradient-boosted tree
    learns the combination. It earns a diagram because its failure is
    informative --- with the \texttt{markov}, \texttt{gru} and \texttt{ann}
    legs it reaches Recall@10 \num{0.085255} (\num{122} of \num{1431}),
    \emph{below its own GRU base leg} at \num{0.12299}, paired $\Delta$
    $-\num{0.0377}$ $[-0.056, -0.020]$ CI$<$0. The red callout names the two
    mechanisms: supervision starvation, since the meta-learner gets one
    labelled event per session where its base leg gets one per prefix
    position, and a train/serve candidate mismatch. The booster is drawn solid
    because it is trained and red because it is refuted; the feature row and
    the ranking step are dashed because they hold no parameters. Visual
    convention as in Figure~\ref{fig:topo-gru}. This is one of four
    independent learned-combiner nulls, the others being reciprocal-rank
    fusion, oracle weights fitted on the test labels, and the jointly-trained
    dual tower of Figure~\ref{fig:topo-dualgru}.}
  \label{fig:topo-stacker}
\end{figure}

%
%  (\graphicspath{{figures/}} governs \includegraphics only, NOT \input, so the
%   figures/ prefix is spelled out above.)
%
%  ------------------------------------------------------------------------
%  REQUIRED PREAMBLE LINES in the main document --- add these two lines after
%  \usepackage{graphicx}, before \begin{document}:
%
%      \usepackage{tikz}
%      \usetikzlibrary{positioning,arrows.meta,fit,backgrounds,calc}
%
%  Nothing else is required: xcolor arrives with tikz, and every colour and
%  every style used below is defined in this file.  \dsid, \Rp and \Cp are
%  \providecommand-guarded here so the file also compiles inside a bare
%  wrapper; when the main document defines them, its definitions win.
%
%  PACKAGES DELIBERATELY *NOT* REQUIRED: amssymb is absent from the house
%  preamble, so nothing below uses \mathbb.  Every picture is laid out with
%  RELATIVE positioning (right=/above=/below= of ...), never with hard
%  coordinates, so no node can collide with its neighbour if the font or the
%  wording changes; the per-figure notes and the legend are anchored to
%  `current bounding box', so they always land underneath the diagram.
%  ------------------------------------------------------------------------
%
%  LABELS DEFINED IN THIS FILE (in file order)
%      fig:topo-gru       seq-nexttrack --- the baseline; ALSO CARRIES THE LEGEND
%      fig:topo-ann       seq-ann       --- pooled feed-forward twin
%      fig:topo-dualgru   seq-dualgru   --- two towers + fusion head + pre-MLP variant
%      fig:topo-embed     seq-embed     --- learned, weight-tied item table (DESIGN ONLY)
%      fig:topo-blend     seq-blend     --- the champion, three legs
%      fig:topo-stacker   seq-stacker   --- learned XGBoost combiner
%
%  VISUAL LANGUAGE (the primary axis is the LINE, not the colour, so it
%  survives greyscale printing; stated for the reader in the caption of
%  fig:topo-gru and re-pointed-at by every other caption)
%      SOLID  border  = TRAINED component (parameters fitted by gradient descent)
%      DASHED border  = FROZEN or PARAMETER-FREE component
%      blue   = the on-record default path
%      purple dashed  = non-parametric table lookup over the dataset
%      dashed blue    = parameter-free, but operating inside a LEARNED geometry
%      green  = variant favoured by the record
%      red    = variant refuted by the record
% ===========================================================================

% ===========================================================================
%  topo-style.tex --- the shared TikZ style block for the next-track topology
%  diagrams AND for the evaluation-protocol diagram in the body of
%  model-comparison.tex.  Extracted from topologies.tex so that the protocol
%  figure, which appears BEFORE \input{figures/topologies}, can use the same
%  visual language.  \input from docs/ (compilation cwd):
%
%      \input{figures/topo-style}
%
%  Idempotent: guarded, so a second \input is a no-op.  Requires
%      \usepackage{tikz}
%      \usetikzlibrary{positioning,arrows.meta,fit,backgrounds,calc}
% ===========================================================================
\ifdefined\nexttracktopostyle\else
\newcommand{\nexttracktopostyle}{}

\providecommand{\dsid}[1]{\texttt{\small #1}}
\providecommand{\Rp}{\ensuremath{R'}}
\providecommand{\Cp}{\ensuremath{C'}}

% --- palette: the same hex values as docs/figures/make_figures.py ----------
\definecolor{topoacc}{HTML}{2563EB}   % ACCENT  --- default / on-record path
\definecolor{topogood}{HTML}{16A34A}  % GOOD    --- favoured variant
\definecolor{topobad}{HTML}{DC2626}   % BAD     --- refuted variant
\definecolor{topomut}{HTML}{9CA3AF}   % MUTED   --- frozen, parameter-free
\definecolor{topolook}{HTML}{9333EA}  % seq-markov purple --- table lookup
\definecolor{topoink}{HTML}{374151}   % annotation ink

\tikzset{
  topo/.style={
    font=\footnotesize,
    every node/.append style={align=center},
  },
  box/.style={rounded corners=1.4pt, inner xsep=3.2pt, inner ysep=2.8pt,
              minimum height=6.4mm, text=topoink},
  dsh/.style={dash pattern=on 2pt off 1.4pt},
  % --- component styles ---------------------------------------------------
  trained/.style={box, draw=topoacc, line width=0.9pt, fill=topoacc!7},
  frozen/.style={box, dsh, draw=topomut, line width=0.8pt, fill=topomut!14},
  lookup/.style={box, dsh, draw=topolook, line width=0.9pt, fill=topolook!7},
  learnspace/.style={box, dsh, draw=topoacc, line width=0.9pt, fill=topoacc!5},
  won/.style={box, draw=topogood, line width=1.0pt, fill=topogood!8},
  refuted/.style={box, draw=topobad, line width=1.0pt, fill=topobad!6},
  op/.style={box, dash pattern=on 1.5pt off 1.2pt, draw=topomut,
             line width=0.7pt, fill=white, rounded corners=3.2pt},
  % --- edges --------------------------------------------------------------
  flow/.style={-{Stealth[length=4.2pt,width=3.4pt]}, draw=topoink,
               line width=0.7pt},
  proj/.style={flow, draw=topoacc},
  data/.style={flow, draw=topomut, dsh},
  tie/.style={flow, draw=topogood, line width=0.8pt},
  grad/.style={flow, draw=topobad, line width=0.75pt,
               dash pattern=on 1.6pt off 1.3pt},
  variant/.style={flow, draw=topobad, line width=0.75pt,
                  dash pattern=on 1.6pt off 1.3pt},
  % --- text ---------------------------------------------------------------
  tag/.style={font=\scriptsize, text=topoink, inner sep=1.2pt},
  edg/.style={font=\scriptsize, text=topoink, inner sep=1.4pt},
  note/.style={font=\scriptsize, text=topoink, align=flush left,
               text width=134mm, inner sep=1.4pt, anchor=north west},
  grp/.style={draw=topomut!80, line width=0.5pt, rounded corners=2.4pt,
              dash pattern=on 1pt off 1pt, inner sep=3.4pt},
  swatch/.style={minimum width=5.2mm, minimum height=3.2mm, inner sep=0pt},
}
\fi




% ===========================================================================
% 1. seq-nexttrack --- the baseline (and the legend)
% ===========================================================================
\begin{figure}[ht]
  \centering
  \begin{tikzpicture}[topo]
    % ---- the frozen item-latent table, read TWICE ------------------------
    \node[frozen, minimum width=20mm, minimum height=12mm] (tab)
      {frozen item\\latents $Z$\\[1pt]{\scriptsize$19{,}402\times192$}};
    % ---- the prefix, in order (one column per step) ----------------------
    \node[frozen, minimum width=12mm, right=7mm of tab]  (x1) {$x_1$};
    \node[frozen, minimum width=12mm, right=4mm of x1]   (x2) {$x_2$};
    \node[tag, right=2.5mm of x2]                        (xd) {$\cdots$};
    \node[frozen, minimum width=12mm, right=2.5mm of xd] (x3) {$x_t$};
    \begin{scope}[on background layer]
      \node[grp, fit=(x1)(x3)] (pfx) {};
    \end{scope}
    \draw[data] (tab.east) -- (pfx.west);
    \node[tag, below=1.2mm of pfx] (pfxlab)
      {prefix latents --- \emph{order preserved}};
    % ---- the recurrence --------------------------------------------------
    \node[trained, minimum width=12mm, above=7mm of x1] (g1) {GRU};
    \node[trained, minimum width=12mm, above=7mm of x2] (g2) {GRU};
    \node[tag] (gd) at (g1 -| xd) {$\cdots$};
    \node[trained, minimum width=12mm, above=7mm of x3] (g3) {GRU};
    \draw[flow] (x1) -- (g1);
    \draw[flow] (x2) -- (g2);
    \draw[flow] (x3) -- (g3);
    \draw[flow] (g1) -- node[edg, above, inner sep=0.7pt] {$h_t$} (g2);
    \draw[flow] (g2) -- (gd);
    \draw[flow] (gd) -- (g3);
    % ---- head and retrieval ---------------------------------------------
    \node[trained, right=6mm of g3] (head)
      {linear head\\{\scriptsize$256\to192$}};
    \node[op, right=7mm of head] (rk)
      {rank the catalogue by $\cos(\hat z, z_v)$\\{\scriptsize prefix excluded $\to$ top-10}};
    \draw[flow] (g3) -- (head);
    \draw[flow] (head) -- node[edg, above, inner sep=1pt] {$\hat z$} (rk);
    % ---- the SAME frozen table is the retrieval target -------------------
    \coordinate (c1) at ([yshift=-3mm]tab.south |- current bounding box.south);
    \coordinate (c2) at (rk.south |- c1);
    \draw[data] (tab.south) -- (c1);
    \draw[data] (c1) -- node[tag, above]
      {the \emph{same} frozen table is the retrieval target} (c2);
    \draw[data] (c2) -- (rk.south);
    % ---- supervision note ------------------------------------------------
    \node[note] (nt) at ([yshift=-3.5mm]current bounding box.south west)
      {\textbf{Supervision --- per-step teacher forcing.} The head is shared
       across positions and must predict the next item's latent at
       \emph{every} prefix position, so one session of length $L$ yields
       $L-1$ examples: $73{,}632$ (prefix, next-item) pairs on the train
       split. The loss is masked in-batch InfoNCE at $\tau=0.07$ --- each
       step's true next latent must outscore the other steps' targets in the
       same minibatch, which supplies negatives for free.};
    % ---- LEGEND (two rows so it cannot widen the picture) ----------------
    \node[trained, swatch, anchor=north west] (l1)
      at ([yshift=-3mm]current bounding box.south west) {};
    \node[tag, right=1.1mm of l1] (t1) {trained (solid)};
    \node[frozen, swatch, right=3.4mm of t1] (l2) {};
    \node[tag, right=1.1mm of l2] (t2) {frozen / parameter-free (dashed)};
    \node[lookup, swatch, right=3.4mm of t2] (l3) {};
    \node[tag, right=1.1mm of l3] (t3) {table lookup};
    \node[learnspace, swatch, anchor=north west] (l4)
      at ([yshift=-1.6mm]l1.south west) {};
    \node[tag, right=1.1mm of l4] (t4) {parameter-free, learned space};
    \node[won, swatch, right=3.4mm of t4] (l5) {};
    \node[tag, right=1.1mm of l5] (t5) {favoured variant};
    \node[refuted, swatch, right=3.4mm of t5] (l6) {};
    \node[tag, right=1.1mm of l6] (t6) {refuted variant};
  \end{tikzpicture}
  \caption{\texttt{seq-nexttrack}, the single-model baseline every other
    topology in this document is measured against: a frozen item space, a
    recurrence over the \emph{ordered} prefix, a linear head, and cosine
    retrieval back into the \emph{same} frozen space. Nothing about the item
    geometry is learned here --- the $19{,}402\times192$ latent table (offline
    PCA-192 on \dsid{seq-20260715-131139}) is read twice, once to embed the
    prefix and once as the retrieval target, and only the GRU and its linear
    head carry parameters (\num{394944} at $h=256$). \textbf{The visual
    language used throughout this document is fixed here.} A \emph{solid}
    border marks a trained component and a \emph{dashed} border a frozen or
    parameter-free one; that distinction rides on the line style, not the
    colour, so it survives greyscale. Colour then carries the verdict ---
    blue: the on-record default path; purple dashed: a non-parametric table
    lookup over the dataset; dashed blue: parameter-free but operating inside
    a \emph{learned} geometry; green: a variant the record favours; red: a
    variant the record refutes. On the canonical split this configuration
    scores Recall@10 \num{0.12299} (\num{176} of \num{1431}), reproduced
    bit-identically eleven times.}
  \label{fig:topo-gru}
\end{figure}


% ===========================================================================
% 2. seq-ann --- the pooled feed-forward twin
% ===========================================================================
\begin{figure}[ht]
  \centering
  \begin{tikzpicture}[topo]
    \node[frozen, minimum width=20mm, minimum height=12mm] (tab)
      {frozen item\\latents $Z$\\[1pt]{\scriptsize$19{,}402\times192$}};
    \node[frozen, minimum width=12mm, right=7mm of tab]  (x1) {$x_1$};
    \node[frozen, minimum width=12mm, right=4mm of x1]   (x2) {$x_2$};
    \node[tag, right=2.5mm of x2]                        (xd) {$\cdots$};
    \node[frozen, minimum width=12mm, right=2.5mm of xd] (x3) {$x_t$};
    \begin{scope}[on background layer]
      \node[grp, fit=(x1)(x3)] (pfx) {};
    \end{scope}
    \draw[data] (tab.east) -- (pfx.west);
    % ---- the two parameter-free summaries of the same prefix -------------
    \node[op, right=8mm of pfx, yshift=10mm] (mp)
      {causal mean-pool\\$\frac{1}{t}\sum_{i\le t}x_i$\\[1pt]{\scriptsize\emph{permutation-invariant}}};
    \node[op, right=8mm of pfx, yshift=-10mm] (lt)
      {last item\\$x_t$\\[1pt]{\scriptsize the only recency cue}};
    \draw[flow] (pfx.east) -- (mp.west);
    \draw[flow] (pfx.east) -- (lt.west);
    % ---- the MLP ---------------------------------------------------------
    \node[trained, right=5mm of mp, yshift=-10mm] (mlp)
      {MLP\\$384\to256\to256\to192$\\[1pt]{\scriptsize ReLU $+$ dropout}};
    \draw[flow] (mp.east) -- (mlp.north west);
    \draw[flow] (lt.east) -- (mlp.south west);
    \node[edg] at ([xshift=-4.2mm]mlp.west) {$\|$};
    % ---- retrieval, placed below to keep the picture inside the text ------
    \node[op, below=9mm of mlp] (rk)
      {$\hat z\to$ cosine rank\\{\scriptsize prefix excluded $\to$ top-10}};
    \draw[flow] (mlp) -- (rk);
    \coordinate (c1) at ([yshift=-3mm]tab.south |- current bounding box.south);
    \coordinate (c2) at (rk.south |- c1);
    \draw[data] (tab.south) -- (c1);
    \draw[data] (c1) -- node[tag, above]
      {the same frozen table is the retrieval target} (c2);
    \draw[data] (c2) -- (rk.south);
    \node[note] at ([yshift=-3.5mm]current bounding box.south west)
      {\textbf{The same information, with the order discarded.} Both summaries
       are strictly causal, so the objective, the padded batching and the
       early stopping are reused \emph{verbatim} from
       Figure~\ref{fig:topo-gru} --- but a mean weights the track from twenty
       steps ago exactly like the one before last, and cannot represent ``A
       then B'' differently from ``B then A''. Only the special-cased $x_t$ tap
       carries any notion of recency, and the concatenation ($\|$) is this
       model's entire account of sequence.};
  \end{tikzpicture}
  \caption{\texttt{seq-ann}, the GRU's non-recurrent twin: the recurrence of
    Figure~\ref{fig:topo-gru} is replaced by two parameter-free summaries of
    the same prefix --- a causal mean-pool concatenated with the last item's
    latent --- feeding a plain three-layer ReLU MLP. Everything else is
    literally the same code (\texttt{seq\_ann.py} imports
    \texttt{make\_batches} and \texttt{step\_loss} from
    \texttt{seq\_nexttrack}), which is what makes this pair a clean test of
    the recurrence itself rather than of a training recipe. Visual convention
    as in Figure~\ref{fig:topo-gru}: the two pooling operators are dashed
    because they hold no parameters, and only the MLP is trained. The
    defensible comparison is same-space: on \dsid{seq-20260715-030509}
    (PCA-128) the pooled twin scores Recall@10 \num{0.0901} (\num{129} of
    \num{1431}) against the GRU's \num{0.1146} (\num{164}), paired $\Delta$
    $-\num{0.0245}$ $[-0.039, -0.011]$, a CI \emph{entirely below zero}.
    \texttt{seq-ann} has no run at all on the canonical PCA-192 split, so it
    never belongs in the same column as the \num{0.12299} bar.}
  \label{fig:topo-ann}
\end{figure}


% ===========================================================================
% 3. seq-dualgru --- two towers, a fusion head, and the pre-MLP variant
% ===========================================================================
\begin{figure}[ht]
  \centering
  \begin{tikzpicture}[topo]
    \node[frozen, minimum width=19mm, minimum height=11mm] (pf)
      {prefix\\latents\\$x_1\dots x_t$};
    % ---- the three configurable causal views ----------------------------
    \node[op, right=7mm of pf] (v2) {\texttt{delta}\\$x_t-x_{t-1}$};
    \node[op, above=2.6mm of v2] (v1) {\texttt{latent}\\$x_t$};
    \node[op, below=2.6mm of v2] (v3) {\texttt{cummean}\\$\frac{1}{t}\sum_{i\le t}x_i$};
    \begin{scope}[on background layer]
      \node[grp, fit=(v1)(v3)] (vg) {};
    \end{scope}
    \draw[data] (pf.east) -- (vg.west);
    % ---- two independent towers -----------------------------------------
    \node[trained, right=8mm of vg, yshift=9mm] (ta)
      {tower A on \texttt{view\_a}\\GRU $h_a{=}256$\\[1pt]{\scriptsize own weights}};
    \node[trained, right=8mm of vg, yshift=-9mm] (tb)
      {tower B on \texttt{view\_b}\\GRU $h_b{=}256$\\[1pt]{\scriptsize own weights}};
    \draw[flow] (vg.east) -- (ta.west);
    \draw[flow] (vg.east) -- (tb.west);
    % ---- concat + the two fusion depths ---------------------------------
    \node[op, right=7mm of ta, yshift=-9mm] (cc)
      {concat\\$h_a\|h_b$\\[1pt]{\scriptsize$512$}};
    \draw[flow] (ta.east) -- (cc.north west);
    \draw[flow] (tb.east) -- (cc.south west);
    \node[won, right=7mm of cc, yshift=9mm] (f0)
      {\texttt{fusion\_layers}${}=0$\\$512\to192$, linear};
    \node[refuted, right=7mm of cc, yshift=-9mm] (f1)
      {\texttt{fusion\_layers}${}=1$\\$512\to256\to192$\\[1pt]{\scriptsize$+$ ReLU $+$ dropout}};
    \draw[flow] (cc.north east) -- (f0.west);
    \draw[flow] (cc.south east) -- (f1.west);
    % ---- the refuted per-step pre-MLP variant ---------------------------
    \node[refuted, below=12mm of tb] (pb)
      {\emph{variant} --- per-step pre-MLP\\on a tower:
       Linear$(192{\to}256)$\\$+$ ReLU $+$ Dropout};
    \draw[variant] (vg.south) |- (pb.west);
    \draw[variant] (pb.north) -- (tb.south);
    \node[note] at ([yshift=-3.5mm]current bounding box.south west)
      {\textbf{Both heads emit the same $192$-d $\hat z$ and retrieve exactly
       as in Figure~\ref{fig:topo-gru}.} The towers share nothing --- separate
       modules, separate parameters --- and their per-step hidden states are
       concatenated, so the combiner sees the towers' \emph{representations}
       rather than their scores and gradients flow back into both. There is no
       bidirectional mode by construction: a backward pass would let step $t$
       read item $t+1$ and copy its own teacher-forcing target. All three views
       are strictly causal for the same reason, and so is the pre-MLP, which is
       applied \emph{timestep-wise}; it is drawn solid because it is trained,
       red because it is refuted.};
  \end{tikzpicture}
  \caption{\texttt{seq-dualgru}: two recurrent towers with entirely separate
    weights read the prefix in parallel and a head fuses their hidden states
    --- the strongest available form of ``train the combination end-to-end so
    the parts co-adapt''. Each tower's input is a configurable causal view
    (\texttt{view\_a} / \texttt{view\_b}); the default is
    \texttt{latent}/\texttt{latent}, i.e.\ two independent GRUs over the same
    stream. Two of the boxes are coloured by the record rather than by the
    design. Removing the MLP fusion head (\texttt{fusion\_layers}${}=0$,
    green) \emph{beats} keeping it in three independent measurements, most
    recently by $+\num{0.02655}$ $[+0.01398, +0.03913]$ CI$>$0 --- and with
    some $9$--$10\%$ \emph{fewer} parameters, so the linear readout is the
    family default; and the
    per-step pre-MLP (red) loses the decisive single-tower ablation, Recall@10
    \num{0.10901} (\num{156} of \num{1431}) against the baseline's \num{176},
    paired $\Delta$ $-\num{0.01398}$ $[-0.02657, -0.00140]$ CI$<$0. Visual
    convention as in Figure~\ref{fig:topo-gru}. The best dual arm reaches
    \num{0.11880} and merely \emph{ties} the single-GRU bar, as does a
    $0.098\%$ parameter-matched width control --- so capacity is not the
    explanation.}
  \label{fig:topo-dualgru}
\end{figure}


% ===========================================================================
% 4. seq-embed --- a learned, weight-tied item table (DESIGN ONLY)
% ===========================================================================
\begin{figure}[ht]
  \centering
  \begin{tikzpicture}[topo]
    \node[frozen, minimum width=20mm, minimum height=11mm] (ids)
      {session prefix\\item \emph{ids}\\$i_1\dots i_t$};
    \node[trained, right=10mm of ids, minimum height=11mm] (E)
      {learned item\\table $E$\\[1pt]{\scriptsize$19{,}402\times192$}};
    \node[trained, right=6mm of E] (gru)
      {GRU $h{=}256$\\{\scriptsize$+$ dropout}};
    \node[trained, right=6mm of gru] (hd)
      {linear head\\{\scriptsize$256\to192$}};
    \node[op, right=6mm of hd] (rk)
      {$\cos(\hat z, E_v)$ over all\\$19{,}402$ rows $\to$ top-10};
    \draw[flow] (ids) -- node[edg, above, inner sep=1pt] {gather} (E);
    \draw[flow] (E) -- (gru);
    \draw[flow] (gru) -- (hd);
    \draw[flow] (hd) -- node[edg, above, inner sep=1pt] {$\hat z$} (rk);
    % ---- the three modes -------------------------------------------------
    \node[op, text width=64mm, align=flush left, anchor=south west] (md)
      at ([xshift=4mm,yshift=8mm]E.north east)
      {\texttt{embed\_mode} --- what the table $E$ actually is:\\
       \textbf{\texttt{frozen}} $E=Z$, nothing trained (the reproduction gate)\\
       \textbf{\texttt{residual}} $E=Z+\Delta E$: frozen content vector $+$ a
       trained correction, so a cold item still has a position\\
       \textbf{\texttt{free}} $E$ fully trained, optionally randomly initialised};
    \draw[flow] (md.west) -- (E.north east);
    % ---- weight tying: the same table is the retrieval target -----------
    \coordinate (c1) at ([yshift=-8mm]E.south);
    \coordinate (c2) at (rk.south |- c1);
    \draw[tie] (E.south) -- (c1);
    \draw[tie] (c1) -- node[tag, above, text=topogood]
      {weight tying --- the \emph{same} table is\\the input embedding
       \emph{and} the retrieval target} (c2);
    \draw[tie] (c2) -- (rk.south);
    % ---- the gradient loop that makes the geometry trainable ------------
    \coordinate (t1) at ([yshift=5mm]rk.north |- current bounding box.north);
    \coordinate (t2) at (E.north |- t1);
    \draw[grad] (rk.north) -- (t1);
    \draw[grad] (t1) -- node[tag, above, text=topobad]
      {$\nabla$ --- one InfoNCE step moves the query \emph{and} every target:
       the geometry itself is trainable} (t2);
    \draw[grad] (t2) -- (E.north);
    \node[note] at ([yshift=-3.5mm]current bounding box.south west)
      {\textbf{Why the loop is the whole point.} In
       Figures~\ref{fig:topo-gru}--\ref{fig:topo-dualgru} the retrieval space
       is fixed and only the map into it is learned; here the space is a
       parameter, shaped by next-track adjacency rather than inherited from a
       reconstruction objective. That also creates a collapse hazard absent
       everywhere else: with a learned table a plain cosine objective has a
       trivial optimum at ``send every item to one point'', so only a
       contrastive loss is admissible and \texttt{loss=cosine} is rejected
       outright for the trained modes.};
  \end{tikzpicture}
  \caption{\texttt{seq-embed}, the one axis the record has not tested:
    \textbf{a design, not a result}. No run of this predictor exists on any
    dataset, it is absent from the server's registered predictor list, and
    nothing here is a measured score. The item table $E$ is used twice, as the
    GRU's input embedding and as the retrieval target (weight tying, the
    GRU4Rec/SASRec formulation), so a single gradient step moves the query and
    all the candidates together --- the return arrow along the top is what
    makes the geometry trainable, and it is exactly what every other topology
    in this document lacks. The three \texttt{embed\_mode} settings span the
    space from a reproduction gate to a fully free table. What \emph{is}
    measured is the data budget that bounds them: \num{73632} training steps
    to fit \num{15816} item vectors, a median per-item training frequency of
    \num{1}, only \num{3294} items occurring five or more times, and
    \num{788} of \num{1431} test targets never appearing in any training
    session --- so a \texttt{free} table cannot retrieve the cold
    \SI{55}{\percent} at all and is capped at Recall@10 \num{0.4493}, whereas
    \texttt{residual} keeps the content vector and with it the cold reach.
    Visual convention as in Figure~\ref{fig:topo-gru}.}
  \label{fig:topo-embed}
\end{figure}


% ===========================================================================
% 5. seq-blend --- the champion
% ===========================================================================
\begin{figure}[ht]
  \centering
  \begin{tikzpicture}[topo]
    \node[frozen, minimum width=18mm, minimum height=11mm] (pf)
      {session\\prefix\\$i_1\dots i_t$};
    % ---- the learned metric map: shared by two legs, bypassed by one -----
    \node[trained, right=4mm of pf] (W)
      {learned map $W$\\$192\times192$\\[1pt]{\scriptsize no bias, identity init}\\{\scriptsize
       InfoNCE $\tau{=}0.07$}\\{\scriptsize $36{,}864$ params}};
    % ---- the three legs -------------------------------------------------
    \node[learnspace, right=3mm of W] (C)
      {\Cp\ content-$k$NN in $W$-space\\$\max_{u}\cos(Wz_v,Wz_u)$\\[1pt]{\scriptsize
       no parameters of its own}};
    \node[trained, above=7mm of C] (R)
      {\Rp\ projected GRU\\$h{=}256$, $394{,}944$ params};
    \node[lookup, below=7mm of C] (M)
      {$M$ first-order Markov counts\\$P(v\mid u)$, $u=$ last prefix item\\[1pt]{\scriptsize
       Dirichlet trust $n_u/(n_u{+}8)$,}\\{\scriptsize artist/genre back-off; $0$ params}};
    \draw[flow] (pf.east) -- (W.west);
    % ---- the frozen item table W acts ON: the same Z as Figure~\ref{fig:topo-gru} --
    \node[frozen, anchor=south east, minimum width=18mm] (Z)
      at ([yshift=6mm]W.north west)
      {frozen item latents $Z$\\{\scriptsize$19{,}402\times192$}};
    \draw[data] (Z.south east) -- node[edg, left, inner sep=1.2pt] {$z_i$} (W.north west);
    \draw[proj] (W.north east) -- (R.west);
    \draw[proj] (W.east) -- (C.west);
    \draw[data] (pf.south) |- node[tag, above, pos=0.55]
      {item \emph{ids} only ---\\the $M$ leg bypasses $Z$ and $W$} (M.west);
    % ---- the fixed combiner ---------------------------------------------
    \node[op, right=3mm of C] (cb)
      {z-normalise each leg\\over the candidate set,\\then average\\[1pt]{\scriptsize
       equal thirds ---}\\{\scriptsize\emph{no learned weights}}};
    \draw[flow] (R.east) -- node[edg, above, sloped, pos=0.55] {$s_{\Rp}$} (cb.north west);
    \draw[flow] (C.east) -- node[edg, above, inner sep=1pt] {$s_{\Cp}$} (cb.west);
    \draw[flow] (M.east) -- node[edg, above, sloped, pos=0.30] {$s_M$} (cb.south west);
    \node[op, right=2mm of cb] (rk) {rank\\$\to$ top-10};
    \draw[flow] (cb) -- (rk);
    \node[note] at ([yshift=-3.5mm]current bounding box.south west)
      {\textbf{Which legs are learned.} $W$ and the \Rp\ recurrence are
       trained (\num{431808} parameters between them); \Cp\ has no parameters
       of its own but retrieves inside $W$'s learned geometry, which is why it
       is drawn dashed-blue; $M$ is a table of training-set transition counts,
       drawn dashed-purple. The combiner is a fixed rule: each leg's
       full-vocabulary score vector is z-scored over the admissible candidates
       (every item minus those already in the prefix) and the three are
       averaged with hard-coded equal weights. The z-scoring is the
       load-bearing step --- it puts three incomparable score scales onto one
       per-session scale while \emph{preserving} relative magnitude, so a
       confidently-peaked leg can outvote two diffuse ones. Rank-only fusion,
       which discards that magnitude, loses by $-\num{0.029}$ CI$<$0.};
  \end{tikzpicture}
  \caption{\texttt{seq-blend} with \texttt{projection=true}, the champion:
    three full-vocabulary score vectors, z-normalised over the candidate set
    and averaged in fixed equal thirds. $W$ acts on the \emph{frozen item
    latents} $Z$ of Figure~\ref{fig:topo-gru} --- it warps the retrieval geometry,
    it does not read the prefix ids directly --- and it is shared by two of the
    three legs and bypassed by the third; that is the whole mechanism. $W$ is fit by in-batch InfoNCE on consecutive within-train
    pairs, which rebuilds the retrieval metric from ``reconstruct the item
    matrix'' into ``what follows what''; both legs that retrieve by cosine in
    that space improve --- the parameter-free content-$k$NN leg more than
    doubles, \num{0.0685} to \num{0.1516}, and the recurrent leg rises from
    \num{0.1230} to $\approx$\num{0.172} --- while the item-ID lookup is
    untouched at \num{0.106918}. The two projected-leg standalone figures come
    from the off-server driver of the 2026-07-18 campaign and have no server
    run of their own; only the $M$ leg's number is a canonical-split run.
    Together the three legs reach Recall@10 \num{0.21174} (\num{303} of
    \num{1431}), paired $\Delta$ over the unprojected three-leg blend
    $+\num{0.0259}$ $[+0.012, +0.041]$, a CI entirely above zero and robust
    across three seeds. Visual convention as in Figure~\ref{fig:topo-gru}.}
  \label{fig:topo-blend}
\end{figure}


% ===========================================================================
% 6. seq-stacker --- the learned combiner
% ===========================================================================
\begin{figure}[ht]
  \centering
  \begin{tikzpicture}[topo]
    \node[frozen, minimum width=18mm, minimum height=11mm] (pf)
      {session\\prefix\\$i_1\dots i_t$};
    \node[trained, right=6mm of pf] (R)
      {$R$ GRU $h{=}256$\\[1pt]{\scriptsize trained separately}};
    \node[lookup, above=6mm of R] (M) {$M$ Markov counts};
    \node[trained, below=6mm of R] (A)
      {$A$ pooled MLP\\[1pt]{\scriptsize trained separately}};
    \draw[flow] (pf.east) -- (M.west);
    \draw[flow] (pf.east) -- (R.west);
    \draw[flow] (pf.east) -- (A.west);
    % ---- the per-candidate feature row ----------------------------------
    \node[op, right=6mm of R] (fx)
      {per-candidate feature row\\$[\,s_M,\ s_R,\ s_A,\ \log(1{+}\text{plays}),$\\$\text{recency},\
       \text{artist match},\ \text{genre match}\,]$};
    \draw[flow] (M.east) -- (fx.north west);
    \draw[flow] (R.east) -- (fx.west);
    \draw[flow] (A.east) -- (fx.south west);
    \node[refuted, right=5mm of fx] (xgb)
      {XGBoost meta-learner\\$200$ rounds, depth $6$\\[1pt]{\scriptsize\texttt{binary:logistic}}};
    \draw[flow] (fx) -- (xgb);
    \node[op, below=9mm of xgb] (rk)
      {score all $19{,}402$ items\\$\to$ top-10};
    \draw[flow] (xgb) -- (rk);
    % ---- the two failure mechanisms, marked on the diagram --------------
    \node[refuted, text width=56mm, align=flush left, anchor=north west] (tr)
      at ([yshift=-10mm]fx.south west)
      {\textbf{Training rows}: $1$ positive $+$ $50$ popularity-sampled
       negatives per train session, i.e.\ \num{5723} positives, against the
       \num{73632} teacher-forced steps its own $R$ leg trains on
       ($12.9\times$ less signal) --- and a decision surface calibrated on
       $51$ candidates ($0.26\%$ of the catalogue) is then asked to rank all
       \num{19402}.};
    \draw[grad] (tr.north east) -- (xgb.south west);
    \node[note] at ([yshift=-3.5mm]current bounding box.south west)
      {\textbf{Why the shape matters.} A $200$-round depth-$6$ booster is
       piecewise \emph{constant}, so ranking \num{19402} candidates with it
       produces large tied blocks that the harness breaks by item index. The
       fixed z-blend of Figure~\ref{fig:topo-blend} is instead a strictly
       monotone continuous function of the leg scores: it preserves each leg's
       internal ordering exactly and can only trade the legs off against one
       another.};
  \end{tikzpicture}
  \caption{\texttt{seq-stacker}, the learned counterpart of
    Figure~\ref{fig:topo-blend}'s fixed rule: the same separately-trained base
    legs, but their per-candidate scores become columns of a feature row
    alongside four cheap candidate features, and a gradient-boosted tree
    learns the combination. It earns a diagram because its failure is
    informative --- with the \texttt{markov}, \texttt{gru} and \texttt{ann}
    legs it reaches Recall@10 \num{0.085255} (\num{122} of \num{1431}),
    \emph{below its own GRU base leg} at \num{0.12299}, paired $\Delta$
    $-\num{0.0377}$ $[-0.056, -0.020]$ CI$<$0. The red callout names the two
    mechanisms: supervision starvation, since the meta-learner gets one
    labelled event per session where its base leg gets one per prefix
    position, and a train/serve candidate mismatch. The booster is drawn solid
    because it is trained and red because it is refuted; the feature row and
    the ranking step are dashed because they hold no parameters. Visual
    convention as in Figure~\ref{fig:topo-gru}. This is one of four
    independent learned-combiner nulls, the others being reciprocal-rank
    fusion, oracle weights fitted on the test labels, and the jointly-trained
    dual tower of Figure~\ref{fig:topo-dualgru}.}
  \label{fig:topo-stacker}
\end{figure}
, can use the same
%  visual language.  \input from docs/ (compilation cwd):
%
%      % ===========================================================================
%  topo-style.tex --- the shared TikZ style block for the next-track topology
%  diagrams AND for the evaluation-protocol diagram in the body of
%  model-comparison.tex.  Extracted from topologies.tex so that the protocol
%  figure, which appears BEFORE % ===========================================================================
%  topologies.tex --- TikZ network-topology diagrams for the next-track
%  MODEL-COMPARISON document.  Not a standalone document: \input it from the
%  body of the main .tex, which must be compiled from docs/:
%
%      \input{figures/topologies}
%
%  (\graphicspath{{figures/}} governs \includegraphics only, NOT \input, so the
%   figures/ prefix is spelled out above.)
%
%  ------------------------------------------------------------------------
%  REQUIRED PREAMBLE LINES in the main document --- add these two lines after
%  \usepackage{graphicx}, before \begin{document}:
%
%      \usepackage{tikz}
%      \usetikzlibrary{positioning,arrows.meta,fit,backgrounds,calc}
%
%  Nothing else is required: xcolor arrives with tikz, and every colour and
%  every style used below is defined in this file.  \dsid, \Rp and \Cp are
%  \providecommand-guarded here so the file also compiles inside a bare
%  wrapper; when the main document defines them, its definitions win.
%
%  PACKAGES DELIBERATELY *NOT* REQUIRED: amssymb is absent from the house
%  preamble, so nothing below uses \mathbb.  Every picture is laid out with
%  RELATIVE positioning (right=/above=/below= of ...), never with hard
%  coordinates, so no node can collide with its neighbour if the font or the
%  wording changes; the per-figure notes and the legend are anchored to
%  `current bounding box', so they always land underneath the diagram.
%  ------------------------------------------------------------------------
%
%  LABELS DEFINED IN THIS FILE (in file order)
%      fig:topo-gru       seq-nexttrack --- the baseline; ALSO CARRIES THE LEGEND
%      fig:topo-ann       seq-ann       --- pooled feed-forward twin
%      fig:topo-dualgru   seq-dualgru   --- two towers + fusion head + pre-MLP variant
%      fig:topo-embed     seq-embed     --- learned, weight-tied item table (DESIGN ONLY)
%      fig:topo-blend     seq-blend     --- the champion, three legs
%      fig:topo-stacker   seq-stacker   --- learned XGBoost combiner
%
%  VISUAL LANGUAGE (the primary axis is the LINE, not the colour, so it
%  survives greyscale printing; stated for the reader in the caption of
%  fig:topo-gru and re-pointed-at by every other caption)
%      SOLID  border  = TRAINED component (parameters fitted by gradient descent)
%      DASHED border  = FROZEN or PARAMETER-FREE component
%      blue   = the on-record default path
%      purple dashed  = non-parametric table lookup over the dataset
%      dashed blue    = parameter-free, but operating inside a LEARNED geometry
%      green  = variant favoured by the record
%      red    = variant refuted by the record
% ===========================================================================

\input{figures/topo-style}



% ===========================================================================
% 1. seq-nexttrack --- the baseline (and the legend)
% ===========================================================================
\begin{figure}[ht]
  \centering
  \begin{tikzpicture}[topo]
    % ---- the frozen item-latent table, read TWICE ------------------------
    \node[frozen, minimum width=20mm, minimum height=12mm] (tab)
      {frozen item\\latents $Z$\\[1pt]{\scriptsize$19{,}402\times192$}};
    % ---- the prefix, in order (one column per step) ----------------------
    \node[frozen, minimum width=12mm, right=7mm of tab]  (x1) {$x_1$};
    \node[frozen, minimum width=12mm, right=4mm of x1]   (x2) {$x_2$};
    \node[tag, right=2.5mm of x2]                        (xd) {$\cdots$};
    \node[frozen, minimum width=12mm, right=2.5mm of xd] (x3) {$x_t$};
    \begin{scope}[on background layer]
      \node[grp, fit=(x1)(x3)] (pfx) {};
    \end{scope}
    \draw[data] (tab.east) -- (pfx.west);
    \node[tag, below=1.2mm of pfx] (pfxlab)
      {prefix latents --- \emph{order preserved}};
    % ---- the recurrence --------------------------------------------------
    \node[trained, minimum width=12mm, above=7mm of x1] (g1) {GRU};
    \node[trained, minimum width=12mm, above=7mm of x2] (g2) {GRU};
    \node[tag] (gd) at (g1 -| xd) {$\cdots$};
    \node[trained, minimum width=12mm, above=7mm of x3] (g3) {GRU};
    \draw[flow] (x1) -- (g1);
    \draw[flow] (x2) -- (g2);
    \draw[flow] (x3) -- (g3);
    \draw[flow] (g1) -- node[edg, above, inner sep=0.7pt] {$h_t$} (g2);
    \draw[flow] (g2) -- (gd);
    \draw[flow] (gd) -- (g3);
    % ---- head and retrieval ---------------------------------------------
    \node[trained, right=6mm of g3] (head)
      {linear head\\{\scriptsize$256\to192$}};
    \node[op, right=7mm of head] (rk)
      {rank the catalogue by $\cos(\hat z, z_v)$\\{\scriptsize prefix excluded $\to$ top-10}};
    \draw[flow] (g3) -- (head);
    \draw[flow] (head) -- node[edg, above, inner sep=1pt] {$\hat z$} (rk);
    % ---- the SAME frozen table is the retrieval target -------------------
    \coordinate (c1) at ([yshift=-3mm]tab.south |- current bounding box.south);
    \coordinate (c2) at (rk.south |- c1);
    \draw[data] (tab.south) -- (c1);
    \draw[data] (c1) -- node[tag, above]
      {the \emph{same} frozen table is the retrieval target} (c2);
    \draw[data] (c2) -- (rk.south);
    % ---- supervision note ------------------------------------------------
    \node[note] (nt) at ([yshift=-3.5mm]current bounding box.south west)
      {\textbf{Supervision --- per-step teacher forcing.} The head is shared
       across positions and must predict the next item's latent at
       \emph{every} prefix position, so one session of length $L$ yields
       $L-1$ examples: $73{,}632$ (prefix, next-item) pairs on the train
       split. The loss is masked in-batch InfoNCE at $\tau=0.07$ --- each
       step's true next latent must outscore the other steps' targets in the
       same minibatch, which supplies negatives for free.};
    % ---- LEGEND (two rows so it cannot widen the picture) ----------------
    \node[trained, swatch, anchor=north west] (l1)
      at ([yshift=-3mm]current bounding box.south west) {};
    \node[tag, right=1.1mm of l1] (t1) {trained (solid)};
    \node[frozen, swatch, right=3.4mm of t1] (l2) {};
    \node[tag, right=1.1mm of l2] (t2) {frozen / parameter-free (dashed)};
    \node[lookup, swatch, right=3.4mm of t2] (l3) {};
    \node[tag, right=1.1mm of l3] (t3) {table lookup};
    \node[learnspace, swatch, anchor=north west] (l4)
      at ([yshift=-1.6mm]l1.south west) {};
    \node[tag, right=1.1mm of l4] (t4) {parameter-free, learned space};
    \node[won, swatch, right=3.4mm of t4] (l5) {};
    \node[tag, right=1.1mm of l5] (t5) {favoured variant};
    \node[refuted, swatch, right=3.4mm of t5] (l6) {};
    \node[tag, right=1.1mm of l6] (t6) {refuted variant};
  \end{tikzpicture}
  \caption{\texttt{seq-nexttrack}, the single-model baseline every other
    topology in this document is measured against: a frozen item space, a
    recurrence over the \emph{ordered} prefix, a linear head, and cosine
    retrieval back into the \emph{same} frozen space. Nothing about the item
    geometry is learned here --- the $19{,}402\times192$ latent table (offline
    PCA-192 on \dsid{seq-20260715-131139}) is read twice, once to embed the
    prefix and once as the retrieval target, and only the GRU and its linear
    head carry parameters (\num{394944} at $h=256$). \textbf{The visual
    language used throughout this document is fixed here.} A \emph{solid}
    border marks a trained component and a \emph{dashed} border a frozen or
    parameter-free one; that distinction rides on the line style, not the
    colour, so it survives greyscale. Colour then carries the verdict ---
    blue: the on-record default path; purple dashed: a non-parametric table
    lookup over the dataset; dashed blue: parameter-free but operating inside
    a \emph{learned} geometry; green: a variant the record favours; red: a
    variant the record refutes. On the canonical split this configuration
    scores Recall@10 \num{0.12299} (\num{176} of \num{1431}), reproduced
    bit-identically eleven times.}
  \label{fig:topo-gru}
\end{figure}


% ===========================================================================
% 2. seq-ann --- the pooled feed-forward twin
% ===========================================================================
\begin{figure}[ht]
  \centering
  \begin{tikzpicture}[topo]
    \node[frozen, minimum width=20mm, minimum height=12mm] (tab)
      {frozen item\\latents $Z$\\[1pt]{\scriptsize$19{,}402\times192$}};
    \node[frozen, minimum width=12mm, right=7mm of tab]  (x1) {$x_1$};
    \node[frozen, minimum width=12mm, right=4mm of x1]   (x2) {$x_2$};
    \node[tag, right=2.5mm of x2]                        (xd) {$\cdots$};
    \node[frozen, minimum width=12mm, right=2.5mm of xd] (x3) {$x_t$};
    \begin{scope}[on background layer]
      \node[grp, fit=(x1)(x3)] (pfx) {};
    \end{scope}
    \draw[data] (tab.east) -- (pfx.west);
    % ---- the two parameter-free summaries of the same prefix -------------
    \node[op, right=8mm of pfx, yshift=10mm] (mp)
      {causal mean-pool\\$\frac{1}{t}\sum_{i\le t}x_i$\\[1pt]{\scriptsize\emph{permutation-invariant}}};
    \node[op, right=8mm of pfx, yshift=-10mm] (lt)
      {last item\\$x_t$\\[1pt]{\scriptsize the only recency cue}};
    \draw[flow] (pfx.east) -- (mp.west);
    \draw[flow] (pfx.east) -- (lt.west);
    % ---- the MLP ---------------------------------------------------------
    \node[trained, right=5mm of mp, yshift=-10mm] (mlp)
      {MLP\\$384\to256\to256\to192$\\[1pt]{\scriptsize ReLU $+$ dropout}};
    \draw[flow] (mp.east) -- (mlp.north west);
    \draw[flow] (lt.east) -- (mlp.south west);
    \node[edg] at ([xshift=-4.2mm]mlp.west) {$\|$};
    % ---- retrieval, placed below to keep the picture inside the text ------
    \node[op, below=9mm of mlp] (rk)
      {$\hat z\to$ cosine rank\\{\scriptsize prefix excluded $\to$ top-10}};
    \draw[flow] (mlp) -- (rk);
    \coordinate (c1) at ([yshift=-3mm]tab.south |- current bounding box.south);
    \coordinate (c2) at (rk.south |- c1);
    \draw[data] (tab.south) -- (c1);
    \draw[data] (c1) -- node[tag, above]
      {the same frozen table is the retrieval target} (c2);
    \draw[data] (c2) -- (rk.south);
    \node[note] at ([yshift=-3.5mm]current bounding box.south west)
      {\textbf{The same information, with the order discarded.} Both summaries
       are strictly causal, so the objective, the padded batching and the
       early stopping are reused \emph{verbatim} from
       Figure~\ref{fig:topo-gru} --- but a mean weights the track from twenty
       steps ago exactly like the one before last, and cannot represent ``A
       then B'' differently from ``B then A''. Only the special-cased $x_t$ tap
       carries any notion of recency, and the concatenation ($\|$) is this
       model's entire account of sequence.};
  \end{tikzpicture}
  \caption{\texttt{seq-ann}, the GRU's non-recurrent twin: the recurrence of
    Figure~\ref{fig:topo-gru} is replaced by two parameter-free summaries of
    the same prefix --- a causal mean-pool concatenated with the last item's
    latent --- feeding a plain three-layer ReLU MLP. Everything else is
    literally the same code (\texttt{seq\_ann.py} imports
    \texttt{make\_batches} and \texttt{step\_loss} from
    \texttt{seq\_nexttrack}), which is what makes this pair a clean test of
    the recurrence itself rather than of a training recipe. Visual convention
    as in Figure~\ref{fig:topo-gru}: the two pooling operators are dashed
    because they hold no parameters, and only the MLP is trained. The
    defensible comparison is same-space: on \dsid{seq-20260715-030509}
    (PCA-128) the pooled twin scores Recall@10 \num{0.0901} (\num{129} of
    \num{1431}) against the GRU's \num{0.1146} (\num{164}), paired $\Delta$
    $-\num{0.0245}$ $[-0.039, -0.011]$, a CI \emph{entirely below zero}.
    \texttt{seq-ann} has no run at all on the canonical PCA-192 split, so it
    never belongs in the same column as the \num{0.12299} bar.}
  \label{fig:topo-ann}
\end{figure}


% ===========================================================================
% 3. seq-dualgru --- two towers, a fusion head, and the pre-MLP variant
% ===========================================================================
\begin{figure}[ht]
  \centering
  \begin{tikzpicture}[topo]
    \node[frozen, minimum width=19mm, minimum height=11mm] (pf)
      {prefix\\latents\\$x_1\dots x_t$};
    % ---- the three configurable causal views ----------------------------
    \node[op, right=7mm of pf] (v2) {\texttt{delta}\\$x_t-x_{t-1}$};
    \node[op, above=2.6mm of v2] (v1) {\texttt{latent}\\$x_t$};
    \node[op, below=2.6mm of v2] (v3) {\texttt{cummean}\\$\frac{1}{t}\sum_{i\le t}x_i$};
    \begin{scope}[on background layer]
      \node[grp, fit=(v1)(v3)] (vg) {};
    \end{scope}
    \draw[data] (pf.east) -- (vg.west);
    % ---- two independent towers -----------------------------------------
    \node[trained, right=8mm of vg, yshift=9mm] (ta)
      {tower A on \texttt{view\_a}\\GRU $h_a{=}256$\\[1pt]{\scriptsize own weights}};
    \node[trained, right=8mm of vg, yshift=-9mm] (tb)
      {tower B on \texttt{view\_b}\\GRU $h_b{=}256$\\[1pt]{\scriptsize own weights}};
    \draw[flow] (vg.east) -- (ta.west);
    \draw[flow] (vg.east) -- (tb.west);
    % ---- concat + the two fusion depths ---------------------------------
    \node[op, right=7mm of ta, yshift=-9mm] (cc)
      {concat\\$h_a\|h_b$\\[1pt]{\scriptsize$512$}};
    \draw[flow] (ta.east) -- (cc.north west);
    \draw[flow] (tb.east) -- (cc.south west);
    \node[won, right=7mm of cc, yshift=9mm] (f0)
      {\texttt{fusion\_layers}${}=0$\\$512\to192$, linear};
    \node[refuted, right=7mm of cc, yshift=-9mm] (f1)
      {\texttt{fusion\_layers}${}=1$\\$512\to256\to192$\\[1pt]{\scriptsize$+$ ReLU $+$ dropout}};
    \draw[flow] (cc.north east) -- (f0.west);
    \draw[flow] (cc.south east) -- (f1.west);
    % ---- the refuted per-step pre-MLP variant ---------------------------
    \node[refuted, below=12mm of tb] (pb)
      {\emph{variant} --- per-step pre-MLP\\on a tower:
       Linear$(192{\to}256)$\\$+$ ReLU $+$ Dropout};
    \draw[variant] (vg.south) |- (pb.west);
    \draw[variant] (pb.north) -- (tb.south);
    \node[note] at ([yshift=-3.5mm]current bounding box.south west)
      {\textbf{Both heads emit the same $192$-d $\hat z$ and retrieve exactly
       as in Figure~\ref{fig:topo-gru}.} The towers share nothing --- separate
       modules, separate parameters --- and their per-step hidden states are
       concatenated, so the combiner sees the towers' \emph{representations}
       rather than their scores and gradients flow back into both. There is no
       bidirectional mode by construction: a backward pass would let step $t$
       read item $t+1$ and copy its own teacher-forcing target. All three views
       are strictly causal for the same reason, and so is the pre-MLP, which is
       applied \emph{timestep-wise}; it is drawn solid because it is trained,
       red because it is refuted.};
  \end{tikzpicture}
  \caption{\texttt{seq-dualgru}: two recurrent towers with entirely separate
    weights read the prefix in parallel and a head fuses their hidden states
    --- the strongest available form of ``train the combination end-to-end so
    the parts co-adapt''. Each tower's input is a configurable causal view
    (\texttt{view\_a} / \texttt{view\_b}); the default is
    \texttt{latent}/\texttt{latent}, i.e.\ two independent GRUs over the same
    stream. Two of the boxes are coloured by the record rather than by the
    design. Removing the MLP fusion head (\texttt{fusion\_layers}${}=0$,
    green) \emph{beats} keeping it in three independent measurements, most
    recently by $+\num{0.02655}$ $[+0.01398, +0.03913]$ CI$>$0 --- and with
    some $9$--$10\%$ \emph{fewer} parameters, so the linear readout is the
    family default; and the
    per-step pre-MLP (red) loses the decisive single-tower ablation, Recall@10
    \num{0.10901} (\num{156} of \num{1431}) against the baseline's \num{176},
    paired $\Delta$ $-\num{0.01398}$ $[-0.02657, -0.00140]$ CI$<$0. Visual
    convention as in Figure~\ref{fig:topo-gru}. The best dual arm reaches
    \num{0.11880} and merely \emph{ties} the single-GRU bar, as does a
    $0.098\%$ parameter-matched width control --- so capacity is not the
    explanation.}
  \label{fig:topo-dualgru}
\end{figure}


% ===========================================================================
% 4. seq-embed --- a learned, weight-tied item table (DESIGN ONLY)
% ===========================================================================
\begin{figure}[ht]
  \centering
  \begin{tikzpicture}[topo]
    \node[frozen, minimum width=20mm, minimum height=11mm] (ids)
      {session prefix\\item \emph{ids}\\$i_1\dots i_t$};
    \node[trained, right=10mm of ids, minimum height=11mm] (E)
      {learned item\\table $E$\\[1pt]{\scriptsize$19{,}402\times192$}};
    \node[trained, right=6mm of E] (gru)
      {GRU $h{=}256$\\{\scriptsize$+$ dropout}};
    \node[trained, right=6mm of gru] (hd)
      {linear head\\{\scriptsize$256\to192$}};
    \node[op, right=6mm of hd] (rk)
      {$\cos(\hat z, E_v)$ over all\\$19{,}402$ rows $\to$ top-10};
    \draw[flow] (ids) -- node[edg, above, inner sep=1pt] {gather} (E);
    \draw[flow] (E) -- (gru);
    \draw[flow] (gru) -- (hd);
    \draw[flow] (hd) -- node[edg, above, inner sep=1pt] {$\hat z$} (rk);
    % ---- the three modes -------------------------------------------------
    \node[op, text width=64mm, align=flush left, anchor=south west] (md)
      at ([xshift=4mm,yshift=8mm]E.north east)
      {\texttt{embed\_mode} --- what the table $E$ actually is:\\
       \textbf{\texttt{frozen}} $E=Z$, nothing trained (the reproduction gate)\\
       \textbf{\texttt{residual}} $E=Z+\Delta E$: frozen content vector $+$ a
       trained correction, so a cold item still has a position\\
       \textbf{\texttt{free}} $E$ fully trained, optionally randomly initialised};
    \draw[flow] (md.west) -- (E.north east);
    % ---- weight tying: the same table is the retrieval target -----------
    \coordinate (c1) at ([yshift=-8mm]E.south);
    \coordinate (c2) at (rk.south |- c1);
    \draw[tie] (E.south) -- (c1);
    \draw[tie] (c1) -- node[tag, above, text=topogood]
      {weight tying --- the \emph{same} table is\\the input embedding
       \emph{and} the retrieval target} (c2);
    \draw[tie] (c2) -- (rk.south);
    % ---- the gradient loop that makes the geometry trainable ------------
    \coordinate (t1) at ([yshift=5mm]rk.north |- current bounding box.north);
    \coordinate (t2) at (E.north |- t1);
    \draw[grad] (rk.north) -- (t1);
    \draw[grad] (t1) -- node[tag, above, text=topobad]
      {$\nabla$ --- one InfoNCE step moves the query \emph{and} every target:
       the geometry itself is trainable} (t2);
    \draw[grad] (t2) -- (E.north);
    \node[note] at ([yshift=-3.5mm]current bounding box.south west)
      {\textbf{Why the loop is the whole point.} In
       Figures~\ref{fig:topo-gru}--\ref{fig:topo-dualgru} the retrieval space
       is fixed and only the map into it is learned; here the space is a
       parameter, shaped by next-track adjacency rather than inherited from a
       reconstruction objective. That also creates a collapse hazard absent
       everywhere else: with a learned table a plain cosine objective has a
       trivial optimum at ``send every item to one point'', so only a
       contrastive loss is admissible and \texttt{loss=cosine} is rejected
       outright for the trained modes.};
  \end{tikzpicture}
  \caption{\texttt{seq-embed}, the one axis the record has not tested:
    \textbf{a design, not a result}. No run of this predictor exists on any
    dataset, it is absent from the server's registered predictor list, and
    nothing here is a measured score. The item table $E$ is used twice, as the
    GRU's input embedding and as the retrieval target (weight tying, the
    GRU4Rec/SASRec formulation), so a single gradient step moves the query and
    all the candidates together --- the return arrow along the top is what
    makes the geometry trainable, and it is exactly what every other topology
    in this document lacks. The three \texttt{embed\_mode} settings span the
    space from a reproduction gate to a fully free table. What \emph{is}
    measured is the data budget that bounds them: \num{73632} training steps
    to fit \num{15816} item vectors, a median per-item training frequency of
    \num{1}, only \num{3294} items occurring five or more times, and
    \num{788} of \num{1431} test targets never appearing in any training
    session --- so a \texttt{free} table cannot retrieve the cold
    \SI{55}{\percent} at all and is capped at Recall@10 \num{0.4493}, whereas
    \texttt{residual} keeps the content vector and with it the cold reach.
    Visual convention as in Figure~\ref{fig:topo-gru}.}
  \label{fig:topo-embed}
\end{figure}


% ===========================================================================
% 5. seq-blend --- the champion
% ===========================================================================
\begin{figure}[ht]
  \centering
  \begin{tikzpicture}[topo]
    \node[frozen, minimum width=18mm, minimum height=11mm] (pf)
      {session\\prefix\\$i_1\dots i_t$};
    % ---- the learned metric map: shared by two legs, bypassed by one -----
    \node[trained, right=4mm of pf] (W)
      {learned map $W$\\$192\times192$\\[1pt]{\scriptsize no bias, identity init}\\{\scriptsize
       InfoNCE $\tau{=}0.07$}\\{\scriptsize $36{,}864$ params}};
    % ---- the three legs -------------------------------------------------
    \node[learnspace, right=3mm of W] (C)
      {\Cp\ content-$k$NN in $W$-space\\$\max_{u}\cos(Wz_v,Wz_u)$\\[1pt]{\scriptsize
       no parameters of its own}};
    \node[trained, above=7mm of C] (R)
      {\Rp\ projected GRU\\$h{=}256$, $394{,}944$ params};
    \node[lookup, below=7mm of C] (M)
      {$M$ first-order Markov counts\\$P(v\mid u)$, $u=$ last prefix item\\[1pt]{\scriptsize
       Dirichlet trust $n_u/(n_u{+}8)$,}\\{\scriptsize artist/genre back-off; $0$ params}};
    \draw[flow] (pf.east) -- (W.west);
    % ---- the frozen item table W acts ON: the same Z as Figure~\ref{fig:topo-gru} --
    \node[frozen, anchor=south east, minimum width=18mm] (Z)
      at ([yshift=6mm]W.north west)
      {frozen item latents $Z$\\{\scriptsize$19{,}402\times192$}};
    \draw[data] (Z.south east) -- node[edg, left, inner sep=1.2pt] {$z_i$} (W.north west);
    \draw[proj] (W.north east) -- (R.west);
    \draw[proj] (W.east) -- (C.west);
    \draw[data] (pf.south) |- node[tag, above, pos=0.55]
      {item \emph{ids} only ---\\the $M$ leg bypasses $Z$ and $W$} (M.west);
    % ---- the fixed combiner ---------------------------------------------
    \node[op, right=3mm of C] (cb)
      {z-normalise each leg\\over the candidate set,\\then average\\[1pt]{\scriptsize
       equal thirds ---}\\{\scriptsize\emph{no learned weights}}};
    \draw[flow] (R.east) -- node[edg, above, sloped, pos=0.55] {$s_{\Rp}$} (cb.north west);
    \draw[flow] (C.east) -- node[edg, above, inner sep=1pt] {$s_{\Cp}$} (cb.west);
    \draw[flow] (M.east) -- node[edg, above, sloped, pos=0.30] {$s_M$} (cb.south west);
    \node[op, right=2mm of cb] (rk) {rank\\$\to$ top-10};
    \draw[flow] (cb) -- (rk);
    \node[note] at ([yshift=-3.5mm]current bounding box.south west)
      {\textbf{Which legs are learned.} $W$ and the \Rp\ recurrence are
       trained (\num{431808} parameters between them); \Cp\ has no parameters
       of its own but retrieves inside $W$'s learned geometry, which is why it
       is drawn dashed-blue; $M$ is a table of training-set transition counts,
       drawn dashed-purple. The combiner is a fixed rule: each leg's
       full-vocabulary score vector is z-scored over the admissible candidates
       (every item minus those already in the prefix) and the three are
       averaged with hard-coded equal weights. The z-scoring is the
       load-bearing step --- it puts three incomparable score scales onto one
       per-session scale while \emph{preserving} relative magnitude, so a
       confidently-peaked leg can outvote two diffuse ones. Rank-only fusion,
       which discards that magnitude, loses by $-\num{0.029}$ CI$<$0.};
  \end{tikzpicture}
  \caption{\texttt{seq-blend} with \texttt{projection=true}, the champion:
    three full-vocabulary score vectors, z-normalised over the candidate set
    and averaged in fixed equal thirds. $W$ acts on the \emph{frozen item
    latents} $Z$ of Figure~\ref{fig:topo-gru} --- it warps the retrieval geometry,
    it does not read the prefix ids directly --- and it is shared by two of the
    three legs and bypassed by the third; that is the whole mechanism. $W$ is fit by in-batch InfoNCE on consecutive within-train
    pairs, which rebuilds the retrieval metric from ``reconstruct the item
    matrix'' into ``what follows what''; both legs that retrieve by cosine in
    that space improve --- the parameter-free content-$k$NN leg more than
    doubles, \num{0.0685} to \num{0.1516}, and the recurrent leg rises from
    \num{0.1230} to $\approx$\num{0.172} --- while the item-ID lookup is
    untouched at \num{0.106918}. The two projected-leg standalone figures come
    from the off-server driver of the 2026-07-18 campaign and have no server
    run of their own; only the $M$ leg's number is a canonical-split run.
    Together the three legs reach Recall@10 \num{0.21174} (\num{303} of
    \num{1431}), paired $\Delta$ over the unprojected three-leg blend
    $+\num{0.0259}$ $[+0.012, +0.041]$, a CI entirely above zero and robust
    across three seeds. Visual convention as in Figure~\ref{fig:topo-gru}.}
  \label{fig:topo-blend}
\end{figure}


% ===========================================================================
% 6. seq-stacker --- the learned combiner
% ===========================================================================
\begin{figure}[ht]
  \centering
  \begin{tikzpicture}[topo]
    \node[frozen, minimum width=18mm, minimum height=11mm] (pf)
      {session\\prefix\\$i_1\dots i_t$};
    \node[trained, right=6mm of pf] (R)
      {$R$ GRU $h{=}256$\\[1pt]{\scriptsize trained separately}};
    \node[lookup, above=6mm of R] (M) {$M$ Markov counts};
    \node[trained, below=6mm of R] (A)
      {$A$ pooled MLP\\[1pt]{\scriptsize trained separately}};
    \draw[flow] (pf.east) -- (M.west);
    \draw[flow] (pf.east) -- (R.west);
    \draw[flow] (pf.east) -- (A.west);
    % ---- the per-candidate feature row ----------------------------------
    \node[op, right=6mm of R] (fx)
      {per-candidate feature row\\$[\,s_M,\ s_R,\ s_A,\ \log(1{+}\text{plays}),$\\$\text{recency},\
       \text{artist match},\ \text{genre match}\,]$};
    \draw[flow] (M.east) -- (fx.north west);
    \draw[flow] (R.east) -- (fx.west);
    \draw[flow] (A.east) -- (fx.south west);
    \node[refuted, right=5mm of fx] (xgb)
      {XGBoost meta-learner\\$200$ rounds, depth $6$\\[1pt]{\scriptsize\texttt{binary:logistic}}};
    \draw[flow] (fx) -- (xgb);
    \node[op, below=9mm of xgb] (rk)
      {score all $19{,}402$ items\\$\to$ top-10};
    \draw[flow] (xgb) -- (rk);
    % ---- the two failure mechanisms, marked on the diagram --------------
    \node[refuted, text width=56mm, align=flush left, anchor=north west] (tr)
      at ([yshift=-10mm]fx.south west)
      {\textbf{Training rows}: $1$ positive $+$ $50$ popularity-sampled
       negatives per train session, i.e.\ \num{5723} positives, against the
       \num{73632} teacher-forced steps its own $R$ leg trains on
       ($12.9\times$ less signal) --- and a decision surface calibrated on
       $51$ candidates ($0.26\%$ of the catalogue) is then asked to rank all
       \num{19402}.};
    \draw[grad] (tr.north east) -- (xgb.south west);
    \node[note] at ([yshift=-3.5mm]current bounding box.south west)
      {\textbf{Why the shape matters.} A $200$-round depth-$6$ booster is
       piecewise \emph{constant}, so ranking \num{19402} candidates with it
       produces large tied blocks that the harness breaks by item index. The
       fixed z-blend of Figure~\ref{fig:topo-blend} is instead a strictly
       monotone continuous function of the leg scores: it preserves each leg's
       internal ordering exactly and can only trade the legs off against one
       another.};
  \end{tikzpicture}
  \caption{\texttt{seq-stacker}, the learned counterpart of
    Figure~\ref{fig:topo-blend}'s fixed rule: the same separately-trained base
    legs, but their per-candidate scores become columns of a feature row
    alongside four cheap candidate features, and a gradient-boosted tree
    learns the combination. It earns a diagram because its failure is
    informative --- with the \texttt{markov}, \texttt{gru} and \texttt{ann}
    legs it reaches Recall@10 \num{0.085255} (\num{122} of \num{1431}),
    \emph{below its own GRU base leg} at \num{0.12299}, paired $\Delta$
    $-\num{0.0377}$ $[-0.056, -0.020]$ CI$<$0. The red callout names the two
    mechanisms: supervision starvation, since the meta-learner gets one
    labelled event per session where its base leg gets one per prefix
    position, and a train/serve candidate mismatch. The booster is drawn solid
    because it is trained and red because it is refuted; the feature row and
    the ranking step are dashed because they hold no parameters. Visual
    convention as in Figure~\ref{fig:topo-gru}. This is one of four
    independent learned-combiner nulls, the others being reciprocal-rank
    fusion, oracle weights fitted on the test labels, and the jointly-trained
    dual tower of Figure~\ref{fig:topo-dualgru}.}
  \label{fig:topo-stacker}
\end{figure}
, can use the same
%  visual language.  \input from docs/ (compilation cwd):
%
%      % ===========================================================================
%  topo-style.tex --- the shared TikZ style block for the next-track topology
%  diagrams AND for the evaluation-protocol diagram in the body of
%  model-comparison.tex.  Extracted from topologies.tex so that the protocol
%  figure, which appears BEFORE \input{figures/topologies}, can use the same
%  visual language.  \input from docs/ (compilation cwd):
%
%      \input{figures/topo-style}
%
%  Idempotent: guarded, so a second \input is a no-op.  Requires
%      \usepackage{tikz}
%      \usetikzlibrary{positioning,arrows.meta,fit,backgrounds,calc}
% ===========================================================================
\ifdefined\nexttracktopostyle\else
\newcommand{\nexttracktopostyle}{}

\providecommand{\dsid}[1]{\texttt{\small #1}}
\providecommand{\Rp}{\ensuremath{R'}}
\providecommand{\Cp}{\ensuremath{C'}}

% --- palette: the same hex values as docs/figures/make_figures.py ----------
\definecolor{topoacc}{HTML}{2563EB}   % ACCENT  --- default / on-record path
\definecolor{topogood}{HTML}{16A34A}  % GOOD    --- favoured variant
\definecolor{topobad}{HTML}{DC2626}   % BAD     --- refuted variant
\definecolor{topomut}{HTML}{9CA3AF}   % MUTED   --- frozen, parameter-free
\definecolor{topolook}{HTML}{9333EA}  % seq-markov purple --- table lookup
\definecolor{topoink}{HTML}{374151}   % annotation ink

\tikzset{
  topo/.style={
    font=\footnotesize,
    every node/.append style={align=center},
  },
  box/.style={rounded corners=1.4pt, inner xsep=3.2pt, inner ysep=2.8pt,
              minimum height=6.4mm, text=topoink},
  dsh/.style={dash pattern=on 2pt off 1.4pt},
  % --- component styles ---------------------------------------------------
  trained/.style={box, draw=topoacc, line width=0.9pt, fill=topoacc!7},
  frozen/.style={box, dsh, draw=topomut, line width=0.8pt, fill=topomut!14},
  lookup/.style={box, dsh, draw=topolook, line width=0.9pt, fill=topolook!7},
  learnspace/.style={box, dsh, draw=topoacc, line width=0.9pt, fill=topoacc!5},
  won/.style={box, draw=topogood, line width=1.0pt, fill=topogood!8},
  refuted/.style={box, draw=topobad, line width=1.0pt, fill=topobad!6},
  op/.style={box, dash pattern=on 1.5pt off 1.2pt, draw=topomut,
             line width=0.7pt, fill=white, rounded corners=3.2pt},
  % --- edges --------------------------------------------------------------
  flow/.style={-{Stealth[length=4.2pt,width=3.4pt]}, draw=topoink,
               line width=0.7pt},
  proj/.style={flow, draw=topoacc},
  data/.style={flow, draw=topomut, dsh},
  tie/.style={flow, draw=topogood, line width=0.8pt},
  grad/.style={flow, draw=topobad, line width=0.75pt,
               dash pattern=on 1.6pt off 1.3pt},
  variant/.style={flow, draw=topobad, line width=0.75pt,
                  dash pattern=on 1.6pt off 1.3pt},
  % --- text ---------------------------------------------------------------
  tag/.style={font=\scriptsize, text=topoink, inner sep=1.2pt},
  edg/.style={font=\scriptsize, text=topoink, inner sep=1.4pt},
  note/.style={font=\scriptsize, text=topoink, align=flush left,
               text width=134mm, inner sep=1.4pt, anchor=north west},
  grp/.style={draw=topomut!80, line width=0.5pt, rounded corners=2.4pt,
              dash pattern=on 1pt off 1pt, inner sep=3.4pt},
  swatch/.style={minimum width=5.2mm, minimum height=3.2mm, inner sep=0pt},
}
\fi

%
%  Idempotent: guarded, so a second \input is a no-op.  Requires
%      \usepackage{tikz}
%      \usetikzlibrary{positioning,arrows.meta,fit,backgrounds,calc}
% ===========================================================================
\ifdefined\nexttracktopostyle\else
\newcommand{\nexttracktopostyle}{}

\providecommand{\dsid}[1]{\texttt{\small #1}}
\providecommand{\Rp}{\ensuremath{R'}}
\providecommand{\Cp}{\ensuremath{C'}}

% --- palette: the same hex values as docs/figures/make_figures.py ----------
\definecolor{topoacc}{HTML}{2563EB}   % ACCENT  --- default / on-record path
\definecolor{topogood}{HTML}{16A34A}  % GOOD    --- favoured variant
\definecolor{topobad}{HTML}{DC2626}   % BAD     --- refuted variant
\definecolor{topomut}{HTML}{9CA3AF}   % MUTED   --- frozen, parameter-free
\definecolor{topolook}{HTML}{9333EA}  % seq-markov purple --- table lookup
\definecolor{topoink}{HTML}{374151}   % annotation ink

\tikzset{
  topo/.style={
    font=\footnotesize,
    every node/.append style={align=center},
  },
  box/.style={rounded corners=1.4pt, inner xsep=3.2pt, inner ysep=2.8pt,
              minimum height=6.4mm, text=topoink},
  dsh/.style={dash pattern=on 2pt off 1.4pt},
  % --- component styles ---------------------------------------------------
  trained/.style={box, draw=topoacc, line width=0.9pt, fill=topoacc!7},
  frozen/.style={box, dsh, draw=topomut, line width=0.8pt, fill=topomut!14},
  lookup/.style={box, dsh, draw=topolook, line width=0.9pt, fill=topolook!7},
  learnspace/.style={box, dsh, draw=topoacc, line width=0.9pt, fill=topoacc!5},
  won/.style={box, draw=topogood, line width=1.0pt, fill=topogood!8},
  refuted/.style={box, draw=topobad, line width=1.0pt, fill=topobad!6},
  op/.style={box, dash pattern=on 1.5pt off 1.2pt, draw=topomut,
             line width=0.7pt, fill=white, rounded corners=3.2pt},
  % --- edges --------------------------------------------------------------
  flow/.style={-{Stealth[length=4.2pt,width=3.4pt]}, draw=topoink,
               line width=0.7pt},
  proj/.style={flow, draw=topoacc},
  data/.style={flow, draw=topomut, dsh},
  tie/.style={flow, draw=topogood, line width=0.8pt},
  grad/.style={flow, draw=topobad, line width=0.75pt,
               dash pattern=on 1.6pt off 1.3pt},
  variant/.style={flow, draw=topobad, line width=0.75pt,
                  dash pattern=on 1.6pt off 1.3pt},
  % --- text ---------------------------------------------------------------
  tag/.style={font=\scriptsize, text=topoink, inner sep=1.2pt},
  edg/.style={font=\scriptsize, text=topoink, inner sep=1.4pt},
  note/.style={font=\scriptsize, text=topoink, align=flush left,
               text width=134mm, inner sep=1.4pt, anchor=north west},
  grp/.style={draw=topomut!80, line width=0.5pt, rounded corners=2.4pt,
              dash pattern=on 1pt off 1pt, inner sep=3.4pt},
  swatch/.style={minimum width=5.2mm, minimum height=3.2mm, inner sep=0pt},
}
\fi

%
%  Idempotent: guarded, so a second \input is a no-op.  Requires
%      \usepackage{tikz}
%      \usetikzlibrary{positioning,arrows.meta,fit,backgrounds,calc}
% ===========================================================================
\ifdefined\nexttracktopostyle\else
\newcommand{\nexttracktopostyle}{}

\providecommand{\dsid}[1]{\texttt{\small #1}}
\providecommand{\Rp}{\ensuremath{R'}}
\providecommand{\Cp}{\ensuremath{C'}}

% --- palette: the same hex values as docs/figures/make_figures.py ----------
\definecolor{topoacc}{HTML}{2563EB}   % ACCENT  --- default / on-record path
\definecolor{topogood}{HTML}{16A34A}  % GOOD    --- favoured variant
\definecolor{topobad}{HTML}{DC2626}   % BAD     --- refuted variant
\definecolor{topomut}{HTML}{9CA3AF}   % MUTED   --- frozen, parameter-free
\definecolor{topolook}{HTML}{9333EA}  % seq-markov purple --- table lookup
\definecolor{topoink}{HTML}{374151}   % annotation ink

\tikzset{
  topo/.style={
    font=\footnotesize,
    every node/.append style={align=center},
  },
  box/.style={rounded corners=1.4pt, inner xsep=3.2pt, inner ysep=2.8pt,
              minimum height=6.4mm, text=topoink},
  dsh/.style={dash pattern=on 2pt off 1.4pt},
  % --- component styles ---------------------------------------------------
  trained/.style={box, draw=topoacc, line width=0.9pt, fill=topoacc!7},
  frozen/.style={box, dsh, draw=topomut, line width=0.8pt, fill=topomut!14},
  lookup/.style={box, dsh, draw=topolook, line width=0.9pt, fill=topolook!7},
  learnspace/.style={box, dsh, draw=topoacc, line width=0.9pt, fill=topoacc!5},
  won/.style={box, draw=topogood, line width=1.0pt, fill=topogood!8},
  refuted/.style={box, draw=topobad, line width=1.0pt, fill=topobad!6},
  op/.style={box, dash pattern=on 1.5pt off 1.2pt, draw=topomut,
             line width=0.7pt, fill=white, rounded corners=3.2pt},
  % --- edges --------------------------------------------------------------
  flow/.style={-{Stealth[length=4.2pt,width=3.4pt]}, draw=topoink,
               line width=0.7pt},
  proj/.style={flow, draw=topoacc},
  data/.style={flow, draw=topomut, dsh},
  tie/.style={flow, draw=topogood, line width=0.8pt},
  grad/.style={flow, draw=topobad, line width=0.75pt,
               dash pattern=on 1.6pt off 1.3pt},
  variant/.style={flow, draw=topobad, line width=0.75pt,
                  dash pattern=on 1.6pt off 1.3pt},
  % --- text ---------------------------------------------------------------
  tag/.style={font=\scriptsize, text=topoink, inner sep=1.2pt},
  edg/.style={font=\scriptsize, text=topoink, inner sep=1.4pt},
  note/.style={font=\scriptsize, text=topoink, align=flush left,
               text width=134mm, inner sep=1.4pt, anchor=north west},
  grp/.style={draw=topomut!80, line width=0.5pt, rounded corners=2.4pt,
              dash pattern=on 1pt off 1pt, inner sep=3.4pt},
  swatch/.style={minimum width=5.2mm, minimum height=3.2mm, inner sep=0pt},
}
\fi


This section is for a reader who knows mathematics and machine learning but not
recommender-systems conventions. Nothing below is assumed later.
Figure~\ref{fig:protocol} draws the evaluation protocol every number in this
document is produced by; the topology diagrams of Section~\ref{sec:zoo} all
presuppose it.

\begin{figure}[ht]
  \centering
  \begin{tikzpicture}[topo]
    % ---- row 1: the session is cut into a prefix and one held-out target ----
    \node[frozen, minimum width=38mm, minimum height=8.5mm] (sess)
      {one \emph{test} session, in play order\\$[\,i_1,\ i_2,\ \dots,\ i_{L-1}\ \mid\ i_L\,]$};
    \node[frozen, right=8mm of sess, minimum width=22mm, minimum height=8.5mm] (pfx)
      {\textbf{prefix} $i_1\dots i_{L-1}$};
    \node[trained, right=8mm of pfx, minimum width=28mm, minimum height=8.5mm] (mod)
      {model $\to\ \hat{y}\in\mathbf{R}^{192}$};
    \draw[flow] (sess.east) -- (pfx.west);
    \draw[flow] (pfx.east) -- (mod.west);
    % ---- the held-out target: the other half of the same cut ---------------
    \node[won, anchor=north west, minimum width=28mm, minimum height=6.4mm]
      (tgt) at ([yshift=-4.4mm]sess.south west)
      {\textbf{target} $i_L$ --- the next \emph{distinct} track};
    \draw[data] ([xshift=-5mm]sess.south) -- ([xshift=-5mm]tgt.north);
    % ---- row 2: rank the catalog, cut to ten, score against the target -----
    \node[op, anchor=north west, minimum width=38mm, minimum height=8.5mm]
      (rank) at ([yshift=-5.5mm]tgt.south west)
      {rank all \num{19402} items minus\\the prefix, by $\cos(\hat{y},z_i)$};
    \node[op, right=8mm of rank, minimum width=22mm, minimum height=8.5mm] (top)
      {top-\num{10} list};
    \node[won, right=8mm of top, minimum width=28mm, minimum height=6mm] (hit)
      {$i_L$ \emph{in} it $\Rightarrow$ \textbf{hit}};
    \node[refuted, below=2.6mm of hit, minimum width=28mm, minimum height=6mm] (miss)
      {$i_L$ \emph{not} in it $\Rightarrow$ \textbf{miss}};
    \draw[flow] (mod.south) -- ++(0,-5.4mm) -| ([xshift=-4mm]rank.north east);
    \draw[flow] (rank.east) -- (top.west);
    \draw[tie] (top.east) -- (hit.west);
    \draw[variant] (top.east) -- (miss.west);
    \node[tag, below=1.0mm of rank]
      {the split cut is chronological --- no test session is ever trained on};
    \node[note] at ([yshift=-3.5mm]current bounding box.south west)
      {\textbf{How every number in this document is produced.} Repeat the above
       for each of the \num{1431} test sessions. \emph{Recall@10} is the fraction
       that are hits (reported as ``value (raw hits)'', e.g.\ \num{0.21174}
       (\num{303})); \emph{MRR} averages $1/\mathrm{rank}(i_L)$ over the full
       ranking; the \emph{graded} diagnostics ask instead whether some top-ten
       candidate shares $i_L$'s artist or genre, and music@10 how close the
       nearest of the ten is to $i_L$ in the content metric $M$. Each session is
       additionally labelled \textbf{warm} or \textbf{cold} by whether $i_L$'s
       item index occurs anywhere in the \num{5723} training sessions ---
       \num{643} warm, \num{788} cold --- which is the partition
       Section~\ref{sec:fusion} turns on. During \emph{training} the same picture
       is applied at \emph{every} prefix position of every training session
       rather than only the last (``teacher forcing'', below), which is where the
       \num{73632} supervision events come from.};
  \end{tikzpicture}
  \caption{The evaluation protocol, once. One test session is cut into an
    ordered prefix and a single held-out target; the model turns the prefix into
    one \num{192}-dimensional point; the whole catalog minus the prefix is ranked
    by cosine to that point; the top ten either contains the target or does not.
    Everything else in this document is an aggregate over \num{1431} such trials,
    or a property of the ten returned items rather than of the target.
    Solid borders mark trained components, dashed borders frozen or
    parameter-free ones --- the convention reused by
    Figures~\ref{fig:topo-gru}--\ref{fig:topo-stacker}.}
  \label{fig:protocol}
\end{figure}

\begin{description}
  \item[Cosine retrieval in a frozen item space.] Every model here scores
  candidates the same way: it emits a single predicted vector $\hat{y}\in
  \mathbf{R}^{192}$ (or a full score column, for the lookup legs), and candidate
  $i$ is ranked by $\cos(\hat{y}, z_i)$ against the \emph{frozen} item latent
  $z_i$. The item table is not a parameter --- it is read twice, once to embed
  the prefix and once as the retrieval target. This is why the geometry of that
  table is load-bearing, and why the single largest win on record is a learned
  \emph{warp} of it (Section~\ref{sec:projection}).
  \item[Teacher forcing (per-step supervision).] For a session $[a,b,c,d]$ the
  supervised steps are $b\mid[a]$, $c\mid[a,b]$, $d\mid[a,b,c]$: at every
  timestep the model predicts the \emph{next} item from the prefix so far, so one
  session of length $L$ yields $L-1$ training examples. On this split that is
  \num{73632} supervision events rather than \num{5723}. The distinction is not
  bookkeeping --- two of the worst results on record are caused by an
  architecture that forced a switch to a one-target-per-session objective
  (Sections~\ref{sec:fail-combiners} and~\ref{sec:fail-topology}).
  \item[In-batch InfoNCE.] The contrastive objective every trained model here
  uses. For a minibatch of $B$ (prefix, target) pairs, embed the predictions
  $a_i$ and the targets $p_j$, $L2$-normalize both, form the $B\times B$ logit
  matrix $(a\cdot p^{\top})/\tau$ and apply cross-entropy with the
  \emph{diagonal} as the labels. The $B-1$ off-diagonal targets are the
  ``in-batch negatives'' --- free negatives obtained by reusing every other row's
  positive. The temperature $\tau$ (\num{0.07} throughout unless stated) is the
  softmax sharpness: smaller $\tau$ demands a larger margin over the negatives.
  InfoNCE matters because ranking rewards a representation that \emph{separates}
  the true next item from distractors, which a contrastive loss optimizes
  directly and a pointwise cosine regression does not.
  \item[The z-blend (score-level fusion).] Each leg produces a full score column
  over the admissible candidates (all catalog items minus the prefix). Each
  column is z-scored \emph{over that candidate set} and the columns are averaged
  with hard-coded weights: equal thirds for the three-leg recipe,
  $\alpha z(\text{GRU}) + (1-\alpha)z(\text{Markov})$ with $\alpha=\num{0.5}$
  for the two-leg one. The z-scoring is the load-bearing detail: it puts legs
  whose raw scores live on incomparable scales onto a common per-session scale
  \emph{while preserving relative magnitude}, so a leg that is confidently right
  about one candidate can outvote two diffusely uncertain legs.
  \item[Two spaces, and which one each quantity lives in.] This is the single
  most confusing thing in the setup and it is stated once, here. \textbf{Retrieval}
  happens in the frozen PCA-192 item space: a candidate's rank is
  $\cos(\hat{y},z_i)$ with $z_i\in\mathbf{R}^{192}$. \textbf{Every list-quality
  facet}, by contrast, is a cosine in a \emph{different}, larger space: the
  \num{517}-dimensional intrinsic \emph{content metric} $M$ (the
  \texttt{balanced} dial of \dsid{spotify\_tracks\_content\_metric}). music@10,
  mood\_coh@10, ild@10 \emph{and} MMR's diversity penalty are all cosines in
  $M$; nothing in that group is computed in PCA-192. That split is exactly why
  the facets in Section~\ref{sec:mmr} move together --- one greedy penalty in
  $M$ moves three quantities that are all defined in $M$ --- and why an
  eval-time re-rank can raise them while barely touching a ranking produced in
  PCA-192.
  \item[MMR (maximal marginal relevance) re-ranking.] A purely \emph{eval-time}
  policy: take the top \texttt{mmr\_pool} candidates in base-score order,
  min-max normalize their scores to $\mathrm{rel}\in[0,1]$ over that pool, then
  build the top-$k$ greedily by
  \[
    \arg\max_i\ \lambda\,\mathrm{rel}(i) - (1-\lambda)\max_{j\in\text{selected}}
      \cos(M_i, M_j),
  \]
  with $M$ the \num{517}-dimensional content metric just defined. In words: keep
  taking the most relevant remaining track, minus a penalty equal to its
  similarity to whichever already-chosen track it most resembles. The penalty is
  \emph{continuous} --- there is no similarity threshold and no cutoff radius;
  a highly relevant track can survive close to an already-picked one.
  $\lambda=1$ recovers the base order exactly, so it is an exact identity gate;
  $\lambda\to0$ ignores relevance and merely spreads the list out sonically.
  Crucially the \emph{trained weights are bit-identical} across $\lambda$ arms of
  one model, which is what makes those comparisons far sharper than
  architecture comparisons (Section~\ref{sec:power}).
  \item[Participation ratio.] Used once, in
  Figure~\ref{fig:mc-candidate-eigen}. Take the top candidates' latents,
  weight each by its share of the score mass, centre them, and let $s_1\ge s_2\ge
  \dots$ be the resulting covariance eigenvalues normalized to sum to one. Then
  $\mathrm{pr}=1/\sum_i s_i^2$ is the \emph{effective number of directions} the
  belief spreads over: $\mathrm{pr}=1$ when all the variance is on one axis, and
  $\mathrm{pr}=m$ when $m$ directions carry equal shares. It is not entropy over
  items --- a tight cluster of fifty near-identical tracks has high item entropy
  but $\mathrm{pr}\approx1$.
  \item[CKA (centered kernel alignment).] A scalar in $[0,1]$ measuring how
  similar two layers' \emph{representation geometries} are, computed from their
  centred Gram matrices; $1$ means the two carry the same pairwise-similarity
  structure up to rotation and scaling, $0$ means unrelated. A middling value
  such as \num{0.458} therefore says the two representations are \emph{partly}
  distinct, not that a fraction of them is new information.
  \item[Paired bootstrap.] The decisive statistic throughout is a \emph{paired}
  per-session $\Delta$: resample the \num{1431} test sessions with replacement
  (\num{2000} resamples, rng 1337), recompute both arms' metrics on the
  \emph{same} resampled index set, and report the 95\% interval of the
  difference. A CI entirely above zero is a win, entirely below is a loss, and a
  CI straddling zero is a \emph{tie}. Marginal (unpaired) CIs are also reported
  for the champion rows but are much wider and much weaker.
  \item[The $\pm\num{0.015}$ band --- a policy constant, not a measurement.] One
  number does two jobs in this record and it is worth separating them from the
  measured half-width. As a \emph{practical tie band} it declares any Recall@10
  difference under \num{0.015} not worth acting on; as an \emph{anti-degenerate
  recall floor} it says a holisticness gain does not count unless Recall@10 stays
  within \num{0.015} of the reference. In both roles it is a chosen bright line,
  deliberately set \emph{above} the measured cross-model paired half-width of
  ${\sim}\num{0.0124}$ (Section~\ref{sec:power}) so that it is conservative
  rather than borderline. It is not pre-registered as a noise estimate and it
  should never be quoted as one --- the measured per-comparison half-widths in
  this document range from \num{0.000382} to \num{0.018169} depending on the
  pairing structure.
\end{description}

{\rightskip=0pt plus 2em \emph{Source: \texttt{PROJECT-FACTS.md}.}\par}

\clearpage
\section{Two boards, three definitions: Recall@10 is the currency, holisticness@10 is the declared primary}
\label{sec:metrics}

\emph{Rationale.} Two things are true at once and the document would be
dishonest if it collapsed them. The experimental record's \emph{comparability
currency} is Recall@10: the whole leaderboard is denominated in it, every paired
decision rule is stated on it, and it is stable across every code change of the
past month. The live server's \emph{declared} primary metric, however, is
holisticness@10: \texttt{GET /api/best-models} returns
\texttt{primary\_metric} $=$ \texttt{holisticness@10} and
\texttt{domain.toml [metrics].primary} has read the same since 2026-07-23. The
two orderings are close to inverted, so which board you read decides who is
champion.

\begin{description}
  \item[Recall@10 ($\uparrow$).] The fraction of test sessions whose held-out
  next track appears in the top ten retrieved. Reported throughout as
  ``\num{0.21174} (\num{303})'' --- the value and the raw hit count out of
  \num{1431}.
  \item[MRR ($\uparrow$).] Mean reciprocal rank of the true next track --- the
  average of $1/\mathrm{rank}$, so it rewards getting the truth to position 1
  rather than position 9.
  \item[artist@10 / genre@10 ($\uparrow$).] \emph{Graded} relevance --- credit
  in degrees rather than the binary hit/miss of Recall@10. Both ask whether
  \emph{some} top-ten candidate shares the true next track's artist,
  respectively its primary genre. Exact recall is modest here by construction,
  and landing the right artist ${\sim}\SI{33}{\percent}$ of the time and the
  right genre ${\sim}\SI{40}{\percent}$ is what the models actually deliver.
  \textbf{These two columns must not be read as absolute quality.} Their chance
  floors are high and \emph{have not been measured} on this corpus --- there are
  far fewer distinct genres than items, so a random top ten hits the right genre
  vastly more often than it hits the right track. For contrast, chance Recall@10
  is $10/\num{19402}\approx\num{0.00052}$, and the popularity lookup ---
  which is not random but is the crudest real baseline --- lands \num{1} hit in
  \num{1431} (Recall@10 \num{0.00070}) at genre@10 well above zero. Use
  artist@10 / genre@10 only to compare models against each other.
  \item[music@10 ($\uparrow$).] The \emph{grounding} metric: the mean, over
  sessions, of the best cosine \emph{in the \num{517}-dimensional content metric
  $M$} (Section~\ref{sec:machinery}) between the \emph{true} next track and the
  ten retrieved candidates. It equals
  \num{1.0} whenever the exact truth is in the list, so it is a smooth superset
  of Recall@10 that also credits landing something that merely \emph{sounds
  like} the real continuation.
  \item[mood\_coh@10 ($\uparrow$).] Mean cosine, again in $M$, of the top ten to
  the \emph{prefix}'s mood centroid --- the mean $M$-vector of the tracks already
  played. It asks whether the list belongs to the session it continues, and it
  can go \emph{negative} for a list that is unrelated to the prefix.
  \item[ild@10 ($\uparrow$).] Intra-list diversity: mean pairwise $1-\cos$
  \emph{within} the top ten, again in $M$. Note that this is exactly the
  quantity MMR's penalty term (also a cosine in $M$) pushes down, which is why
  Section~\ref{sec:mmr}'s lever raises it by construction.
  \item[artist\_adj@10 ($\downarrow$).] The ``eagerness'' or seed-echo rate: the
  fraction of the top ten sharing the artist of the \emph{seed} track --- the
  \emph{last item of the prefix}, the track notionally still playing. It measures
  how much a ranker simply parrots what is on. \textbf{Lower is better.}
  \item[artist\_conc@10 ($\downarrow$).] Normalized Herfindahl concentration
  over the top ten's artists: with $p_a$ the share of slots held by artist $a$
  and $n=10$, $(\sum_a p_a^2 - 1/n)/(1-1/n)$, which is \num{0} when every slot
  is a different artist and \num{1} when one artist holds the whole list. It
  exists because artist\_adj@10 is blind to a list of ten tracks by one
  \emph{other} artist --- a blind spot that is widest precisely on the
  pathological models. The out-of-scope \texttt{seq-mood} predictor (promoted
  under the definition name \texttt{mood-session}; called \texttt{seq-mood}
  throughout this document) reads artist\_adj \num{0.003564}, apparently
  ``perfectly non-eager'', at artist\_conc \num{0.329063} --- a ${\sim}90\times$
  gap. Note that the \emph{popularity} lookup also reads artist\_adj
  ${\approx}\num{0.003}$ (\num{0.002655}), for the opposite reason: it returns
  the same ten globally popular tracks for every session and so echoes no seed at
  all. Two different pathologies, one apparently virtuous number.
  \textbf{Lower is better.}
\end{description}

\paragraph{The crown, and its three generations.} holisticness@10 multiplies
four aggregate facets. The \emph{live} (third) definition is
\[
  \begin{aligned}
    H = {}& \mathrm{clamp}(\text{mood\_coh},0)\ \times\ \text{ild}
      \ \times\ \bigl(1-\mathrm{clamp}(\max(\text{artist\_adj},\text{artist\_conc}),0,1)\bigr)\\
      &{}\times\ \mathrm{clamp}\!\left(\frac{\text{music@10}-\num{0.34}}{\num{0.66}},0,1\right).
  \end{aligned}
\]
Three of those four factors are properties of the \emph{returned list} rather
than of its correctness, so a ranker that has given up on identifying the actual
next track can maximize them all at once. That is not a fear, it is a
measurement: under the original three-factor form
$\mathrm{clamp}(\text{mood\_coh},0)\times\text{ild}\times(1-\text{artist\_adj})$
the out-of-scope \texttt{seq-mood} model held best-models rank 1 at $H$
\num{0.2332} on Recall@10 \num{0.0238} (\num{34} of \num{1431}), and a
degenerate two-tower arm that reads only the slow session centroid took rank 2
at Recall@10 \num{0.06639} (\num{95} of \num{1431}). Two fixes followed, both
load-bearing: the \emph{grounding} factor built on music@10 (rescaled, because
music@10's raw spread across models is far narrower than the three-factor
product's, so multiplying in the raw value does not stop the ignore-the-user
mode), and an anti-degenerate \textbf{recall floor} applied as a decision rule
--- an $H$ gain does not count unless Recall@10 stays within $\pm\num{0.015}$ of
the reference. Table~\ref{tab:hgen} shows what those revisions did to identical
predictions.

\begin{table}[ht]
  \centering
  \caption{The same predictions under all three holisticness@10 generations ---
    gen 1 (three-factor, to 2026-07-24), gen 2 (grounded four-factor with
    music floor \num{0.185}, 2026-07-25), gen 3 (concentration-aware, eagerness
    becomes $\max(\text{artist\_adj},\text{artist\_conc})$ and the floor rises to
    \num{0.34}; live since 2026-07-26). Values change by up to $6.4\times$ and
    \emph{so does the ordering}: a stored \texttt{holisticness\_at\_k} is
    comparable only \emph{within} one generation. This is survivable because
    every input facet is stored per row in each run's \texttt{predictions.json},
    so any generation is exactly recomputable for any past run --- the
    generation labels used in this document were re-derived that way and matched
    by exact residual (\numrange{0}{3e-7}). Recall@10 is given for scale.
    \textsuperscript{*}~run \protect\runid{run-20260726-161801-0df6f} and its
    bit-identical twin \protect\runid{run-20260725-143502-12fe6} are the same
    configuration with the same facets, and they \emph{store} $H$ \num{0.022841}
    and \num{0.146354} respectively --- the cleanest single demonstration that a
    stored value carries its generation with it.
    Every cell here is a value that exists: gen-1 and gen-2 figures are as
    recorded in the campaign reports, gen-3 figures as returned by the live API.
    The only ``---'' in the table is the one genuinely uncomputable cell, the
    degenerate arm's gen-3 value: that run is remote-only, never emitted
    \texttt{artist\_conc\_at\_k}, and so cannot be rescored under a formula that
    needs it. Note also \texttt{seq-mood}'s gen-3 value: it is \num{0.0}
    \emph{exactly}, because its music@10 of \num{0.33837} falls below the
    grounding floor \num{0.34} and the fourth factor clamps to zero. That single
    number is the grounding factor working as designed.
    \emph{Sources: \texttt{2026-\allowbreak 07-\allowbreak 25-\allowbreak nexttrack-\allowbreak dual-\allowbreak tower-\allowbreak fusion-\allowbreak scan.md},
    \texttt{2026-\allowbreak 07-\allowbreak 26-\allowbreak nexttrack-\allowbreak pre-\allowbreak encoder-\allowbreak scan.md},
    and \texttt{2026-\allowbreak 07-\allowbreak 26-\allowbreak nexttrack-\allowbreak mmr-\allowbreak lambda-\allowbreak frontier.md}.}}
  \label{tab:hgen}
  {\small\setlength{\tabcolsep}{6pt}\begin{tabular}{lcccc}
    \toprule
    Run & R@10 (raw) & $H$ gen 1 & $H$ gen 2 & $H$ gen 3 (live) \\
    \midrule
    GRU h256 baseline & 0.12299 (176) & 0.141723 & 0.0470862 & 0.024538 \\
    GRU h425 (width control) & 0.12089 (173) & 0.13850 & 0.046677 & 0.024872 \\
    MLP256$\to$GRU256 (S1) & 0.10901 (156) & 0.14499 & 0.047132 & 0.024102 \\
    Dual latent/latent, $f{=}1$\textsuperscript{*} & 0.10273 (147) & 0.146354 & 0.046612 & 0.022841 \\
    Dual cummean/cummean (degenerate) & 0.06639 (95) & 0.19715 & 0.050436 & --- \\
    \texttt{seq-mood} (out of scope) & 0.0238 (34) & 0.23319 & 0.043881 & 0.0 \\
    \bottomrule
  \end{tabular}}
\end{table}

\emph{Results.} On the live board the recall champion ranks \emph{eighth of the
nine} current best-models entries ($H$ \num{0.008603}, above only the two-leg
R+M blend at \num{0.008259}), while the leader is an eval-time re-ranking policy
applied to the plain single GRU ($H$ \num{0.034171} at \num{164} of \num{1431}).
Figure~\ref{fig:mc-two-boards} plots the near-inversion;
Figure~\ref{fig:nexttrack-holisticness-antiwin} shows the failure mode the
recall floor was written for.

\emph{Results --- the two statements that look contradictory, reconciled.} The
record also says that the crown ordering carries \emph{no information}, and both
things are true because they are statements about different parts of the axis.
Take the ten separately-trained architecture arms that carry a live $H$: they
span \numrange{0.02127}{0.02487}, a spread of ${\sim}\num{0.0036}$ against a
measured $\mathrm{hw}(\Delta H)$ of ${\sim}\num{0.0018}$ --- so within that
\emph{healthy band} the arms are statistically indistinguishable and their rank
order is noise, exactly as \texttt{PROJECT-FACTS.md} states. The inversion is not
inside that band. The recall champion sits at \num{0.008603}, roughly seven
half-widths \emph{below} the band's floor, and the MMR arms sit up at
\numrange{0.026761}{0.034171}, well above its ceiling and resolved by a far
sharper instrument (Section~\ref{sec:power}). So: the champion-versus-band gap
and the band-versus-MMR gap are real and large; the ordering \emph{within} the
band is not a ranking of anything.

\emph{Interpretation.} Picking one board fails in both directions: ranking by
$H$ alone crowns models with near-zero recall, while ranking by Recall@10 alone
hides that the recall champion is the most artist-eager model measured
(artist\_adj \num{0.6097}, artist\_conc \num{0.6917} --- inherited from its
Markov leg, whose standalone artist\_adj is \num{0.543}, not from its neural
one, since the single GRU and LSTM sit at \num{0.347} and \num{0.344}). That the
display facets are not merely ``gameable by being good'' is settled by the
popularity floor: it \emph{maximizes} ild (\num{0.898}) and \emph{minimizes}
artist\_adj (\num{0.002655}, by echoing no seed at all rather than by being
diverse) while landing one hit in \num{1431} and a \emph{negative} mood
coherence ($-\num{0.041}$). This document therefore
denominates its board in Recall@10 and reports $H$ beside it, always
generation-tagged.

{\rightskip=0pt plus 2em \emph{Sources: \texttt{2026-\allowbreak 07-\allowbreak 20-\allowbreak musical-\allowbreak distance-\allowbreak metric.md},
\texttt{2026-\allowbreak 07-\allowbreak 22-\allowbreak session-\allowbreak holisticness-\allowbreak and-\allowbreak antieager-\allowbreak levers.md},
and \texttt{2026-\allowbreak 07-\allowbreak 26-\allowbreak nexttrack-\allowbreak mmr-\allowbreak lambda-\allowbreak frontier.md}.}\par}

\begin{figure}[ht]
  \centering
  \includegraphics[width=0.9\textwidth]{mc_two_boards}
  \caption{The two boards are nearly inverted \emph{across} the healthy band,
    and carry no information \emph{within} it. Recall@10 against
    \emph{generation-3 (live)} holisticness@10 for the \num{22} distinct
    canonical-split configurations whose stored $H$ is gen 3 (a further ten
    canonical-split runs are bit-identical repeats of these, and nine store gen-1
    or gen-2 values and are excluded rather than rescaled). The recall champion
    sits bottom-right at $H$ \num{0.008603}; the crown leader is an eval-time MMR
    re-rank of the plain GRU at \num{0.034171} on \num{164} of \num{1431}. MMR
    arms are drawn hollow because they share bit-identical trained weights with
    their parent --- \textbf{none of them is promoted}, and the point labelled
    $\lambda$0.3 on the champion trace is arm A4 of Table~\ref{tab:mmr}, a TIER-3
    \emph{rejected} arm that gets its high $H$ by giving up $\num{0.102}$ of
    recall; it is plotted to map the frontier, not as a candidate. The shaded
    green strip is the \emph{healthy band} \numrange{0.02127}{0.02487} spanned by
    the ten separately-trained architecture arms, and the drawn measurement bar is
    the point: that spread (${\sim}\num{0.0036}$) is barely twice the measured
    $\mathrm{hw}(\Delta H)\approx\num{0.0018}$, so the ordering \emph{inside} the
    strip is noise while the champion's position far below it, and the MMR arms'
    far above it, are not. The shaded red region below Recall@10
    ${\approx}\num{0.108}$ is the degenerate zone the anti-degenerate recall floor
    excludes; the gen-1-only \texttt{cummean/cummean} arm lives there and is
    annotated rather than plotted, since rescaling it across generations would be
    fabrication.}
  \label{fig:mc-two-boards}
\end{figure}

\begin{figure}[ht]
  \centering
  \includegraphics[width=0.8\textwidth]{nexttrack_holisticness_antiwin}
  \caption{Why the crown needs a floor: unconstrained list-quality is maximized
    by predicting \emph{worse}. \textbf{This figure is generation-1
    (three-factor) holisticness} and its $y$ values are not comparable to
    Figure~\ref{fig:mc-two-boards}'s or Table~\ref{tab:hgen}'s gen-3 column.
    Green circles are the five dual-tower arms whose paired $\Delta H$ CI clears
    zero; every one of them falls below the dashed anti-degenerate recall floor
    at $\num{0.12299}-\num{0.015}=\num{0.10799}$ (shaded), so no crown changed
    hands. The arrowed open square is the \texttt{seq-mood} incumbent at $H$
    \num{0.2332} on \num{34} hits of \num{1431} --- the canonical instance of the
    pathology; the figure labels it by its promoted definition name,
    \texttt{mood-session}. Two labels in the figure are the source campaign's
    internal shorthand: its legend's ``Rule 6'' is that campaign's numbering for
    the crown decision rule --- an $H$ gain counts only if Recall@10 holds within
    $\pm\num{0.015}$ --- and the arm codes are C0 the single-GRU baseline, C1 the
    $h425$ width control, and D-\emph{XY} the two-tower arm reading view $X$ in
    tower A and view $Y$ in tower B (L $=$ \texttt{latent}, D $=$ \texttt{delta},
    C $=$ \texttt{cummean}). \emph{Source:
    \texttt{2026-\allowbreak 07-\allowbreak 25-\allowbreak nexttrack-\allowbreak dual-\allowbreak tower-\allowbreak fusion-\allowbreak scan.md}.}}
  \label{fig:nexttrack-holisticness-antiwin}
\end{figure}

\clearpage
\section{The model zoo: six topologies over one frozen geometry}
\label{sec:zoo}

\emph{Rationale.} The six trained families differ in exactly one respect ---
\emph{how the ordered prefix becomes one predicted point (or one score column)}
--- and are otherwise deliberately identical: same item space, same in-batch
InfoNCE, same batching, same held-out-train-slice early stopping, same
$\tau=\num{0.07}$, same seed 1337. That is what makes the comparisons in
Sections~\ref{sec:fusion}--\ref{sec:failures} clean.
Figures~\ref{fig:topo-gru}--\ref{fig:topo-stacker} draw all six; the visual
convention (solid border $=$ trained, dashed $=$ frozen or parameter-free, and
colour carries the verdict) is fixed in Figure~\ref{fig:topo-gru}'s caption and
used by the rest. All six presuppose the evaluation protocol of
Figure~\ref{fig:protocol}: each one is a different answer to ``how does the
prefix become $\hat{y}$?'', and everything downstream of $\hat{y}$ is identical.

\emph{A one-line map before the diagrams.} The six differ as follows.
\texttt{seq-nexttrack} (Figure~\ref{fig:topo-gru}) runs a \emph{recurrence} over
the prefix latents and reads $\hat{y}$ off the final hidden state --- the
workhorse, and the single-model bar. \texttt{seq-ann}
(Figure~\ref{fig:topo-ann}) replaces the recurrence with an order-blind
\emph{pooling} (mean of the prefix plus the last item) through an MLP --- the
control that isolates what recurrence buys. \texttt{seq-dualgru}
(Figure~\ref{fig:topo-dualgru}) runs \emph{two} independent recurrences over
different causal views of the same prefix and fuses them --- refuted, twice over,
with the fusion head itself the larger harm. \texttt{seq-embed}
(Figure~\ref{fig:topo-embed}) would make the item table \emph{trainable} instead
of frozen --- the standard formulation in the literature, and the one family here
with no runs at all. \texttt{seq-blend} (Figure~\ref{fig:topo-blend}) does not
predict a point at all: it fuses three full-vocabulary \emph{score columns} with
fixed weights --- the champion. \texttt{seq-stacker}
(Figure~\ref{fig:topo-stacker}) takes the same score columns and \emph{learns}
the combination with a gradient-boosted tree ensemble --- the fourth of four
learned-combiner nulls. Read the pair \texttt{seq-blend} versus
\texttt{seq-stacker} as the document's central contrast: same inputs, fixed
versus learned combination, and a \num{0.13} recall gap in favour of fixed.

% ===========================================================================
%  topologies.tex --- TikZ network-topology diagrams for the next-track
%  MODEL-COMPARISON document.  Not a standalone document: \input it from the
%  body of the main .tex, which must be compiled from docs/:
%
%      % ===========================================================================
%  topologies.tex --- TikZ network-topology diagrams for the next-track
%  MODEL-COMPARISON document.  Not a standalone document: \input it from the
%  body of the main .tex, which must be compiled from docs/:
%
%      % ===========================================================================
%  topologies.tex --- TikZ network-topology diagrams for the next-track
%  MODEL-COMPARISON document.  Not a standalone document: \input it from the
%  body of the main .tex, which must be compiled from docs/:
%
%      \input{figures/topologies}
%
%  (\graphicspath{{figures/}} governs \includegraphics only, NOT \input, so the
%   figures/ prefix is spelled out above.)
%
%  ------------------------------------------------------------------------
%  REQUIRED PREAMBLE LINES in the main document --- add these two lines after
%  \usepackage{graphicx}, before \begin{document}:
%
%      \usepackage{tikz}
%      \usetikzlibrary{positioning,arrows.meta,fit,backgrounds,calc}
%
%  Nothing else is required: xcolor arrives with tikz, and every colour and
%  every style used below is defined in this file.  \dsid, \Rp and \Cp are
%  \providecommand-guarded here so the file also compiles inside a bare
%  wrapper; when the main document defines them, its definitions win.
%
%  PACKAGES DELIBERATELY *NOT* REQUIRED: amssymb is absent from the house
%  preamble, so nothing below uses \mathbb.  Every picture is laid out with
%  RELATIVE positioning (right=/above=/below= of ...), never with hard
%  coordinates, so no node can collide with its neighbour if the font or the
%  wording changes; the per-figure notes and the legend are anchored to
%  `current bounding box', so they always land underneath the diagram.
%  ------------------------------------------------------------------------
%
%  LABELS DEFINED IN THIS FILE (in file order)
%      fig:topo-gru       seq-nexttrack --- the baseline; ALSO CARRIES THE LEGEND
%      fig:topo-ann       seq-ann       --- pooled feed-forward twin
%      fig:topo-dualgru   seq-dualgru   --- two towers + fusion head + pre-MLP variant
%      fig:topo-embed     seq-embed     --- learned, weight-tied item table (DESIGN ONLY)
%      fig:topo-blend     seq-blend     --- the champion, three legs
%      fig:topo-stacker   seq-stacker   --- learned XGBoost combiner
%
%  VISUAL LANGUAGE (the primary axis is the LINE, not the colour, so it
%  survives greyscale printing; stated for the reader in the caption of
%  fig:topo-gru and re-pointed-at by every other caption)
%      SOLID  border  = TRAINED component (parameters fitted by gradient descent)
%      DASHED border  = FROZEN or PARAMETER-FREE component
%      blue   = the on-record default path
%      purple dashed  = non-parametric table lookup over the dataset
%      dashed blue    = parameter-free, but operating inside a LEARNED geometry
%      green  = variant favoured by the record
%      red    = variant refuted by the record
% ===========================================================================

\input{figures/topo-style}



% ===========================================================================
% 1. seq-nexttrack --- the baseline (and the legend)
% ===========================================================================
\begin{figure}[ht]
  \centering
  \begin{tikzpicture}[topo]
    % ---- the frozen item-latent table, read TWICE ------------------------
    \node[frozen, minimum width=20mm, minimum height=12mm] (tab)
      {frozen item\\latents $Z$\\[1pt]{\scriptsize$19{,}402\times192$}};
    % ---- the prefix, in order (one column per step) ----------------------
    \node[frozen, minimum width=12mm, right=7mm of tab]  (x1) {$x_1$};
    \node[frozen, minimum width=12mm, right=4mm of x1]   (x2) {$x_2$};
    \node[tag, right=2.5mm of x2]                        (xd) {$\cdots$};
    \node[frozen, minimum width=12mm, right=2.5mm of xd] (x3) {$x_t$};
    \begin{scope}[on background layer]
      \node[grp, fit=(x1)(x3)] (pfx) {};
    \end{scope}
    \draw[data] (tab.east) -- (pfx.west);
    \node[tag, below=1.2mm of pfx] (pfxlab)
      {prefix latents --- \emph{order preserved}};
    % ---- the recurrence --------------------------------------------------
    \node[trained, minimum width=12mm, above=7mm of x1] (g1) {GRU};
    \node[trained, minimum width=12mm, above=7mm of x2] (g2) {GRU};
    \node[tag] (gd) at (g1 -| xd) {$\cdots$};
    \node[trained, minimum width=12mm, above=7mm of x3] (g3) {GRU};
    \draw[flow] (x1) -- (g1);
    \draw[flow] (x2) -- (g2);
    \draw[flow] (x3) -- (g3);
    \draw[flow] (g1) -- node[edg, above, inner sep=0.7pt] {$h_t$} (g2);
    \draw[flow] (g2) -- (gd);
    \draw[flow] (gd) -- (g3);
    % ---- head and retrieval ---------------------------------------------
    \node[trained, right=6mm of g3] (head)
      {linear head\\{\scriptsize$256\to192$}};
    \node[op, right=7mm of head] (rk)
      {rank the catalogue by $\cos(\hat z, z_v)$\\{\scriptsize prefix excluded $\to$ top-10}};
    \draw[flow] (g3) -- (head);
    \draw[flow] (head) -- node[edg, above, inner sep=1pt] {$\hat z$} (rk);
    % ---- the SAME frozen table is the retrieval target -------------------
    \coordinate (c1) at ([yshift=-3mm]tab.south |- current bounding box.south);
    \coordinate (c2) at (rk.south |- c1);
    \draw[data] (tab.south) -- (c1);
    \draw[data] (c1) -- node[tag, above]
      {the \emph{same} frozen table is the retrieval target} (c2);
    \draw[data] (c2) -- (rk.south);
    % ---- supervision note ------------------------------------------------
    \node[note] (nt) at ([yshift=-3.5mm]current bounding box.south west)
      {\textbf{Supervision --- per-step teacher forcing.} The head is shared
       across positions and must predict the next item's latent at
       \emph{every} prefix position, so one session of length $L$ yields
       $L-1$ examples: $73{,}632$ (prefix, next-item) pairs on the train
       split. The loss is masked in-batch InfoNCE at $\tau=0.07$ --- each
       step's true next latent must outscore the other steps' targets in the
       same minibatch, which supplies negatives for free.};
    % ---- LEGEND (two rows so it cannot widen the picture) ----------------
    \node[trained, swatch, anchor=north west] (l1)
      at ([yshift=-3mm]current bounding box.south west) {};
    \node[tag, right=1.1mm of l1] (t1) {trained (solid)};
    \node[frozen, swatch, right=3.4mm of t1] (l2) {};
    \node[tag, right=1.1mm of l2] (t2) {frozen / parameter-free (dashed)};
    \node[lookup, swatch, right=3.4mm of t2] (l3) {};
    \node[tag, right=1.1mm of l3] (t3) {table lookup};
    \node[learnspace, swatch, anchor=north west] (l4)
      at ([yshift=-1.6mm]l1.south west) {};
    \node[tag, right=1.1mm of l4] (t4) {parameter-free, learned space};
    \node[won, swatch, right=3.4mm of t4] (l5) {};
    \node[tag, right=1.1mm of l5] (t5) {favoured variant};
    \node[refuted, swatch, right=3.4mm of t5] (l6) {};
    \node[tag, right=1.1mm of l6] (t6) {refuted variant};
  \end{tikzpicture}
  \caption{\texttt{seq-nexttrack}, the single-model baseline every other
    topology in this document is measured against: a frozen item space, a
    recurrence over the \emph{ordered} prefix, a linear head, and cosine
    retrieval back into the \emph{same} frozen space. Nothing about the item
    geometry is learned here --- the $19{,}402\times192$ latent table (offline
    PCA-192 on \dsid{seq-20260715-131139}) is read twice, once to embed the
    prefix and once as the retrieval target, and only the GRU and its linear
    head carry parameters (\num{394944} at $h=256$). \textbf{The visual
    language used throughout this document is fixed here.} A \emph{solid}
    border marks a trained component and a \emph{dashed} border a frozen or
    parameter-free one; that distinction rides on the line style, not the
    colour, so it survives greyscale. Colour then carries the verdict ---
    blue: the on-record default path; purple dashed: a non-parametric table
    lookup over the dataset; dashed blue: parameter-free but operating inside
    a \emph{learned} geometry; green: a variant the record favours; red: a
    variant the record refutes. On the canonical split this configuration
    scores Recall@10 \num{0.12299} (\num{176} of \num{1431}), reproduced
    bit-identically eleven times.}
  \label{fig:topo-gru}
\end{figure}


% ===========================================================================
% 2. seq-ann --- the pooled feed-forward twin
% ===========================================================================
\begin{figure}[ht]
  \centering
  \begin{tikzpicture}[topo]
    \node[frozen, minimum width=20mm, minimum height=12mm] (tab)
      {frozen item\\latents $Z$\\[1pt]{\scriptsize$19{,}402\times192$}};
    \node[frozen, minimum width=12mm, right=7mm of tab]  (x1) {$x_1$};
    \node[frozen, minimum width=12mm, right=4mm of x1]   (x2) {$x_2$};
    \node[tag, right=2.5mm of x2]                        (xd) {$\cdots$};
    \node[frozen, minimum width=12mm, right=2.5mm of xd] (x3) {$x_t$};
    \begin{scope}[on background layer]
      \node[grp, fit=(x1)(x3)] (pfx) {};
    \end{scope}
    \draw[data] (tab.east) -- (pfx.west);
    % ---- the two parameter-free summaries of the same prefix -------------
    \node[op, right=8mm of pfx, yshift=10mm] (mp)
      {causal mean-pool\\$\frac{1}{t}\sum_{i\le t}x_i$\\[1pt]{\scriptsize\emph{permutation-invariant}}};
    \node[op, right=8mm of pfx, yshift=-10mm] (lt)
      {last item\\$x_t$\\[1pt]{\scriptsize the only recency cue}};
    \draw[flow] (pfx.east) -- (mp.west);
    \draw[flow] (pfx.east) -- (lt.west);
    % ---- the MLP ---------------------------------------------------------
    \node[trained, right=5mm of mp, yshift=-10mm] (mlp)
      {MLP\\$384\to256\to256\to192$\\[1pt]{\scriptsize ReLU $+$ dropout}};
    \draw[flow] (mp.east) -- (mlp.north west);
    \draw[flow] (lt.east) -- (mlp.south west);
    \node[edg] at ([xshift=-4.2mm]mlp.west) {$\|$};
    % ---- retrieval, placed below to keep the picture inside the text ------
    \node[op, below=9mm of mlp] (rk)
      {$\hat z\to$ cosine rank\\{\scriptsize prefix excluded $\to$ top-10}};
    \draw[flow] (mlp) -- (rk);
    \coordinate (c1) at ([yshift=-3mm]tab.south |- current bounding box.south);
    \coordinate (c2) at (rk.south |- c1);
    \draw[data] (tab.south) -- (c1);
    \draw[data] (c1) -- node[tag, above]
      {the same frozen table is the retrieval target} (c2);
    \draw[data] (c2) -- (rk.south);
    \node[note] at ([yshift=-3.5mm]current bounding box.south west)
      {\textbf{The same information, with the order discarded.} Both summaries
       are strictly causal, so the objective, the padded batching and the
       early stopping are reused \emph{verbatim} from
       Figure~\ref{fig:topo-gru} --- but a mean weights the track from twenty
       steps ago exactly like the one before last, and cannot represent ``A
       then B'' differently from ``B then A''. Only the special-cased $x_t$ tap
       carries any notion of recency, and the concatenation ($\|$) is this
       model's entire account of sequence.};
  \end{tikzpicture}
  \caption{\texttt{seq-ann}, the GRU's non-recurrent twin: the recurrence of
    Figure~\ref{fig:topo-gru} is replaced by two parameter-free summaries of
    the same prefix --- a causal mean-pool concatenated with the last item's
    latent --- feeding a plain three-layer ReLU MLP. Everything else is
    literally the same code (\texttt{seq\_ann.py} imports
    \texttt{make\_batches} and \texttt{step\_loss} from
    \texttt{seq\_nexttrack}), which is what makes this pair a clean test of
    the recurrence itself rather than of a training recipe. Visual convention
    as in Figure~\ref{fig:topo-gru}: the two pooling operators are dashed
    because they hold no parameters, and only the MLP is trained. The
    defensible comparison is same-space: on \dsid{seq-20260715-030509}
    (PCA-128) the pooled twin scores Recall@10 \num{0.0901} (\num{129} of
    \num{1431}) against the GRU's \num{0.1146} (\num{164}), paired $\Delta$
    $-\num{0.0245}$ $[-0.039, -0.011]$, a CI \emph{entirely below zero}.
    \texttt{seq-ann} has no run at all on the canonical PCA-192 split, so it
    never belongs in the same column as the \num{0.12299} bar.}
  \label{fig:topo-ann}
\end{figure}


% ===========================================================================
% 3. seq-dualgru --- two towers, a fusion head, and the pre-MLP variant
% ===========================================================================
\begin{figure}[ht]
  \centering
  \begin{tikzpicture}[topo]
    \node[frozen, minimum width=19mm, minimum height=11mm] (pf)
      {prefix\\latents\\$x_1\dots x_t$};
    % ---- the three configurable causal views ----------------------------
    \node[op, right=7mm of pf] (v2) {\texttt{delta}\\$x_t-x_{t-1}$};
    \node[op, above=2.6mm of v2] (v1) {\texttt{latent}\\$x_t$};
    \node[op, below=2.6mm of v2] (v3) {\texttt{cummean}\\$\frac{1}{t}\sum_{i\le t}x_i$};
    \begin{scope}[on background layer]
      \node[grp, fit=(v1)(v3)] (vg) {};
    \end{scope}
    \draw[data] (pf.east) -- (vg.west);
    % ---- two independent towers -----------------------------------------
    \node[trained, right=8mm of vg, yshift=9mm] (ta)
      {tower A on \texttt{view\_a}\\GRU $h_a{=}256$\\[1pt]{\scriptsize own weights}};
    \node[trained, right=8mm of vg, yshift=-9mm] (tb)
      {tower B on \texttt{view\_b}\\GRU $h_b{=}256$\\[1pt]{\scriptsize own weights}};
    \draw[flow] (vg.east) -- (ta.west);
    \draw[flow] (vg.east) -- (tb.west);
    % ---- concat + the two fusion depths ---------------------------------
    \node[op, right=7mm of ta, yshift=-9mm] (cc)
      {concat\\$h_a\|h_b$\\[1pt]{\scriptsize$512$}};
    \draw[flow] (ta.east) -- (cc.north west);
    \draw[flow] (tb.east) -- (cc.south west);
    \node[won, right=7mm of cc, yshift=9mm] (f0)
      {\texttt{fusion\_layers}${}=0$\\$512\to192$, linear};
    \node[refuted, right=7mm of cc, yshift=-9mm] (f1)
      {\texttt{fusion\_layers}${}=1$\\$512\to256\to192$\\[1pt]{\scriptsize$+$ ReLU $+$ dropout}};
    \draw[flow] (cc.north east) -- (f0.west);
    \draw[flow] (cc.south east) -- (f1.west);
    % ---- the refuted per-step pre-MLP variant ---------------------------
    \node[refuted, below=12mm of tb] (pb)
      {\emph{variant} --- per-step pre-MLP\\on a tower:
       Linear$(192{\to}256)$\\$+$ ReLU $+$ Dropout};
    \draw[variant] (vg.south) |- (pb.west);
    \draw[variant] (pb.north) -- (tb.south);
    \node[note] at ([yshift=-3.5mm]current bounding box.south west)
      {\textbf{Both heads emit the same $192$-d $\hat z$ and retrieve exactly
       as in Figure~\ref{fig:topo-gru}.} The towers share nothing --- separate
       modules, separate parameters --- and their per-step hidden states are
       concatenated, so the combiner sees the towers' \emph{representations}
       rather than their scores and gradients flow back into both. There is no
       bidirectional mode by construction: a backward pass would let step $t$
       read item $t+1$ and copy its own teacher-forcing target. All three views
       are strictly causal for the same reason, and so is the pre-MLP, which is
       applied \emph{timestep-wise}; it is drawn solid because it is trained,
       red because it is refuted.};
  \end{tikzpicture}
  \caption{\texttt{seq-dualgru}: two recurrent towers with entirely separate
    weights read the prefix in parallel and a head fuses their hidden states
    --- the strongest available form of ``train the combination end-to-end so
    the parts co-adapt''. Each tower's input is a configurable causal view
    (\texttt{view\_a} / \texttt{view\_b}); the default is
    \texttt{latent}/\texttt{latent}, i.e.\ two independent GRUs over the same
    stream. Two of the boxes are coloured by the record rather than by the
    design. Removing the MLP fusion head (\texttt{fusion\_layers}${}=0$,
    green) \emph{beats} keeping it in three independent measurements, most
    recently by $+\num{0.02655}$ $[+0.01398, +0.03913]$ CI$>$0 --- and with
    some $9$--$10\%$ \emph{fewer} parameters, so the linear readout is the
    family default; and the
    per-step pre-MLP (red) loses the decisive single-tower ablation, Recall@10
    \num{0.10901} (\num{156} of \num{1431}) against the baseline's \num{176},
    paired $\Delta$ $-\num{0.01398}$ $[-0.02657, -0.00140]$ CI$<$0. Visual
    convention as in Figure~\ref{fig:topo-gru}. The best dual arm reaches
    \num{0.11880} and merely \emph{ties} the single-GRU bar, as does a
    $0.098\%$ parameter-matched width control --- so capacity is not the
    explanation.}
  \label{fig:topo-dualgru}
\end{figure}


% ===========================================================================
% 4. seq-embed --- a learned, weight-tied item table (DESIGN ONLY)
% ===========================================================================
\begin{figure}[ht]
  \centering
  \begin{tikzpicture}[topo]
    \node[frozen, minimum width=20mm, minimum height=11mm] (ids)
      {session prefix\\item \emph{ids}\\$i_1\dots i_t$};
    \node[trained, right=10mm of ids, minimum height=11mm] (E)
      {learned item\\table $E$\\[1pt]{\scriptsize$19{,}402\times192$}};
    \node[trained, right=6mm of E] (gru)
      {GRU $h{=}256$\\{\scriptsize$+$ dropout}};
    \node[trained, right=6mm of gru] (hd)
      {linear head\\{\scriptsize$256\to192$}};
    \node[op, right=6mm of hd] (rk)
      {$\cos(\hat z, E_v)$ over all\\$19{,}402$ rows $\to$ top-10};
    \draw[flow] (ids) -- node[edg, above, inner sep=1pt] {gather} (E);
    \draw[flow] (E) -- (gru);
    \draw[flow] (gru) -- (hd);
    \draw[flow] (hd) -- node[edg, above, inner sep=1pt] {$\hat z$} (rk);
    % ---- the three modes -------------------------------------------------
    \node[op, text width=64mm, align=flush left, anchor=south west] (md)
      at ([xshift=4mm,yshift=8mm]E.north east)
      {\texttt{embed\_mode} --- what the table $E$ actually is:\\
       \textbf{\texttt{frozen}} $E=Z$, nothing trained (the reproduction gate)\\
       \textbf{\texttt{residual}} $E=Z+\Delta E$: frozen content vector $+$ a
       trained correction, so a cold item still has a position\\
       \textbf{\texttt{free}} $E$ fully trained, optionally randomly initialised};
    \draw[flow] (md.west) -- (E.north east);
    % ---- weight tying: the same table is the retrieval target -----------
    \coordinate (c1) at ([yshift=-8mm]E.south);
    \coordinate (c2) at (rk.south |- c1);
    \draw[tie] (E.south) -- (c1);
    \draw[tie] (c1) -- node[tag, above, text=topogood]
      {weight tying --- the \emph{same} table is\\the input embedding
       \emph{and} the retrieval target} (c2);
    \draw[tie] (c2) -- (rk.south);
    % ---- the gradient loop that makes the geometry trainable ------------
    \coordinate (t1) at ([yshift=5mm]rk.north |- current bounding box.north);
    \coordinate (t2) at (E.north |- t1);
    \draw[grad] (rk.north) -- (t1);
    \draw[grad] (t1) -- node[tag, above, text=topobad]
      {$\nabla$ --- one InfoNCE step moves the query \emph{and} every target:
       the geometry itself is trainable} (t2);
    \draw[grad] (t2) -- (E.north);
    \node[note] at ([yshift=-3.5mm]current bounding box.south west)
      {\textbf{Why the loop is the whole point.} In
       Figures~\ref{fig:topo-gru}--\ref{fig:topo-dualgru} the retrieval space
       is fixed and only the map into it is learned; here the space is a
       parameter, shaped by next-track adjacency rather than inherited from a
       reconstruction objective. That also creates a collapse hazard absent
       everywhere else: with a learned table a plain cosine objective has a
       trivial optimum at ``send every item to one point'', so only a
       contrastive loss is admissible and \texttt{loss=cosine} is rejected
       outright for the trained modes.};
  \end{tikzpicture}
  \caption{\texttt{seq-embed}, the one axis the record has not tested:
    \textbf{a design, not a result}. No run of this predictor exists on any
    dataset, it is absent from the server's registered predictor list, and
    nothing here is a measured score. The item table $E$ is used twice, as the
    GRU's input embedding and as the retrieval target (weight tying, the
    GRU4Rec/SASRec formulation), so a single gradient step moves the query and
    all the candidates together --- the return arrow along the top is what
    makes the geometry trainable, and it is exactly what every other topology
    in this document lacks. The three \texttt{embed\_mode} settings span the
    space from a reproduction gate to a fully free table. What \emph{is}
    measured is the data budget that bounds them: \num{73632} training steps
    to fit \num{15816} item vectors, a median per-item training frequency of
    \num{1}, only \num{3294} items occurring five or more times, and
    \num{788} of \num{1431} test targets never appearing in any training
    session --- so a \texttt{free} table cannot retrieve the cold
    \SI{55}{\percent} at all and is capped at Recall@10 \num{0.4493}, whereas
    \texttt{residual} keeps the content vector and with it the cold reach.
    Visual convention as in Figure~\ref{fig:topo-gru}.}
  \label{fig:topo-embed}
\end{figure}


% ===========================================================================
% 5. seq-blend --- the champion
% ===========================================================================
\begin{figure}[ht]
  \centering
  \begin{tikzpicture}[topo]
    \node[frozen, minimum width=18mm, minimum height=11mm] (pf)
      {session\\prefix\\$i_1\dots i_t$};
    % ---- the learned metric map: shared by two legs, bypassed by one -----
    \node[trained, right=4mm of pf] (W)
      {learned map $W$\\$192\times192$\\[1pt]{\scriptsize no bias, identity init}\\{\scriptsize
       InfoNCE $\tau{=}0.07$}\\{\scriptsize $36{,}864$ params}};
    % ---- the three legs -------------------------------------------------
    \node[learnspace, right=3mm of W] (C)
      {\Cp\ content-$k$NN in $W$-space\\$\max_{u}\cos(Wz_v,Wz_u)$\\[1pt]{\scriptsize
       no parameters of its own}};
    \node[trained, above=7mm of C] (R)
      {\Rp\ projected GRU\\$h{=}256$, $394{,}944$ params};
    \node[lookup, below=7mm of C] (M)
      {$M$ first-order Markov counts\\$P(v\mid u)$, $u=$ last prefix item\\[1pt]{\scriptsize
       Dirichlet trust $n_u/(n_u{+}8)$,}\\{\scriptsize artist/genre back-off; $0$ params}};
    \draw[flow] (pf.east) -- (W.west);
    % ---- the frozen item table W acts ON: the same Z as Figure~\ref{fig:topo-gru} --
    \node[frozen, anchor=south east, minimum width=18mm] (Z)
      at ([yshift=6mm]W.north west)
      {frozen item latents $Z$\\{\scriptsize$19{,}402\times192$}};
    \draw[data] (Z.south east) -- node[edg, left, inner sep=1.2pt] {$z_i$} (W.north west);
    \draw[proj] (W.north east) -- (R.west);
    \draw[proj] (W.east) -- (C.west);
    \draw[data] (pf.south) |- node[tag, above, pos=0.55]
      {item \emph{ids} only ---\\the $M$ leg bypasses $Z$ and $W$} (M.west);
    % ---- the fixed combiner ---------------------------------------------
    \node[op, right=3mm of C] (cb)
      {z-normalise each leg\\over the candidate set,\\then average\\[1pt]{\scriptsize
       equal thirds ---}\\{\scriptsize\emph{no learned weights}}};
    \draw[flow] (R.east) -- node[edg, above, sloped, pos=0.55] {$s_{\Rp}$} (cb.north west);
    \draw[flow] (C.east) -- node[edg, above, inner sep=1pt] {$s_{\Cp}$} (cb.west);
    \draw[flow] (M.east) -- node[edg, above, sloped, pos=0.30] {$s_M$} (cb.south west);
    \node[op, right=2mm of cb] (rk) {rank\\$\to$ top-10};
    \draw[flow] (cb) -- (rk);
    \node[note] at ([yshift=-3.5mm]current bounding box.south west)
      {\textbf{Which legs are learned.} $W$ and the \Rp\ recurrence are
       trained (\num{431808} parameters between them); \Cp\ has no parameters
       of its own but retrieves inside $W$'s learned geometry, which is why it
       is drawn dashed-blue; $M$ is a table of training-set transition counts,
       drawn dashed-purple. The combiner is a fixed rule: each leg's
       full-vocabulary score vector is z-scored over the admissible candidates
       (every item minus those already in the prefix) and the three are
       averaged with hard-coded equal weights. The z-scoring is the
       load-bearing step --- it puts three incomparable score scales onto one
       per-session scale while \emph{preserving} relative magnitude, so a
       confidently-peaked leg can outvote two diffuse ones. Rank-only fusion,
       which discards that magnitude, loses by $-\num{0.029}$ CI$<$0.};
  \end{tikzpicture}
  \caption{\texttt{seq-blend} with \texttt{projection=true}, the champion:
    three full-vocabulary score vectors, z-normalised over the candidate set
    and averaged in fixed equal thirds. $W$ acts on the \emph{frozen item
    latents} $Z$ of Figure~\ref{fig:topo-gru} --- it warps the retrieval geometry,
    it does not read the prefix ids directly --- and it is shared by two of the
    three legs and bypassed by the third; that is the whole mechanism. $W$ is fit by in-batch InfoNCE on consecutive within-train
    pairs, which rebuilds the retrieval metric from ``reconstruct the item
    matrix'' into ``what follows what''; both legs that retrieve by cosine in
    that space improve --- the parameter-free content-$k$NN leg more than
    doubles, \num{0.0685} to \num{0.1516}, and the recurrent leg rises from
    \num{0.1230} to $\approx$\num{0.172} --- while the item-ID lookup is
    untouched at \num{0.106918}. The two projected-leg standalone figures come
    from the off-server driver of the 2026-07-18 campaign and have no server
    run of their own; only the $M$ leg's number is a canonical-split run.
    Together the three legs reach Recall@10 \num{0.21174} (\num{303} of
    \num{1431}), paired $\Delta$ over the unprojected three-leg blend
    $+\num{0.0259}$ $[+0.012, +0.041]$, a CI entirely above zero and robust
    across three seeds. Visual convention as in Figure~\ref{fig:topo-gru}.}
  \label{fig:topo-blend}
\end{figure}


% ===========================================================================
% 6. seq-stacker --- the learned combiner
% ===========================================================================
\begin{figure}[ht]
  \centering
  \begin{tikzpicture}[topo]
    \node[frozen, minimum width=18mm, minimum height=11mm] (pf)
      {session\\prefix\\$i_1\dots i_t$};
    \node[trained, right=6mm of pf] (R)
      {$R$ GRU $h{=}256$\\[1pt]{\scriptsize trained separately}};
    \node[lookup, above=6mm of R] (M) {$M$ Markov counts};
    \node[trained, below=6mm of R] (A)
      {$A$ pooled MLP\\[1pt]{\scriptsize trained separately}};
    \draw[flow] (pf.east) -- (M.west);
    \draw[flow] (pf.east) -- (R.west);
    \draw[flow] (pf.east) -- (A.west);
    % ---- the per-candidate feature row ----------------------------------
    \node[op, right=6mm of R] (fx)
      {per-candidate feature row\\$[\,s_M,\ s_R,\ s_A,\ \log(1{+}\text{plays}),$\\$\text{recency},\
       \text{artist match},\ \text{genre match}\,]$};
    \draw[flow] (M.east) -- (fx.north west);
    \draw[flow] (R.east) -- (fx.west);
    \draw[flow] (A.east) -- (fx.south west);
    \node[refuted, right=5mm of fx] (xgb)
      {XGBoost meta-learner\\$200$ rounds, depth $6$\\[1pt]{\scriptsize\texttt{binary:logistic}}};
    \draw[flow] (fx) -- (xgb);
    \node[op, below=9mm of xgb] (rk)
      {score all $19{,}402$ items\\$\to$ top-10};
    \draw[flow] (xgb) -- (rk);
    % ---- the two failure mechanisms, marked on the diagram --------------
    \node[refuted, text width=56mm, align=flush left, anchor=north west] (tr)
      at ([yshift=-10mm]fx.south west)
      {\textbf{Training rows}: $1$ positive $+$ $50$ popularity-sampled
       negatives per train session, i.e.\ \num{5723} positives, against the
       \num{73632} teacher-forced steps its own $R$ leg trains on
       ($12.9\times$ less signal) --- and a decision surface calibrated on
       $51$ candidates ($0.26\%$ of the catalogue) is then asked to rank all
       \num{19402}.};
    \draw[grad] (tr.north east) -- (xgb.south west);
    \node[note] at ([yshift=-3.5mm]current bounding box.south west)
      {\textbf{Why the shape matters.} A $200$-round depth-$6$ booster is
       piecewise \emph{constant}, so ranking \num{19402} candidates with it
       produces large tied blocks that the harness breaks by item index. The
       fixed z-blend of Figure~\ref{fig:topo-blend} is instead a strictly
       monotone continuous function of the leg scores: it preserves each leg's
       internal ordering exactly and can only trade the legs off against one
       another.};
  \end{tikzpicture}
  \caption{\texttt{seq-stacker}, the learned counterpart of
    Figure~\ref{fig:topo-blend}'s fixed rule: the same separately-trained base
    legs, but their per-candidate scores become columns of a feature row
    alongside four cheap candidate features, and a gradient-boosted tree
    learns the combination. It earns a diagram because its failure is
    informative --- with the \texttt{markov}, \texttt{gru} and \texttt{ann}
    legs it reaches Recall@10 \num{0.085255} (\num{122} of \num{1431}),
    \emph{below its own GRU base leg} at \num{0.12299}, paired $\Delta$
    $-\num{0.0377}$ $[-0.056, -0.020]$ CI$<$0. The red callout names the two
    mechanisms: supervision starvation, since the meta-learner gets one
    labelled event per session where its base leg gets one per prefix
    position, and a train/serve candidate mismatch. The booster is drawn solid
    because it is trained and red because it is refuted; the feature row and
    the ranking step are dashed because they hold no parameters. Visual
    convention as in Figure~\ref{fig:topo-gru}. This is one of four
    independent learned-combiner nulls, the others being reciprocal-rank
    fusion, oracle weights fitted on the test labels, and the jointly-trained
    dual tower of Figure~\ref{fig:topo-dualgru}.}
  \label{fig:topo-stacker}
\end{figure}

%
%  (\graphicspath{{figures/}} governs \includegraphics only, NOT \input, so the
%   figures/ prefix is spelled out above.)
%
%  ------------------------------------------------------------------------
%  REQUIRED PREAMBLE LINES in the main document --- add these two lines after
%  \usepackage{graphicx}, before \begin{document}:
%
%      \usepackage{tikz}
%      \usetikzlibrary{positioning,arrows.meta,fit,backgrounds,calc}
%
%  Nothing else is required: xcolor arrives with tikz, and every colour and
%  every style used below is defined in this file.  \dsid, \Rp and \Cp are
%  \providecommand-guarded here so the file also compiles inside a bare
%  wrapper; when the main document defines them, its definitions win.
%
%  PACKAGES DELIBERATELY *NOT* REQUIRED: amssymb is absent from the house
%  preamble, so nothing below uses \mathbb.  Every picture is laid out with
%  RELATIVE positioning (right=/above=/below= of ...), never with hard
%  coordinates, so no node can collide with its neighbour if the font or the
%  wording changes; the per-figure notes and the legend are anchored to
%  `current bounding box', so they always land underneath the diagram.
%  ------------------------------------------------------------------------
%
%  LABELS DEFINED IN THIS FILE (in file order)
%      fig:topo-gru       seq-nexttrack --- the baseline; ALSO CARRIES THE LEGEND
%      fig:topo-ann       seq-ann       --- pooled feed-forward twin
%      fig:topo-dualgru   seq-dualgru   --- two towers + fusion head + pre-MLP variant
%      fig:topo-embed     seq-embed     --- learned, weight-tied item table (DESIGN ONLY)
%      fig:topo-blend     seq-blend     --- the champion, three legs
%      fig:topo-stacker   seq-stacker   --- learned XGBoost combiner
%
%  VISUAL LANGUAGE (the primary axis is the LINE, not the colour, so it
%  survives greyscale printing; stated for the reader in the caption of
%  fig:topo-gru and re-pointed-at by every other caption)
%      SOLID  border  = TRAINED component (parameters fitted by gradient descent)
%      DASHED border  = FROZEN or PARAMETER-FREE component
%      blue   = the on-record default path
%      purple dashed  = non-parametric table lookup over the dataset
%      dashed blue    = parameter-free, but operating inside a LEARNED geometry
%      green  = variant favoured by the record
%      red    = variant refuted by the record
% ===========================================================================

% ===========================================================================
%  topo-style.tex --- the shared TikZ style block for the next-track topology
%  diagrams AND for the evaluation-protocol diagram in the body of
%  model-comparison.tex.  Extracted from topologies.tex so that the protocol
%  figure, which appears BEFORE \input{figures/topologies}, can use the same
%  visual language.  \input from docs/ (compilation cwd):
%
%      \input{figures/topo-style}
%
%  Idempotent: guarded, so a second \input is a no-op.  Requires
%      \usepackage{tikz}
%      \usetikzlibrary{positioning,arrows.meta,fit,backgrounds,calc}
% ===========================================================================
\ifdefined\nexttracktopostyle\else
\newcommand{\nexttracktopostyle}{}

\providecommand{\dsid}[1]{\texttt{\small #1}}
\providecommand{\Rp}{\ensuremath{R'}}
\providecommand{\Cp}{\ensuremath{C'}}

% --- palette: the same hex values as docs/figures/make_figures.py ----------
\definecolor{topoacc}{HTML}{2563EB}   % ACCENT  --- default / on-record path
\definecolor{topogood}{HTML}{16A34A}  % GOOD    --- favoured variant
\definecolor{topobad}{HTML}{DC2626}   % BAD     --- refuted variant
\definecolor{topomut}{HTML}{9CA3AF}   % MUTED   --- frozen, parameter-free
\definecolor{topolook}{HTML}{9333EA}  % seq-markov purple --- table lookup
\definecolor{topoink}{HTML}{374151}   % annotation ink

\tikzset{
  topo/.style={
    font=\footnotesize,
    every node/.append style={align=center},
  },
  box/.style={rounded corners=1.4pt, inner xsep=3.2pt, inner ysep=2.8pt,
              minimum height=6.4mm, text=topoink},
  dsh/.style={dash pattern=on 2pt off 1.4pt},
  % --- component styles ---------------------------------------------------
  trained/.style={box, draw=topoacc, line width=0.9pt, fill=topoacc!7},
  frozen/.style={box, dsh, draw=topomut, line width=0.8pt, fill=topomut!14},
  lookup/.style={box, dsh, draw=topolook, line width=0.9pt, fill=topolook!7},
  learnspace/.style={box, dsh, draw=topoacc, line width=0.9pt, fill=topoacc!5},
  won/.style={box, draw=topogood, line width=1.0pt, fill=topogood!8},
  refuted/.style={box, draw=topobad, line width=1.0pt, fill=topobad!6},
  op/.style={box, dash pattern=on 1.5pt off 1.2pt, draw=topomut,
             line width=0.7pt, fill=white, rounded corners=3.2pt},
  % --- edges --------------------------------------------------------------
  flow/.style={-{Stealth[length=4.2pt,width=3.4pt]}, draw=topoink,
               line width=0.7pt},
  proj/.style={flow, draw=topoacc},
  data/.style={flow, draw=topomut, dsh},
  tie/.style={flow, draw=topogood, line width=0.8pt},
  grad/.style={flow, draw=topobad, line width=0.75pt,
               dash pattern=on 1.6pt off 1.3pt},
  variant/.style={flow, draw=topobad, line width=0.75pt,
                  dash pattern=on 1.6pt off 1.3pt},
  % --- text ---------------------------------------------------------------
  tag/.style={font=\scriptsize, text=topoink, inner sep=1.2pt},
  edg/.style={font=\scriptsize, text=topoink, inner sep=1.4pt},
  note/.style={font=\scriptsize, text=topoink, align=flush left,
               text width=134mm, inner sep=1.4pt, anchor=north west},
  grp/.style={draw=topomut!80, line width=0.5pt, rounded corners=2.4pt,
              dash pattern=on 1pt off 1pt, inner sep=3.4pt},
  swatch/.style={minimum width=5.2mm, minimum height=3.2mm, inner sep=0pt},
}
\fi




% ===========================================================================
% 1. seq-nexttrack --- the baseline (and the legend)
% ===========================================================================
\begin{figure}[ht]
  \centering
  \begin{tikzpicture}[topo]
    % ---- the frozen item-latent table, read TWICE ------------------------
    \node[frozen, minimum width=20mm, minimum height=12mm] (tab)
      {frozen item\\latents $Z$\\[1pt]{\scriptsize$19{,}402\times192$}};
    % ---- the prefix, in order (one column per step) ----------------------
    \node[frozen, minimum width=12mm, right=7mm of tab]  (x1) {$x_1$};
    \node[frozen, minimum width=12mm, right=4mm of x1]   (x2) {$x_2$};
    \node[tag, right=2.5mm of x2]                        (xd) {$\cdots$};
    \node[frozen, minimum width=12mm, right=2.5mm of xd] (x3) {$x_t$};
    \begin{scope}[on background layer]
      \node[grp, fit=(x1)(x3)] (pfx) {};
    \end{scope}
    \draw[data] (tab.east) -- (pfx.west);
    \node[tag, below=1.2mm of pfx] (pfxlab)
      {prefix latents --- \emph{order preserved}};
    % ---- the recurrence --------------------------------------------------
    \node[trained, minimum width=12mm, above=7mm of x1] (g1) {GRU};
    \node[trained, minimum width=12mm, above=7mm of x2] (g2) {GRU};
    \node[tag] (gd) at (g1 -| xd) {$\cdots$};
    \node[trained, minimum width=12mm, above=7mm of x3] (g3) {GRU};
    \draw[flow] (x1) -- (g1);
    \draw[flow] (x2) -- (g2);
    \draw[flow] (x3) -- (g3);
    \draw[flow] (g1) -- node[edg, above, inner sep=0.7pt] {$h_t$} (g2);
    \draw[flow] (g2) -- (gd);
    \draw[flow] (gd) -- (g3);
    % ---- head and retrieval ---------------------------------------------
    \node[trained, right=6mm of g3] (head)
      {linear head\\{\scriptsize$256\to192$}};
    \node[op, right=7mm of head] (rk)
      {rank the catalogue by $\cos(\hat z, z_v)$\\{\scriptsize prefix excluded $\to$ top-10}};
    \draw[flow] (g3) -- (head);
    \draw[flow] (head) -- node[edg, above, inner sep=1pt] {$\hat z$} (rk);
    % ---- the SAME frozen table is the retrieval target -------------------
    \coordinate (c1) at ([yshift=-3mm]tab.south |- current bounding box.south);
    \coordinate (c2) at (rk.south |- c1);
    \draw[data] (tab.south) -- (c1);
    \draw[data] (c1) -- node[tag, above]
      {the \emph{same} frozen table is the retrieval target} (c2);
    \draw[data] (c2) -- (rk.south);
    % ---- supervision note ------------------------------------------------
    \node[note] (nt) at ([yshift=-3.5mm]current bounding box.south west)
      {\textbf{Supervision --- per-step teacher forcing.} The head is shared
       across positions and must predict the next item's latent at
       \emph{every} prefix position, so one session of length $L$ yields
       $L-1$ examples: $73{,}632$ (prefix, next-item) pairs on the train
       split. The loss is masked in-batch InfoNCE at $\tau=0.07$ --- each
       step's true next latent must outscore the other steps' targets in the
       same minibatch, which supplies negatives for free.};
    % ---- LEGEND (two rows so it cannot widen the picture) ----------------
    \node[trained, swatch, anchor=north west] (l1)
      at ([yshift=-3mm]current bounding box.south west) {};
    \node[tag, right=1.1mm of l1] (t1) {trained (solid)};
    \node[frozen, swatch, right=3.4mm of t1] (l2) {};
    \node[tag, right=1.1mm of l2] (t2) {frozen / parameter-free (dashed)};
    \node[lookup, swatch, right=3.4mm of t2] (l3) {};
    \node[tag, right=1.1mm of l3] (t3) {table lookup};
    \node[learnspace, swatch, anchor=north west] (l4)
      at ([yshift=-1.6mm]l1.south west) {};
    \node[tag, right=1.1mm of l4] (t4) {parameter-free, learned space};
    \node[won, swatch, right=3.4mm of t4] (l5) {};
    \node[tag, right=1.1mm of l5] (t5) {favoured variant};
    \node[refuted, swatch, right=3.4mm of t5] (l6) {};
    \node[tag, right=1.1mm of l6] (t6) {refuted variant};
  \end{tikzpicture}
  \caption{\texttt{seq-nexttrack}, the single-model baseline every other
    topology in this document is measured against: a frozen item space, a
    recurrence over the \emph{ordered} prefix, a linear head, and cosine
    retrieval back into the \emph{same} frozen space. Nothing about the item
    geometry is learned here --- the $19{,}402\times192$ latent table (offline
    PCA-192 on \dsid{seq-20260715-131139}) is read twice, once to embed the
    prefix and once as the retrieval target, and only the GRU and its linear
    head carry parameters (\num{394944} at $h=256$). \textbf{The visual
    language used throughout this document is fixed here.} A \emph{solid}
    border marks a trained component and a \emph{dashed} border a frozen or
    parameter-free one; that distinction rides on the line style, not the
    colour, so it survives greyscale. Colour then carries the verdict ---
    blue: the on-record default path; purple dashed: a non-parametric table
    lookup over the dataset; dashed blue: parameter-free but operating inside
    a \emph{learned} geometry; green: a variant the record favours; red: a
    variant the record refutes. On the canonical split this configuration
    scores Recall@10 \num{0.12299} (\num{176} of \num{1431}), reproduced
    bit-identically eleven times.}
  \label{fig:topo-gru}
\end{figure}


% ===========================================================================
% 2. seq-ann --- the pooled feed-forward twin
% ===========================================================================
\begin{figure}[ht]
  \centering
  \begin{tikzpicture}[topo]
    \node[frozen, minimum width=20mm, minimum height=12mm] (tab)
      {frozen item\\latents $Z$\\[1pt]{\scriptsize$19{,}402\times192$}};
    \node[frozen, minimum width=12mm, right=7mm of tab]  (x1) {$x_1$};
    \node[frozen, minimum width=12mm, right=4mm of x1]   (x2) {$x_2$};
    \node[tag, right=2.5mm of x2]                        (xd) {$\cdots$};
    \node[frozen, minimum width=12mm, right=2.5mm of xd] (x3) {$x_t$};
    \begin{scope}[on background layer]
      \node[grp, fit=(x1)(x3)] (pfx) {};
    \end{scope}
    \draw[data] (tab.east) -- (pfx.west);
    % ---- the two parameter-free summaries of the same prefix -------------
    \node[op, right=8mm of pfx, yshift=10mm] (mp)
      {causal mean-pool\\$\frac{1}{t}\sum_{i\le t}x_i$\\[1pt]{\scriptsize\emph{permutation-invariant}}};
    \node[op, right=8mm of pfx, yshift=-10mm] (lt)
      {last item\\$x_t$\\[1pt]{\scriptsize the only recency cue}};
    \draw[flow] (pfx.east) -- (mp.west);
    \draw[flow] (pfx.east) -- (lt.west);
    % ---- the MLP ---------------------------------------------------------
    \node[trained, right=5mm of mp, yshift=-10mm] (mlp)
      {MLP\\$384\to256\to256\to192$\\[1pt]{\scriptsize ReLU $+$ dropout}};
    \draw[flow] (mp.east) -- (mlp.north west);
    \draw[flow] (lt.east) -- (mlp.south west);
    \node[edg] at ([xshift=-4.2mm]mlp.west) {$\|$};
    % ---- retrieval, placed below to keep the picture inside the text ------
    \node[op, below=9mm of mlp] (rk)
      {$\hat z\to$ cosine rank\\{\scriptsize prefix excluded $\to$ top-10}};
    \draw[flow] (mlp) -- (rk);
    \coordinate (c1) at ([yshift=-3mm]tab.south |- current bounding box.south);
    \coordinate (c2) at (rk.south |- c1);
    \draw[data] (tab.south) -- (c1);
    \draw[data] (c1) -- node[tag, above]
      {the same frozen table is the retrieval target} (c2);
    \draw[data] (c2) -- (rk.south);
    \node[note] at ([yshift=-3.5mm]current bounding box.south west)
      {\textbf{The same information, with the order discarded.} Both summaries
       are strictly causal, so the objective, the padded batching and the
       early stopping are reused \emph{verbatim} from
       Figure~\ref{fig:topo-gru} --- but a mean weights the track from twenty
       steps ago exactly like the one before last, and cannot represent ``A
       then B'' differently from ``B then A''. Only the special-cased $x_t$ tap
       carries any notion of recency, and the concatenation ($\|$) is this
       model's entire account of sequence.};
  \end{tikzpicture}
  \caption{\texttt{seq-ann}, the GRU's non-recurrent twin: the recurrence of
    Figure~\ref{fig:topo-gru} is replaced by two parameter-free summaries of
    the same prefix --- a causal mean-pool concatenated with the last item's
    latent --- feeding a plain three-layer ReLU MLP. Everything else is
    literally the same code (\texttt{seq\_ann.py} imports
    \texttt{make\_batches} and \texttt{step\_loss} from
    \texttt{seq\_nexttrack}), which is what makes this pair a clean test of
    the recurrence itself rather than of a training recipe. Visual convention
    as in Figure~\ref{fig:topo-gru}: the two pooling operators are dashed
    because they hold no parameters, and only the MLP is trained. The
    defensible comparison is same-space: on \dsid{seq-20260715-030509}
    (PCA-128) the pooled twin scores Recall@10 \num{0.0901} (\num{129} of
    \num{1431}) against the GRU's \num{0.1146} (\num{164}), paired $\Delta$
    $-\num{0.0245}$ $[-0.039, -0.011]$, a CI \emph{entirely below zero}.
    \texttt{seq-ann} has no run at all on the canonical PCA-192 split, so it
    never belongs in the same column as the \num{0.12299} bar.}
  \label{fig:topo-ann}
\end{figure}


% ===========================================================================
% 3. seq-dualgru --- two towers, a fusion head, and the pre-MLP variant
% ===========================================================================
\begin{figure}[ht]
  \centering
  \begin{tikzpicture}[topo]
    \node[frozen, minimum width=19mm, minimum height=11mm] (pf)
      {prefix\\latents\\$x_1\dots x_t$};
    % ---- the three configurable causal views ----------------------------
    \node[op, right=7mm of pf] (v2) {\texttt{delta}\\$x_t-x_{t-1}$};
    \node[op, above=2.6mm of v2] (v1) {\texttt{latent}\\$x_t$};
    \node[op, below=2.6mm of v2] (v3) {\texttt{cummean}\\$\frac{1}{t}\sum_{i\le t}x_i$};
    \begin{scope}[on background layer]
      \node[grp, fit=(v1)(v3)] (vg) {};
    \end{scope}
    \draw[data] (pf.east) -- (vg.west);
    % ---- two independent towers -----------------------------------------
    \node[trained, right=8mm of vg, yshift=9mm] (ta)
      {tower A on \texttt{view\_a}\\GRU $h_a{=}256$\\[1pt]{\scriptsize own weights}};
    \node[trained, right=8mm of vg, yshift=-9mm] (tb)
      {tower B on \texttt{view\_b}\\GRU $h_b{=}256$\\[1pt]{\scriptsize own weights}};
    \draw[flow] (vg.east) -- (ta.west);
    \draw[flow] (vg.east) -- (tb.west);
    % ---- concat + the two fusion depths ---------------------------------
    \node[op, right=7mm of ta, yshift=-9mm] (cc)
      {concat\\$h_a\|h_b$\\[1pt]{\scriptsize$512$}};
    \draw[flow] (ta.east) -- (cc.north west);
    \draw[flow] (tb.east) -- (cc.south west);
    \node[won, right=7mm of cc, yshift=9mm] (f0)
      {\texttt{fusion\_layers}${}=0$\\$512\to192$, linear};
    \node[refuted, right=7mm of cc, yshift=-9mm] (f1)
      {\texttt{fusion\_layers}${}=1$\\$512\to256\to192$\\[1pt]{\scriptsize$+$ ReLU $+$ dropout}};
    \draw[flow] (cc.north east) -- (f0.west);
    \draw[flow] (cc.south east) -- (f1.west);
    % ---- the refuted per-step pre-MLP variant ---------------------------
    \node[refuted, below=12mm of tb] (pb)
      {\emph{variant} --- per-step pre-MLP\\on a tower:
       Linear$(192{\to}256)$\\$+$ ReLU $+$ Dropout};
    \draw[variant] (vg.south) |- (pb.west);
    \draw[variant] (pb.north) -- (tb.south);
    \node[note] at ([yshift=-3.5mm]current bounding box.south west)
      {\textbf{Both heads emit the same $192$-d $\hat z$ and retrieve exactly
       as in Figure~\ref{fig:topo-gru}.} The towers share nothing --- separate
       modules, separate parameters --- and their per-step hidden states are
       concatenated, so the combiner sees the towers' \emph{representations}
       rather than their scores and gradients flow back into both. There is no
       bidirectional mode by construction: a backward pass would let step $t$
       read item $t+1$ and copy its own teacher-forcing target. All three views
       are strictly causal for the same reason, and so is the pre-MLP, which is
       applied \emph{timestep-wise}; it is drawn solid because it is trained,
       red because it is refuted.};
  \end{tikzpicture}
  \caption{\texttt{seq-dualgru}: two recurrent towers with entirely separate
    weights read the prefix in parallel and a head fuses their hidden states
    --- the strongest available form of ``train the combination end-to-end so
    the parts co-adapt''. Each tower's input is a configurable causal view
    (\texttt{view\_a} / \texttt{view\_b}); the default is
    \texttt{latent}/\texttt{latent}, i.e.\ two independent GRUs over the same
    stream. Two of the boxes are coloured by the record rather than by the
    design. Removing the MLP fusion head (\texttt{fusion\_layers}${}=0$,
    green) \emph{beats} keeping it in three independent measurements, most
    recently by $+\num{0.02655}$ $[+0.01398, +0.03913]$ CI$>$0 --- and with
    some $9$--$10\%$ \emph{fewer} parameters, so the linear readout is the
    family default; and the
    per-step pre-MLP (red) loses the decisive single-tower ablation, Recall@10
    \num{0.10901} (\num{156} of \num{1431}) against the baseline's \num{176},
    paired $\Delta$ $-\num{0.01398}$ $[-0.02657, -0.00140]$ CI$<$0. Visual
    convention as in Figure~\ref{fig:topo-gru}. The best dual arm reaches
    \num{0.11880} and merely \emph{ties} the single-GRU bar, as does a
    $0.098\%$ parameter-matched width control --- so capacity is not the
    explanation.}
  \label{fig:topo-dualgru}
\end{figure}


% ===========================================================================
% 4. seq-embed --- a learned, weight-tied item table (DESIGN ONLY)
% ===========================================================================
\begin{figure}[ht]
  \centering
  \begin{tikzpicture}[topo]
    \node[frozen, minimum width=20mm, minimum height=11mm] (ids)
      {session prefix\\item \emph{ids}\\$i_1\dots i_t$};
    \node[trained, right=10mm of ids, minimum height=11mm] (E)
      {learned item\\table $E$\\[1pt]{\scriptsize$19{,}402\times192$}};
    \node[trained, right=6mm of E] (gru)
      {GRU $h{=}256$\\{\scriptsize$+$ dropout}};
    \node[trained, right=6mm of gru] (hd)
      {linear head\\{\scriptsize$256\to192$}};
    \node[op, right=6mm of hd] (rk)
      {$\cos(\hat z, E_v)$ over all\\$19{,}402$ rows $\to$ top-10};
    \draw[flow] (ids) -- node[edg, above, inner sep=1pt] {gather} (E);
    \draw[flow] (E) -- (gru);
    \draw[flow] (gru) -- (hd);
    \draw[flow] (hd) -- node[edg, above, inner sep=1pt] {$\hat z$} (rk);
    % ---- the three modes -------------------------------------------------
    \node[op, text width=64mm, align=flush left, anchor=south west] (md)
      at ([xshift=4mm,yshift=8mm]E.north east)
      {\texttt{embed\_mode} --- what the table $E$ actually is:\\
       \textbf{\texttt{frozen}} $E=Z$, nothing trained (the reproduction gate)\\
       \textbf{\texttt{residual}} $E=Z+\Delta E$: frozen content vector $+$ a
       trained correction, so a cold item still has a position\\
       \textbf{\texttt{free}} $E$ fully trained, optionally randomly initialised};
    \draw[flow] (md.west) -- (E.north east);
    % ---- weight tying: the same table is the retrieval target -----------
    \coordinate (c1) at ([yshift=-8mm]E.south);
    \coordinate (c2) at (rk.south |- c1);
    \draw[tie] (E.south) -- (c1);
    \draw[tie] (c1) -- node[tag, above, text=topogood]
      {weight tying --- the \emph{same} table is\\the input embedding
       \emph{and} the retrieval target} (c2);
    \draw[tie] (c2) -- (rk.south);
    % ---- the gradient loop that makes the geometry trainable ------------
    \coordinate (t1) at ([yshift=5mm]rk.north |- current bounding box.north);
    \coordinate (t2) at (E.north |- t1);
    \draw[grad] (rk.north) -- (t1);
    \draw[grad] (t1) -- node[tag, above, text=topobad]
      {$\nabla$ --- one InfoNCE step moves the query \emph{and} every target:
       the geometry itself is trainable} (t2);
    \draw[grad] (t2) -- (E.north);
    \node[note] at ([yshift=-3.5mm]current bounding box.south west)
      {\textbf{Why the loop is the whole point.} In
       Figures~\ref{fig:topo-gru}--\ref{fig:topo-dualgru} the retrieval space
       is fixed and only the map into it is learned; here the space is a
       parameter, shaped by next-track adjacency rather than inherited from a
       reconstruction objective. That also creates a collapse hazard absent
       everywhere else: with a learned table a plain cosine objective has a
       trivial optimum at ``send every item to one point'', so only a
       contrastive loss is admissible and \texttt{loss=cosine} is rejected
       outright for the trained modes.};
  \end{tikzpicture}
  \caption{\texttt{seq-embed}, the one axis the record has not tested:
    \textbf{a design, not a result}. No run of this predictor exists on any
    dataset, it is absent from the server's registered predictor list, and
    nothing here is a measured score. The item table $E$ is used twice, as the
    GRU's input embedding and as the retrieval target (weight tying, the
    GRU4Rec/SASRec formulation), so a single gradient step moves the query and
    all the candidates together --- the return arrow along the top is what
    makes the geometry trainable, and it is exactly what every other topology
    in this document lacks. The three \texttt{embed\_mode} settings span the
    space from a reproduction gate to a fully free table. What \emph{is}
    measured is the data budget that bounds them: \num{73632} training steps
    to fit \num{15816} item vectors, a median per-item training frequency of
    \num{1}, only \num{3294} items occurring five or more times, and
    \num{788} of \num{1431} test targets never appearing in any training
    session --- so a \texttt{free} table cannot retrieve the cold
    \SI{55}{\percent} at all and is capped at Recall@10 \num{0.4493}, whereas
    \texttt{residual} keeps the content vector and with it the cold reach.
    Visual convention as in Figure~\ref{fig:topo-gru}.}
  \label{fig:topo-embed}
\end{figure}


% ===========================================================================
% 5. seq-blend --- the champion
% ===========================================================================
\begin{figure}[ht]
  \centering
  \begin{tikzpicture}[topo]
    \node[frozen, minimum width=18mm, minimum height=11mm] (pf)
      {session\\prefix\\$i_1\dots i_t$};
    % ---- the learned metric map: shared by two legs, bypassed by one -----
    \node[trained, right=4mm of pf] (W)
      {learned map $W$\\$192\times192$\\[1pt]{\scriptsize no bias, identity init}\\{\scriptsize
       InfoNCE $\tau{=}0.07$}\\{\scriptsize $36{,}864$ params}};
    % ---- the three legs -------------------------------------------------
    \node[learnspace, right=3mm of W] (C)
      {\Cp\ content-$k$NN in $W$-space\\$\max_{u}\cos(Wz_v,Wz_u)$\\[1pt]{\scriptsize
       no parameters of its own}};
    \node[trained, above=7mm of C] (R)
      {\Rp\ projected GRU\\$h{=}256$, $394{,}944$ params};
    \node[lookup, below=7mm of C] (M)
      {$M$ first-order Markov counts\\$P(v\mid u)$, $u=$ last prefix item\\[1pt]{\scriptsize
       Dirichlet trust $n_u/(n_u{+}8)$,}\\{\scriptsize artist/genre back-off; $0$ params}};
    \draw[flow] (pf.east) -- (W.west);
    % ---- the frozen item table W acts ON: the same Z as Figure~\ref{fig:topo-gru} --
    \node[frozen, anchor=south east, minimum width=18mm] (Z)
      at ([yshift=6mm]W.north west)
      {frozen item latents $Z$\\{\scriptsize$19{,}402\times192$}};
    \draw[data] (Z.south east) -- node[edg, left, inner sep=1.2pt] {$z_i$} (W.north west);
    \draw[proj] (W.north east) -- (R.west);
    \draw[proj] (W.east) -- (C.west);
    \draw[data] (pf.south) |- node[tag, above, pos=0.55]
      {item \emph{ids} only ---\\the $M$ leg bypasses $Z$ and $W$} (M.west);
    % ---- the fixed combiner ---------------------------------------------
    \node[op, right=3mm of C] (cb)
      {z-normalise each leg\\over the candidate set,\\then average\\[1pt]{\scriptsize
       equal thirds ---}\\{\scriptsize\emph{no learned weights}}};
    \draw[flow] (R.east) -- node[edg, above, sloped, pos=0.55] {$s_{\Rp}$} (cb.north west);
    \draw[flow] (C.east) -- node[edg, above, inner sep=1pt] {$s_{\Cp}$} (cb.west);
    \draw[flow] (M.east) -- node[edg, above, sloped, pos=0.30] {$s_M$} (cb.south west);
    \node[op, right=2mm of cb] (rk) {rank\\$\to$ top-10};
    \draw[flow] (cb) -- (rk);
    \node[note] at ([yshift=-3.5mm]current bounding box.south west)
      {\textbf{Which legs are learned.} $W$ and the \Rp\ recurrence are
       trained (\num{431808} parameters between them); \Cp\ has no parameters
       of its own but retrieves inside $W$'s learned geometry, which is why it
       is drawn dashed-blue; $M$ is a table of training-set transition counts,
       drawn dashed-purple. The combiner is a fixed rule: each leg's
       full-vocabulary score vector is z-scored over the admissible candidates
       (every item minus those already in the prefix) and the three are
       averaged with hard-coded equal weights. The z-scoring is the
       load-bearing step --- it puts three incomparable score scales onto one
       per-session scale while \emph{preserving} relative magnitude, so a
       confidently-peaked leg can outvote two diffuse ones. Rank-only fusion,
       which discards that magnitude, loses by $-\num{0.029}$ CI$<$0.};
  \end{tikzpicture}
  \caption{\texttt{seq-blend} with \texttt{projection=true}, the champion:
    three full-vocabulary score vectors, z-normalised over the candidate set
    and averaged in fixed equal thirds. $W$ acts on the \emph{frozen item
    latents} $Z$ of Figure~\ref{fig:topo-gru} --- it warps the retrieval geometry,
    it does not read the prefix ids directly --- and it is shared by two of the
    three legs and bypassed by the third; that is the whole mechanism. $W$ is fit by in-batch InfoNCE on consecutive within-train
    pairs, which rebuilds the retrieval metric from ``reconstruct the item
    matrix'' into ``what follows what''; both legs that retrieve by cosine in
    that space improve --- the parameter-free content-$k$NN leg more than
    doubles, \num{0.0685} to \num{0.1516}, and the recurrent leg rises from
    \num{0.1230} to $\approx$\num{0.172} --- while the item-ID lookup is
    untouched at \num{0.106918}. The two projected-leg standalone figures come
    from the off-server driver of the 2026-07-18 campaign and have no server
    run of their own; only the $M$ leg's number is a canonical-split run.
    Together the three legs reach Recall@10 \num{0.21174} (\num{303} of
    \num{1431}), paired $\Delta$ over the unprojected three-leg blend
    $+\num{0.0259}$ $[+0.012, +0.041]$, a CI entirely above zero and robust
    across three seeds. Visual convention as in Figure~\ref{fig:topo-gru}.}
  \label{fig:topo-blend}
\end{figure}


% ===========================================================================
% 6. seq-stacker --- the learned combiner
% ===========================================================================
\begin{figure}[ht]
  \centering
  \begin{tikzpicture}[topo]
    \node[frozen, minimum width=18mm, minimum height=11mm] (pf)
      {session\\prefix\\$i_1\dots i_t$};
    \node[trained, right=6mm of pf] (R)
      {$R$ GRU $h{=}256$\\[1pt]{\scriptsize trained separately}};
    \node[lookup, above=6mm of R] (M) {$M$ Markov counts};
    \node[trained, below=6mm of R] (A)
      {$A$ pooled MLP\\[1pt]{\scriptsize trained separately}};
    \draw[flow] (pf.east) -- (M.west);
    \draw[flow] (pf.east) -- (R.west);
    \draw[flow] (pf.east) -- (A.west);
    % ---- the per-candidate feature row ----------------------------------
    \node[op, right=6mm of R] (fx)
      {per-candidate feature row\\$[\,s_M,\ s_R,\ s_A,\ \log(1{+}\text{plays}),$\\$\text{recency},\
       \text{artist match},\ \text{genre match}\,]$};
    \draw[flow] (M.east) -- (fx.north west);
    \draw[flow] (R.east) -- (fx.west);
    \draw[flow] (A.east) -- (fx.south west);
    \node[refuted, right=5mm of fx] (xgb)
      {XGBoost meta-learner\\$200$ rounds, depth $6$\\[1pt]{\scriptsize\texttt{binary:logistic}}};
    \draw[flow] (fx) -- (xgb);
    \node[op, below=9mm of xgb] (rk)
      {score all $19{,}402$ items\\$\to$ top-10};
    \draw[flow] (xgb) -- (rk);
    % ---- the two failure mechanisms, marked on the diagram --------------
    \node[refuted, text width=56mm, align=flush left, anchor=north west] (tr)
      at ([yshift=-10mm]fx.south west)
      {\textbf{Training rows}: $1$ positive $+$ $50$ popularity-sampled
       negatives per train session, i.e.\ \num{5723} positives, against the
       \num{73632} teacher-forced steps its own $R$ leg trains on
       ($12.9\times$ less signal) --- and a decision surface calibrated on
       $51$ candidates ($0.26\%$ of the catalogue) is then asked to rank all
       \num{19402}.};
    \draw[grad] (tr.north east) -- (xgb.south west);
    \node[note] at ([yshift=-3.5mm]current bounding box.south west)
      {\textbf{Why the shape matters.} A $200$-round depth-$6$ booster is
       piecewise \emph{constant}, so ranking \num{19402} candidates with it
       produces large tied blocks that the harness breaks by item index. The
       fixed z-blend of Figure~\ref{fig:topo-blend} is instead a strictly
       monotone continuous function of the leg scores: it preserves each leg's
       internal ordering exactly and can only trade the legs off against one
       another.};
  \end{tikzpicture}
  \caption{\texttt{seq-stacker}, the learned counterpart of
    Figure~\ref{fig:topo-blend}'s fixed rule: the same separately-trained base
    legs, but their per-candidate scores become columns of a feature row
    alongside four cheap candidate features, and a gradient-boosted tree
    learns the combination. It earns a diagram because its failure is
    informative --- with the \texttt{markov}, \texttt{gru} and \texttt{ann}
    legs it reaches Recall@10 \num{0.085255} (\num{122} of \num{1431}),
    \emph{below its own GRU base leg} at \num{0.12299}, paired $\Delta$
    $-\num{0.0377}$ $[-0.056, -0.020]$ CI$<$0. The red callout names the two
    mechanisms: supervision starvation, since the meta-learner gets one
    labelled event per session where its base leg gets one per prefix
    position, and a train/serve candidate mismatch. The booster is drawn solid
    because it is trained and red because it is refuted; the feature row and
    the ranking step are dashed because they hold no parameters. Visual
    convention as in Figure~\ref{fig:topo-gru}. This is one of four
    independent learned-combiner nulls, the others being reciprocal-rank
    fusion, oracle weights fitted on the test labels, and the jointly-trained
    dual tower of Figure~\ref{fig:topo-dualgru}.}
  \label{fig:topo-stacker}
\end{figure}

%
%  (\graphicspath{{figures/}} governs \includegraphics only, NOT \input, so the
%   figures/ prefix is spelled out above.)
%
%  ------------------------------------------------------------------------
%  REQUIRED PREAMBLE LINES in the main document --- add these two lines after
%  \usepackage{graphicx}, before \begin{document}:
%
%      \usepackage{tikz}
%      \usetikzlibrary{positioning,arrows.meta,fit,backgrounds,calc}
%
%  Nothing else is required: xcolor arrives with tikz, and every colour and
%  every style used below is defined in this file.  \dsid, \Rp and \Cp are
%  \providecommand-guarded here so the file also compiles inside a bare
%  wrapper; when the main document defines them, its definitions win.
%
%  PACKAGES DELIBERATELY *NOT* REQUIRED: amssymb is absent from the house
%  preamble, so nothing below uses \mathbb.  Every picture is laid out with
%  RELATIVE positioning (right=/above=/below= of ...), never with hard
%  coordinates, so no node can collide with its neighbour if the font or the
%  wording changes; the per-figure notes and the legend are anchored to
%  `current bounding box', so they always land underneath the diagram.
%  ------------------------------------------------------------------------
%
%  LABELS DEFINED IN THIS FILE (in file order)
%      fig:topo-gru       seq-nexttrack --- the baseline; ALSO CARRIES THE LEGEND
%      fig:topo-ann       seq-ann       --- pooled feed-forward twin
%      fig:topo-dualgru   seq-dualgru   --- two towers + fusion head + pre-MLP variant
%      fig:topo-embed     seq-embed     --- learned, weight-tied item table (DESIGN ONLY)
%      fig:topo-blend     seq-blend     --- the champion, three legs
%      fig:topo-stacker   seq-stacker   --- learned XGBoost combiner
%
%  VISUAL LANGUAGE (the primary axis is the LINE, not the colour, so it
%  survives greyscale printing; stated for the reader in the caption of
%  fig:topo-gru and re-pointed-at by every other caption)
%      SOLID  border  = TRAINED component (parameters fitted by gradient descent)
%      DASHED border  = FROZEN or PARAMETER-FREE component
%      blue   = the on-record default path
%      purple dashed  = non-parametric table lookup over the dataset
%      dashed blue    = parameter-free, but operating inside a LEARNED geometry
%      green  = variant favoured by the record
%      red    = variant refuted by the record
% ===========================================================================

% ===========================================================================
%  topo-style.tex --- the shared TikZ style block for the next-track topology
%  diagrams AND for the evaluation-protocol diagram in the body of
%  model-comparison.tex.  Extracted from topologies.tex so that the protocol
%  figure, which appears BEFORE % ===========================================================================
%  topologies.tex --- TikZ network-topology diagrams for the next-track
%  MODEL-COMPARISON document.  Not a standalone document: \input it from the
%  body of the main .tex, which must be compiled from docs/:
%
%      \input{figures/topologies}
%
%  (\graphicspath{{figures/}} governs \includegraphics only, NOT \input, so the
%   figures/ prefix is spelled out above.)
%
%  ------------------------------------------------------------------------
%  REQUIRED PREAMBLE LINES in the main document --- add these two lines after
%  \usepackage{graphicx}, before \begin{document}:
%
%      \usepackage{tikz}
%      \usetikzlibrary{positioning,arrows.meta,fit,backgrounds,calc}
%
%  Nothing else is required: xcolor arrives with tikz, and every colour and
%  every style used below is defined in this file.  \dsid, \Rp and \Cp are
%  \providecommand-guarded here so the file also compiles inside a bare
%  wrapper; when the main document defines them, its definitions win.
%
%  PACKAGES DELIBERATELY *NOT* REQUIRED: amssymb is absent from the house
%  preamble, so nothing below uses \mathbb.  Every picture is laid out with
%  RELATIVE positioning (right=/above=/below= of ...), never with hard
%  coordinates, so no node can collide with its neighbour if the font or the
%  wording changes; the per-figure notes and the legend are anchored to
%  `current bounding box', so they always land underneath the diagram.
%  ------------------------------------------------------------------------
%
%  LABELS DEFINED IN THIS FILE (in file order)
%      fig:topo-gru       seq-nexttrack --- the baseline; ALSO CARRIES THE LEGEND
%      fig:topo-ann       seq-ann       --- pooled feed-forward twin
%      fig:topo-dualgru   seq-dualgru   --- two towers + fusion head + pre-MLP variant
%      fig:topo-embed     seq-embed     --- learned, weight-tied item table (DESIGN ONLY)
%      fig:topo-blend     seq-blend     --- the champion, three legs
%      fig:topo-stacker   seq-stacker   --- learned XGBoost combiner
%
%  VISUAL LANGUAGE (the primary axis is the LINE, not the colour, so it
%  survives greyscale printing; stated for the reader in the caption of
%  fig:topo-gru and re-pointed-at by every other caption)
%      SOLID  border  = TRAINED component (parameters fitted by gradient descent)
%      DASHED border  = FROZEN or PARAMETER-FREE component
%      blue   = the on-record default path
%      purple dashed  = non-parametric table lookup over the dataset
%      dashed blue    = parameter-free, but operating inside a LEARNED geometry
%      green  = variant favoured by the record
%      red    = variant refuted by the record
% ===========================================================================

\input{figures/topo-style}



% ===========================================================================
% 1. seq-nexttrack --- the baseline (and the legend)
% ===========================================================================
\begin{figure}[ht]
  \centering
  \begin{tikzpicture}[topo]
    % ---- the frozen item-latent table, read TWICE ------------------------
    \node[frozen, minimum width=20mm, minimum height=12mm] (tab)
      {frozen item\\latents $Z$\\[1pt]{\scriptsize$19{,}402\times192$}};
    % ---- the prefix, in order (one column per step) ----------------------
    \node[frozen, minimum width=12mm, right=7mm of tab]  (x1) {$x_1$};
    \node[frozen, minimum width=12mm, right=4mm of x1]   (x2) {$x_2$};
    \node[tag, right=2.5mm of x2]                        (xd) {$\cdots$};
    \node[frozen, minimum width=12mm, right=2.5mm of xd] (x3) {$x_t$};
    \begin{scope}[on background layer]
      \node[grp, fit=(x1)(x3)] (pfx) {};
    \end{scope}
    \draw[data] (tab.east) -- (pfx.west);
    \node[tag, below=1.2mm of pfx] (pfxlab)
      {prefix latents --- \emph{order preserved}};
    % ---- the recurrence --------------------------------------------------
    \node[trained, minimum width=12mm, above=7mm of x1] (g1) {GRU};
    \node[trained, minimum width=12mm, above=7mm of x2] (g2) {GRU};
    \node[tag] (gd) at (g1 -| xd) {$\cdots$};
    \node[trained, minimum width=12mm, above=7mm of x3] (g3) {GRU};
    \draw[flow] (x1) -- (g1);
    \draw[flow] (x2) -- (g2);
    \draw[flow] (x3) -- (g3);
    \draw[flow] (g1) -- node[edg, above, inner sep=0.7pt] {$h_t$} (g2);
    \draw[flow] (g2) -- (gd);
    \draw[flow] (gd) -- (g3);
    % ---- head and retrieval ---------------------------------------------
    \node[trained, right=6mm of g3] (head)
      {linear head\\{\scriptsize$256\to192$}};
    \node[op, right=7mm of head] (rk)
      {rank the catalogue by $\cos(\hat z, z_v)$\\{\scriptsize prefix excluded $\to$ top-10}};
    \draw[flow] (g3) -- (head);
    \draw[flow] (head) -- node[edg, above, inner sep=1pt] {$\hat z$} (rk);
    % ---- the SAME frozen table is the retrieval target -------------------
    \coordinate (c1) at ([yshift=-3mm]tab.south |- current bounding box.south);
    \coordinate (c2) at (rk.south |- c1);
    \draw[data] (tab.south) -- (c1);
    \draw[data] (c1) -- node[tag, above]
      {the \emph{same} frozen table is the retrieval target} (c2);
    \draw[data] (c2) -- (rk.south);
    % ---- supervision note ------------------------------------------------
    \node[note] (nt) at ([yshift=-3.5mm]current bounding box.south west)
      {\textbf{Supervision --- per-step teacher forcing.} The head is shared
       across positions and must predict the next item's latent at
       \emph{every} prefix position, so one session of length $L$ yields
       $L-1$ examples: $73{,}632$ (prefix, next-item) pairs on the train
       split. The loss is masked in-batch InfoNCE at $\tau=0.07$ --- each
       step's true next latent must outscore the other steps' targets in the
       same minibatch, which supplies negatives for free.};
    % ---- LEGEND (two rows so it cannot widen the picture) ----------------
    \node[trained, swatch, anchor=north west] (l1)
      at ([yshift=-3mm]current bounding box.south west) {};
    \node[tag, right=1.1mm of l1] (t1) {trained (solid)};
    \node[frozen, swatch, right=3.4mm of t1] (l2) {};
    \node[tag, right=1.1mm of l2] (t2) {frozen / parameter-free (dashed)};
    \node[lookup, swatch, right=3.4mm of t2] (l3) {};
    \node[tag, right=1.1mm of l3] (t3) {table lookup};
    \node[learnspace, swatch, anchor=north west] (l4)
      at ([yshift=-1.6mm]l1.south west) {};
    \node[tag, right=1.1mm of l4] (t4) {parameter-free, learned space};
    \node[won, swatch, right=3.4mm of t4] (l5) {};
    \node[tag, right=1.1mm of l5] (t5) {favoured variant};
    \node[refuted, swatch, right=3.4mm of t5] (l6) {};
    \node[tag, right=1.1mm of l6] (t6) {refuted variant};
  \end{tikzpicture}
  \caption{\texttt{seq-nexttrack}, the single-model baseline every other
    topology in this document is measured against: a frozen item space, a
    recurrence over the \emph{ordered} prefix, a linear head, and cosine
    retrieval back into the \emph{same} frozen space. Nothing about the item
    geometry is learned here --- the $19{,}402\times192$ latent table (offline
    PCA-192 on \dsid{seq-20260715-131139}) is read twice, once to embed the
    prefix and once as the retrieval target, and only the GRU and its linear
    head carry parameters (\num{394944} at $h=256$). \textbf{The visual
    language used throughout this document is fixed here.} A \emph{solid}
    border marks a trained component and a \emph{dashed} border a frozen or
    parameter-free one; that distinction rides on the line style, not the
    colour, so it survives greyscale. Colour then carries the verdict ---
    blue: the on-record default path; purple dashed: a non-parametric table
    lookup over the dataset; dashed blue: parameter-free but operating inside
    a \emph{learned} geometry; green: a variant the record favours; red: a
    variant the record refutes. On the canonical split this configuration
    scores Recall@10 \num{0.12299} (\num{176} of \num{1431}), reproduced
    bit-identically eleven times.}
  \label{fig:topo-gru}
\end{figure}


% ===========================================================================
% 2. seq-ann --- the pooled feed-forward twin
% ===========================================================================
\begin{figure}[ht]
  \centering
  \begin{tikzpicture}[topo]
    \node[frozen, minimum width=20mm, minimum height=12mm] (tab)
      {frozen item\\latents $Z$\\[1pt]{\scriptsize$19{,}402\times192$}};
    \node[frozen, minimum width=12mm, right=7mm of tab]  (x1) {$x_1$};
    \node[frozen, minimum width=12mm, right=4mm of x1]   (x2) {$x_2$};
    \node[tag, right=2.5mm of x2]                        (xd) {$\cdots$};
    \node[frozen, minimum width=12mm, right=2.5mm of xd] (x3) {$x_t$};
    \begin{scope}[on background layer]
      \node[grp, fit=(x1)(x3)] (pfx) {};
    \end{scope}
    \draw[data] (tab.east) -- (pfx.west);
    % ---- the two parameter-free summaries of the same prefix -------------
    \node[op, right=8mm of pfx, yshift=10mm] (mp)
      {causal mean-pool\\$\frac{1}{t}\sum_{i\le t}x_i$\\[1pt]{\scriptsize\emph{permutation-invariant}}};
    \node[op, right=8mm of pfx, yshift=-10mm] (lt)
      {last item\\$x_t$\\[1pt]{\scriptsize the only recency cue}};
    \draw[flow] (pfx.east) -- (mp.west);
    \draw[flow] (pfx.east) -- (lt.west);
    % ---- the MLP ---------------------------------------------------------
    \node[trained, right=5mm of mp, yshift=-10mm] (mlp)
      {MLP\\$384\to256\to256\to192$\\[1pt]{\scriptsize ReLU $+$ dropout}};
    \draw[flow] (mp.east) -- (mlp.north west);
    \draw[flow] (lt.east) -- (mlp.south west);
    \node[edg] at ([xshift=-4.2mm]mlp.west) {$\|$};
    % ---- retrieval, placed below to keep the picture inside the text ------
    \node[op, below=9mm of mlp] (rk)
      {$\hat z\to$ cosine rank\\{\scriptsize prefix excluded $\to$ top-10}};
    \draw[flow] (mlp) -- (rk);
    \coordinate (c1) at ([yshift=-3mm]tab.south |- current bounding box.south);
    \coordinate (c2) at (rk.south |- c1);
    \draw[data] (tab.south) -- (c1);
    \draw[data] (c1) -- node[tag, above]
      {the same frozen table is the retrieval target} (c2);
    \draw[data] (c2) -- (rk.south);
    \node[note] at ([yshift=-3.5mm]current bounding box.south west)
      {\textbf{The same information, with the order discarded.} Both summaries
       are strictly causal, so the objective, the padded batching and the
       early stopping are reused \emph{verbatim} from
       Figure~\ref{fig:topo-gru} --- but a mean weights the track from twenty
       steps ago exactly like the one before last, and cannot represent ``A
       then B'' differently from ``B then A''. Only the special-cased $x_t$ tap
       carries any notion of recency, and the concatenation ($\|$) is this
       model's entire account of sequence.};
  \end{tikzpicture}
  \caption{\texttt{seq-ann}, the GRU's non-recurrent twin: the recurrence of
    Figure~\ref{fig:topo-gru} is replaced by two parameter-free summaries of
    the same prefix --- a causal mean-pool concatenated with the last item's
    latent --- feeding a plain three-layer ReLU MLP. Everything else is
    literally the same code (\texttt{seq\_ann.py} imports
    \texttt{make\_batches} and \texttt{step\_loss} from
    \texttt{seq\_nexttrack}), which is what makes this pair a clean test of
    the recurrence itself rather than of a training recipe. Visual convention
    as in Figure~\ref{fig:topo-gru}: the two pooling operators are dashed
    because they hold no parameters, and only the MLP is trained. The
    defensible comparison is same-space: on \dsid{seq-20260715-030509}
    (PCA-128) the pooled twin scores Recall@10 \num{0.0901} (\num{129} of
    \num{1431}) against the GRU's \num{0.1146} (\num{164}), paired $\Delta$
    $-\num{0.0245}$ $[-0.039, -0.011]$, a CI \emph{entirely below zero}.
    \texttt{seq-ann} has no run at all on the canonical PCA-192 split, so it
    never belongs in the same column as the \num{0.12299} bar.}
  \label{fig:topo-ann}
\end{figure}


% ===========================================================================
% 3. seq-dualgru --- two towers, a fusion head, and the pre-MLP variant
% ===========================================================================
\begin{figure}[ht]
  \centering
  \begin{tikzpicture}[topo]
    \node[frozen, minimum width=19mm, minimum height=11mm] (pf)
      {prefix\\latents\\$x_1\dots x_t$};
    % ---- the three configurable causal views ----------------------------
    \node[op, right=7mm of pf] (v2) {\texttt{delta}\\$x_t-x_{t-1}$};
    \node[op, above=2.6mm of v2] (v1) {\texttt{latent}\\$x_t$};
    \node[op, below=2.6mm of v2] (v3) {\texttt{cummean}\\$\frac{1}{t}\sum_{i\le t}x_i$};
    \begin{scope}[on background layer]
      \node[grp, fit=(v1)(v3)] (vg) {};
    \end{scope}
    \draw[data] (pf.east) -- (vg.west);
    % ---- two independent towers -----------------------------------------
    \node[trained, right=8mm of vg, yshift=9mm] (ta)
      {tower A on \texttt{view\_a}\\GRU $h_a{=}256$\\[1pt]{\scriptsize own weights}};
    \node[trained, right=8mm of vg, yshift=-9mm] (tb)
      {tower B on \texttt{view\_b}\\GRU $h_b{=}256$\\[1pt]{\scriptsize own weights}};
    \draw[flow] (vg.east) -- (ta.west);
    \draw[flow] (vg.east) -- (tb.west);
    % ---- concat + the two fusion depths ---------------------------------
    \node[op, right=7mm of ta, yshift=-9mm] (cc)
      {concat\\$h_a\|h_b$\\[1pt]{\scriptsize$512$}};
    \draw[flow] (ta.east) -- (cc.north west);
    \draw[flow] (tb.east) -- (cc.south west);
    \node[won, right=7mm of cc, yshift=9mm] (f0)
      {\texttt{fusion\_layers}${}=0$\\$512\to192$, linear};
    \node[refuted, right=7mm of cc, yshift=-9mm] (f1)
      {\texttt{fusion\_layers}${}=1$\\$512\to256\to192$\\[1pt]{\scriptsize$+$ ReLU $+$ dropout}};
    \draw[flow] (cc.north east) -- (f0.west);
    \draw[flow] (cc.south east) -- (f1.west);
    % ---- the refuted per-step pre-MLP variant ---------------------------
    \node[refuted, below=12mm of tb] (pb)
      {\emph{variant} --- per-step pre-MLP\\on a tower:
       Linear$(192{\to}256)$\\$+$ ReLU $+$ Dropout};
    \draw[variant] (vg.south) |- (pb.west);
    \draw[variant] (pb.north) -- (tb.south);
    \node[note] at ([yshift=-3.5mm]current bounding box.south west)
      {\textbf{Both heads emit the same $192$-d $\hat z$ and retrieve exactly
       as in Figure~\ref{fig:topo-gru}.} The towers share nothing --- separate
       modules, separate parameters --- and their per-step hidden states are
       concatenated, so the combiner sees the towers' \emph{representations}
       rather than their scores and gradients flow back into both. There is no
       bidirectional mode by construction: a backward pass would let step $t$
       read item $t+1$ and copy its own teacher-forcing target. All three views
       are strictly causal for the same reason, and so is the pre-MLP, which is
       applied \emph{timestep-wise}; it is drawn solid because it is trained,
       red because it is refuted.};
  \end{tikzpicture}
  \caption{\texttt{seq-dualgru}: two recurrent towers with entirely separate
    weights read the prefix in parallel and a head fuses their hidden states
    --- the strongest available form of ``train the combination end-to-end so
    the parts co-adapt''. Each tower's input is a configurable causal view
    (\texttt{view\_a} / \texttt{view\_b}); the default is
    \texttt{latent}/\texttt{latent}, i.e.\ two independent GRUs over the same
    stream. Two of the boxes are coloured by the record rather than by the
    design. Removing the MLP fusion head (\texttt{fusion\_layers}${}=0$,
    green) \emph{beats} keeping it in three independent measurements, most
    recently by $+\num{0.02655}$ $[+0.01398, +0.03913]$ CI$>$0 --- and with
    some $9$--$10\%$ \emph{fewer} parameters, so the linear readout is the
    family default; and the
    per-step pre-MLP (red) loses the decisive single-tower ablation, Recall@10
    \num{0.10901} (\num{156} of \num{1431}) against the baseline's \num{176},
    paired $\Delta$ $-\num{0.01398}$ $[-0.02657, -0.00140]$ CI$<$0. Visual
    convention as in Figure~\ref{fig:topo-gru}. The best dual arm reaches
    \num{0.11880} and merely \emph{ties} the single-GRU bar, as does a
    $0.098\%$ parameter-matched width control --- so capacity is not the
    explanation.}
  \label{fig:topo-dualgru}
\end{figure}


% ===========================================================================
% 4. seq-embed --- a learned, weight-tied item table (DESIGN ONLY)
% ===========================================================================
\begin{figure}[ht]
  \centering
  \begin{tikzpicture}[topo]
    \node[frozen, minimum width=20mm, minimum height=11mm] (ids)
      {session prefix\\item \emph{ids}\\$i_1\dots i_t$};
    \node[trained, right=10mm of ids, minimum height=11mm] (E)
      {learned item\\table $E$\\[1pt]{\scriptsize$19{,}402\times192$}};
    \node[trained, right=6mm of E] (gru)
      {GRU $h{=}256$\\{\scriptsize$+$ dropout}};
    \node[trained, right=6mm of gru] (hd)
      {linear head\\{\scriptsize$256\to192$}};
    \node[op, right=6mm of hd] (rk)
      {$\cos(\hat z, E_v)$ over all\\$19{,}402$ rows $\to$ top-10};
    \draw[flow] (ids) -- node[edg, above, inner sep=1pt] {gather} (E);
    \draw[flow] (E) -- (gru);
    \draw[flow] (gru) -- (hd);
    \draw[flow] (hd) -- node[edg, above, inner sep=1pt] {$\hat z$} (rk);
    % ---- the three modes -------------------------------------------------
    \node[op, text width=64mm, align=flush left, anchor=south west] (md)
      at ([xshift=4mm,yshift=8mm]E.north east)
      {\texttt{embed\_mode} --- what the table $E$ actually is:\\
       \textbf{\texttt{frozen}} $E=Z$, nothing trained (the reproduction gate)\\
       \textbf{\texttt{residual}} $E=Z+\Delta E$: frozen content vector $+$ a
       trained correction, so a cold item still has a position\\
       \textbf{\texttt{free}} $E$ fully trained, optionally randomly initialised};
    \draw[flow] (md.west) -- (E.north east);
    % ---- weight tying: the same table is the retrieval target -----------
    \coordinate (c1) at ([yshift=-8mm]E.south);
    \coordinate (c2) at (rk.south |- c1);
    \draw[tie] (E.south) -- (c1);
    \draw[tie] (c1) -- node[tag, above, text=topogood]
      {weight tying --- the \emph{same} table is\\the input embedding
       \emph{and} the retrieval target} (c2);
    \draw[tie] (c2) -- (rk.south);
    % ---- the gradient loop that makes the geometry trainable ------------
    \coordinate (t1) at ([yshift=5mm]rk.north |- current bounding box.north);
    \coordinate (t2) at (E.north |- t1);
    \draw[grad] (rk.north) -- (t1);
    \draw[grad] (t1) -- node[tag, above, text=topobad]
      {$\nabla$ --- one InfoNCE step moves the query \emph{and} every target:
       the geometry itself is trainable} (t2);
    \draw[grad] (t2) -- (E.north);
    \node[note] at ([yshift=-3.5mm]current bounding box.south west)
      {\textbf{Why the loop is the whole point.} In
       Figures~\ref{fig:topo-gru}--\ref{fig:topo-dualgru} the retrieval space
       is fixed and only the map into it is learned; here the space is a
       parameter, shaped by next-track adjacency rather than inherited from a
       reconstruction objective. That also creates a collapse hazard absent
       everywhere else: with a learned table a plain cosine objective has a
       trivial optimum at ``send every item to one point'', so only a
       contrastive loss is admissible and \texttt{loss=cosine} is rejected
       outright for the trained modes.};
  \end{tikzpicture}
  \caption{\texttt{seq-embed}, the one axis the record has not tested:
    \textbf{a design, not a result}. No run of this predictor exists on any
    dataset, it is absent from the server's registered predictor list, and
    nothing here is a measured score. The item table $E$ is used twice, as the
    GRU's input embedding and as the retrieval target (weight tying, the
    GRU4Rec/SASRec formulation), so a single gradient step moves the query and
    all the candidates together --- the return arrow along the top is what
    makes the geometry trainable, and it is exactly what every other topology
    in this document lacks. The three \texttt{embed\_mode} settings span the
    space from a reproduction gate to a fully free table. What \emph{is}
    measured is the data budget that bounds them: \num{73632} training steps
    to fit \num{15816} item vectors, a median per-item training frequency of
    \num{1}, only \num{3294} items occurring five or more times, and
    \num{788} of \num{1431} test targets never appearing in any training
    session --- so a \texttt{free} table cannot retrieve the cold
    \SI{55}{\percent} at all and is capped at Recall@10 \num{0.4493}, whereas
    \texttt{residual} keeps the content vector and with it the cold reach.
    Visual convention as in Figure~\ref{fig:topo-gru}.}
  \label{fig:topo-embed}
\end{figure}


% ===========================================================================
% 5. seq-blend --- the champion
% ===========================================================================
\begin{figure}[ht]
  \centering
  \begin{tikzpicture}[topo]
    \node[frozen, minimum width=18mm, minimum height=11mm] (pf)
      {session\\prefix\\$i_1\dots i_t$};
    % ---- the learned metric map: shared by two legs, bypassed by one -----
    \node[trained, right=4mm of pf] (W)
      {learned map $W$\\$192\times192$\\[1pt]{\scriptsize no bias, identity init}\\{\scriptsize
       InfoNCE $\tau{=}0.07$}\\{\scriptsize $36{,}864$ params}};
    % ---- the three legs -------------------------------------------------
    \node[learnspace, right=3mm of W] (C)
      {\Cp\ content-$k$NN in $W$-space\\$\max_{u}\cos(Wz_v,Wz_u)$\\[1pt]{\scriptsize
       no parameters of its own}};
    \node[trained, above=7mm of C] (R)
      {\Rp\ projected GRU\\$h{=}256$, $394{,}944$ params};
    \node[lookup, below=7mm of C] (M)
      {$M$ first-order Markov counts\\$P(v\mid u)$, $u=$ last prefix item\\[1pt]{\scriptsize
       Dirichlet trust $n_u/(n_u{+}8)$,}\\{\scriptsize artist/genre back-off; $0$ params}};
    \draw[flow] (pf.east) -- (W.west);
    % ---- the frozen item table W acts ON: the same Z as Figure~\ref{fig:topo-gru} --
    \node[frozen, anchor=south east, minimum width=18mm] (Z)
      at ([yshift=6mm]W.north west)
      {frozen item latents $Z$\\{\scriptsize$19{,}402\times192$}};
    \draw[data] (Z.south east) -- node[edg, left, inner sep=1.2pt] {$z_i$} (W.north west);
    \draw[proj] (W.north east) -- (R.west);
    \draw[proj] (W.east) -- (C.west);
    \draw[data] (pf.south) |- node[tag, above, pos=0.55]
      {item \emph{ids} only ---\\the $M$ leg bypasses $Z$ and $W$} (M.west);
    % ---- the fixed combiner ---------------------------------------------
    \node[op, right=3mm of C] (cb)
      {z-normalise each leg\\over the candidate set,\\then average\\[1pt]{\scriptsize
       equal thirds ---}\\{\scriptsize\emph{no learned weights}}};
    \draw[flow] (R.east) -- node[edg, above, sloped, pos=0.55] {$s_{\Rp}$} (cb.north west);
    \draw[flow] (C.east) -- node[edg, above, inner sep=1pt] {$s_{\Cp}$} (cb.west);
    \draw[flow] (M.east) -- node[edg, above, sloped, pos=0.30] {$s_M$} (cb.south west);
    \node[op, right=2mm of cb] (rk) {rank\\$\to$ top-10};
    \draw[flow] (cb) -- (rk);
    \node[note] at ([yshift=-3.5mm]current bounding box.south west)
      {\textbf{Which legs are learned.} $W$ and the \Rp\ recurrence are
       trained (\num{431808} parameters between them); \Cp\ has no parameters
       of its own but retrieves inside $W$'s learned geometry, which is why it
       is drawn dashed-blue; $M$ is a table of training-set transition counts,
       drawn dashed-purple. The combiner is a fixed rule: each leg's
       full-vocabulary score vector is z-scored over the admissible candidates
       (every item minus those already in the prefix) and the three are
       averaged with hard-coded equal weights. The z-scoring is the
       load-bearing step --- it puts three incomparable score scales onto one
       per-session scale while \emph{preserving} relative magnitude, so a
       confidently-peaked leg can outvote two diffuse ones. Rank-only fusion,
       which discards that magnitude, loses by $-\num{0.029}$ CI$<$0.};
  \end{tikzpicture}
  \caption{\texttt{seq-blend} with \texttt{projection=true}, the champion:
    three full-vocabulary score vectors, z-normalised over the candidate set
    and averaged in fixed equal thirds. $W$ acts on the \emph{frozen item
    latents} $Z$ of Figure~\ref{fig:topo-gru} --- it warps the retrieval geometry,
    it does not read the prefix ids directly --- and it is shared by two of the
    three legs and bypassed by the third; that is the whole mechanism. $W$ is fit by in-batch InfoNCE on consecutive within-train
    pairs, which rebuilds the retrieval metric from ``reconstruct the item
    matrix'' into ``what follows what''; both legs that retrieve by cosine in
    that space improve --- the parameter-free content-$k$NN leg more than
    doubles, \num{0.0685} to \num{0.1516}, and the recurrent leg rises from
    \num{0.1230} to $\approx$\num{0.172} --- while the item-ID lookup is
    untouched at \num{0.106918}. The two projected-leg standalone figures come
    from the off-server driver of the 2026-07-18 campaign and have no server
    run of their own; only the $M$ leg's number is a canonical-split run.
    Together the three legs reach Recall@10 \num{0.21174} (\num{303} of
    \num{1431}), paired $\Delta$ over the unprojected three-leg blend
    $+\num{0.0259}$ $[+0.012, +0.041]$, a CI entirely above zero and robust
    across three seeds. Visual convention as in Figure~\ref{fig:topo-gru}.}
  \label{fig:topo-blend}
\end{figure}


% ===========================================================================
% 6. seq-stacker --- the learned combiner
% ===========================================================================
\begin{figure}[ht]
  \centering
  \begin{tikzpicture}[topo]
    \node[frozen, minimum width=18mm, minimum height=11mm] (pf)
      {session\\prefix\\$i_1\dots i_t$};
    \node[trained, right=6mm of pf] (R)
      {$R$ GRU $h{=}256$\\[1pt]{\scriptsize trained separately}};
    \node[lookup, above=6mm of R] (M) {$M$ Markov counts};
    \node[trained, below=6mm of R] (A)
      {$A$ pooled MLP\\[1pt]{\scriptsize trained separately}};
    \draw[flow] (pf.east) -- (M.west);
    \draw[flow] (pf.east) -- (R.west);
    \draw[flow] (pf.east) -- (A.west);
    % ---- the per-candidate feature row ----------------------------------
    \node[op, right=6mm of R] (fx)
      {per-candidate feature row\\$[\,s_M,\ s_R,\ s_A,\ \log(1{+}\text{plays}),$\\$\text{recency},\
       \text{artist match},\ \text{genre match}\,]$};
    \draw[flow] (M.east) -- (fx.north west);
    \draw[flow] (R.east) -- (fx.west);
    \draw[flow] (A.east) -- (fx.south west);
    \node[refuted, right=5mm of fx] (xgb)
      {XGBoost meta-learner\\$200$ rounds, depth $6$\\[1pt]{\scriptsize\texttt{binary:logistic}}};
    \draw[flow] (fx) -- (xgb);
    \node[op, below=9mm of xgb] (rk)
      {score all $19{,}402$ items\\$\to$ top-10};
    \draw[flow] (xgb) -- (rk);
    % ---- the two failure mechanisms, marked on the diagram --------------
    \node[refuted, text width=56mm, align=flush left, anchor=north west] (tr)
      at ([yshift=-10mm]fx.south west)
      {\textbf{Training rows}: $1$ positive $+$ $50$ popularity-sampled
       negatives per train session, i.e.\ \num{5723} positives, against the
       \num{73632} teacher-forced steps its own $R$ leg trains on
       ($12.9\times$ less signal) --- and a decision surface calibrated on
       $51$ candidates ($0.26\%$ of the catalogue) is then asked to rank all
       \num{19402}.};
    \draw[grad] (tr.north east) -- (xgb.south west);
    \node[note] at ([yshift=-3.5mm]current bounding box.south west)
      {\textbf{Why the shape matters.} A $200$-round depth-$6$ booster is
       piecewise \emph{constant}, so ranking \num{19402} candidates with it
       produces large tied blocks that the harness breaks by item index. The
       fixed z-blend of Figure~\ref{fig:topo-blend} is instead a strictly
       monotone continuous function of the leg scores: it preserves each leg's
       internal ordering exactly and can only trade the legs off against one
       another.};
  \end{tikzpicture}
  \caption{\texttt{seq-stacker}, the learned counterpart of
    Figure~\ref{fig:topo-blend}'s fixed rule: the same separately-trained base
    legs, but their per-candidate scores become columns of a feature row
    alongside four cheap candidate features, and a gradient-boosted tree
    learns the combination. It earns a diagram because its failure is
    informative --- with the \texttt{markov}, \texttt{gru} and \texttt{ann}
    legs it reaches Recall@10 \num{0.085255} (\num{122} of \num{1431}),
    \emph{below its own GRU base leg} at \num{0.12299}, paired $\Delta$
    $-\num{0.0377}$ $[-0.056, -0.020]$ CI$<$0. The red callout names the two
    mechanisms: supervision starvation, since the meta-learner gets one
    labelled event per session where its base leg gets one per prefix
    position, and a train/serve candidate mismatch. The booster is drawn solid
    because it is trained and red because it is refuted; the feature row and
    the ranking step are dashed because they hold no parameters. Visual
    convention as in Figure~\ref{fig:topo-gru}. This is one of four
    independent learned-combiner nulls, the others being reciprocal-rank
    fusion, oracle weights fitted on the test labels, and the jointly-trained
    dual tower of Figure~\ref{fig:topo-dualgru}.}
  \label{fig:topo-stacker}
\end{figure}
, can use the same
%  visual language.  \input from docs/ (compilation cwd):
%
%      % ===========================================================================
%  topo-style.tex --- the shared TikZ style block for the next-track topology
%  diagrams AND for the evaluation-protocol diagram in the body of
%  model-comparison.tex.  Extracted from topologies.tex so that the protocol
%  figure, which appears BEFORE \input{figures/topologies}, can use the same
%  visual language.  \input from docs/ (compilation cwd):
%
%      \input{figures/topo-style}
%
%  Idempotent: guarded, so a second \input is a no-op.  Requires
%      \usepackage{tikz}
%      \usetikzlibrary{positioning,arrows.meta,fit,backgrounds,calc}
% ===========================================================================
\ifdefined\nexttracktopostyle\else
\newcommand{\nexttracktopostyle}{}

\providecommand{\dsid}[1]{\texttt{\small #1}}
\providecommand{\Rp}{\ensuremath{R'}}
\providecommand{\Cp}{\ensuremath{C'}}

% --- palette: the same hex values as docs/figures/make_figures.py ----------
\definecolor{topoacc}{HTML}{2563EB}   % ACCENT  --- default / on-record path
\definecolor{topogood}{HTML}{16A34A}  % GOOD    --- favoured variant
\definecolor{topobad}{HTML}{DC2626}   % BAD     --- refuted variant
\definecolor{topomut}{HTML}{9CA3AF}   % MUTED   --- frozen, parameter-free
\definecolor{topolook}{HTML}{9333EA}  % seq-markov purple --- table lookup
\definecolor{topoink}{HTML}{374151}   % annotation ink

\tikzset{
  topo/.style={
    font=\footnotesize,
    every node/.append style={align=center},
  },
  box/.style={rounded corners=1.4pt, inner xsep=3.2pt, inner ysep=2.8pt,
              minimum height=6.4mm, text=topoink},
  dsh/.style={dash pattern=on 2pt off 1.4pt},
  % --- component styles ---------------------------------------------------
  trained/.style={box, draw=topoacc, line width=0.9pt, fill=topoacc!7},
  frozen/.style={box, dsh, draw=topomut, line width=0.8pt, fill=topomut!14},
  lookup/.style={box, dsh, draw=topolook, line width=0.9pt, fill=topolook!7},
  learnspace/.style={box, dsh, draw=topoacc, line width=0.9pt, fill=topoacc!5},
  won/.style={box, draw=topogood, line width=1.0pt, fill=topogood!8},
  refuted/.style={box, draw=topobad, line width=1.0pt, fill=topobad!6},
  op/.style={box, dash pattern=on 1.5pt off 1.2pt, draw=topomut,
             line width=0.7pt, fill=white, rounded corners=3.2pt},
  % --- edges --------------------------------------------------------------
  flow/.style={-{Stealth[length=4.2pt,width=3.4pt]}, draw=topoink,
               line width=0.7pt},
  proj/.style={flow, draw=topoacc},
  data/.style={flow, draw=topomut, dsh},
  tie/.style={flow, draw=topogood, line width=0.8pt},
  grad/.style={flow, draw=topobad, line width=0.75pt,
               dash pattern=on 1.6pt off 1.3pt},
  variant/.style={flow, draw=topobad, line width=0.75pt,
                  dash pattern=on 1.6pt off 1.3pt},
  % --- text ---------------------------------------------------------------
  tag/.style={font=\scriptsize, text=topoink, inner sep=1.2pt},
  edg/.style={font=\scriptsize, text=topoink, inner sep=1.4pt},
  note/.style={font=\scriptsize, text=topoink, align=flush left,
               text width=134mm, inner sep=1.4pt, anchor=north west},
  grp/.style={draw=topomut!80, line width=0.5pt, rounded corners=2.4pt,
              dash pattern=on 1pt off 1pt, inner sep=3.4pt},
  swatch/.style={minimum width=5.2mm, minimum height=3.2mm, inner sep=0pt},
}
\fi

%
%  Idempotent: guarded, so a second \input is a no-op.  Requires
%      \usepackage{tikz}
%      \usetikzlibrary{positioning,arrows.meta,fit,backgrounds,calc}
% ===========================================================================
\ifdefined\nexttracktopostyle\else
\newcommand{\nexttracktopostyle}{}

\providecommand{\dsid}[1]{\texttt{\small #1}}
\providecommand{\Rp}{\ensuremath{R'}}
\providecommand{\Cp}{\ensuremath{C'}}

% --- palette: the same hex values as docs/figures/make_figures.py ----------
\definecolor{topoacc}{HTML}{2563EB}   % ACCENT  --- default / on-record path
\definecolor{topogood}{HTML}{16A34A}  % GOOD    --- favoured variant
\definecolor{topobad}{HTML}{DC2626}   % BAD     --- refuted variant
\definecolor{topomut}{HTML}{9CA3AF}   % MUTED   --- frozen, parameter-free
\definecolor{topolook}{HTML}{9333EA}  % seq-markov purple --- table lookup
\definecolor{topoink}{HTML}{374151}   % annotation ink

\tikzset{
  topo/.style={
    font=\footnotesize,
    every node/.append style={align=center},
  },
  box/.style={rounded corners=1.4pt, inner xsep=3.2pt, inner ysep=2.8pt,
              minimum height=6.4mm, text=topoink},
  dsh/.style={dash pattern=on 2pt off 1.4pt},
  % --- component styles ---------------------------------------------------
  trained/.style={box, draw=topoacc, line width=0.9pt, fill=topoacc!7},
  frozen/.style={box, dsh, draw=topomut, line width=0.8pt, fill=topomut!14},
  lookup/.style={box, dsh, draw=topolook, line width=0.9pt, fill=topolook!7},
  learnspace/.style={box, dsh, draw=topoacc, line width=0.9pt, fill=topoacc!5},
  won/.style={box, draw=topogood, line width=1.0pt, fill=topogood!8},
  refuted/.style={box, draw=topobad, line width=1.0pt, fill=topobad!6},
  op/.style={box, dash pattern=on 1.5pt off 1.2pt, draw=topomut,
             line width=0.7pt, fill=white, rounded corners=3.2pt},
  % --- edges --------------------------------------------------------------
  flow/.style={-{Stealth[length=4.2pt,width=3.4pt]}, draw=topoink,
               line width=0.7pt},
  proj/.style={flow, draw=topoacc},
  data/.style={flow, draw=topomut, dsh},
  tie/.style={flow, draw=topogood, line width=0.8pt},
  grad/.style={flow, draw=topobad, line width=0.75pt,
               dash pattern=on 1.6pt off 1.3pt},
  variant/.style={flow, draw=topobad, line width=0.75pt,
                  dash pattern=on 1.6pt off 1.3pt},
  % --- text ---------------------------------------------------------------
  tag/.style={font=\scriptsize, text=topoink, inner sep=1.2pt},
  edg/.style={font=\scriptsize, text=topoink, inner sep=1.4pt},
  note/.style={font=\scriptsize, text=topoink, align=flush left,
               text width=134mm, inner sep=1.4pt, anchor=north west},
  grp/.style={draw=topomut!80, line width=0.5pt, rounded corners=2.4pt,
              dash pattern=on 1pt off 1pt, inner sep=3.4pt},
  swatch/.style={minimum width=5.2mm, minimum height=3.2mm, inner sep=0pt},
}
\fi




% ===========================================================================
% 1. seq-nexttrack --- the baseline (and the legend)
% ===========================================================================
\begin{figure}[ht]
  \centering
  \begin{tikzpicture}[topo]
    % ---- the frozen item-latent table, read TWICE ------------------------
    \node[frozen, minimum width=20mm, minimum height=12mm] (tab)
      {frozen item\\latents $Z$\\[1pt]{\scriptsize$19{,}402\times192$}};
    % ---- the prefix, in order (one column per step) ----------------------
    \node[frozen, minimum width=12mm, right=7mm of tab]  (x1) {$x_1$};
    \node[frozen, minimum width=12mm, right=4mm of x1]   (x2) {$x_2$};
    \node[tag, right=2.5mm of x2]                        (xd) {$\cdots$};
    \node[frozen, minimum width=12mm, right=2.5mm of xd] (x3) {$x_t$};
    \begin{scope}[on background layer]
      \node[grp, fit=(x1)(x3)] (pfx) {};
    \end{scope}
    \draw[data] (tab.east) -- (pfx.west);
    \node[tag, below=1.2mm of pfx] (pfxlab)
      {prefix latents --- \emph{order preserved}};
    % ---- the recurrence --------------------------------------------------
    \node[trained, minimum width=12mm, above=7mm of x1] (g1) {GRU};
    \node[trained, minimum width=12mm, above=7mm of x2] (g2) {GRU};
    \node[tag] (gd) at (g1 -| xd) {$\cdots$};
    \node[trained, minimum width=12mm, above=7mm of x3] (g3) {GRU};
    \draw[flow] (x1) -- (g1);
    \draw[flow] (x2) -- (g2);
    \draw[flow] (x3) -- (g3);
    \draw[flow] (g1) -- node[edg, above, inner sep=0.7pt] {$h_t$} (g2);
    \draw[flow] (g2) -- (gd);
    \draw[flow] (gd) -- (g3);
    % ---- head and retrieval ---------------------------------------------
    \node[trained, right=6mm of g3] (head)
      {linear head\\{\scriptsize$256\to192$}};
    \node[op, right=7mm of head] (rk)
      {rank the catalogue by $\cos(\hat z, z_v)$\\{\scriptsize prefix excluded $\to$ top-10}};
    \draw[flow] (g3) -- (head);
    \draw[flow] (head) -- node[edg, above, inner sep=1pt] {$\hat z$} (rk);
    % ---- the SAME frozen table is the retrieval target -------------------
    \coordinate (c1) at ([yshift=-3mm]tab.south |- current bounding box.south);
    \coordinate (c2) at (rk.south |- c1);
    \draw[data] (tab.south) -- (c1);
    \draw[data] (c1) -- node[tag, above]
      {the \emph{same} frozen table is the retrieval target} (c2);
    \draw[data] (c2) -- (rk.south);
    % ---- supervision note ------------------------------------------------
    \node[note] (nt) at ([yshift=-3.5mm]current bounding box.south west)
      {\textbf{Supervision --- per-step teacher forcing.} The head is shared
       across positions and must predict the next item's latent at
       \emph{every} prefix position, so one session of length $L$ yields
       $L-1$ examples: $73{,}632$ (prefix, next-item) pairs on the train
       split. The loss is masked in-batch InfoNCE at $\tau=0.07$ --- each
       step's true next latent must outscore the other steps' targets in the
       same minibatch, which supplies negatives for free.};
    % ---- LEGEND (two rows so it cannot widen the picture) ----------------
    \node[trained, swatch, anchor=north west] (l1)
      at ([yshift=-3mm]current bounding box.south west) {};
    \node[tag, right=1.1mm of l1] (t1) {trained (solid)};
    \node[frozen, swatch, right=3.4mm of t1] (l2) {};
    \node[tag, right=1.1mm of l2] (t2) {frozen / parameter-free (dashed)};
    \node[lookup, swatch, right=3.4mm of t2] (l3) {};
    \node[tag, right=1.1mm of l3] (t3) {table lookup};
    \node[learnspace, swatch, anchor=north west] (l4)
      at ([yshift=-1.6mm]l1.south west) {};
    \node[tag, right=1.1mm of l4] (t4) {parameter-free, learned space};
    \node[won, swatch, right=3.4mm of t4] (l5) {};
    \node[tag, right=1.1mm of l5] (t5) {favoured variant};
    \node[refuted, swatch, right=3.4mm of t5] (l6) {};
    \node[tag, right=1.1mm of l6] (t6) {refuted variant};
  \end{tikzpicture}
  \caption{\texttt{seq-nexttrack}, the single-model baseline every other
    topology in this document is measured against: a frozen item space, a
    recurrence over the \emph{ordered} prefix, a linear head, and cosine
    retrieval back into the \emph{same} frozen space. Nothing about the item
    geometry is learned here --- the $19{,}402\times192$ latent table (offline
    PCA-192 on \dsid{seq-20260715-131139}) is read twice, once to embed the
    prefix and once as the retrieval target, and only the GRU and its linear
    head carry parameters (\num{394944} at $h=256$). \textbf{The visual
    language used throughout this document is fixed here.} A \emph{solid}
    border marks a trained component and a \emph{dashed} border a frozen or
    parameter-free one; that distinction rides on the line style, not the
    colour, so it survives greyscale. Colour then carries the verdict ---
    blue: the on-record default path; purple dashed: a non-parametric table
    lookup over the dataset; dashed blue: parameter-free but operating inside
    a \emph{learned} geometry; green: a variant the record favours; red: a
    variant the record refutes. On the canonical split this configuration
    scores Recall@10 \num{0.12299} (\num{176} of \num{1431}), reproduced
    bit-identically eleven times.}
  \label{fig:topo-gru}
\end{figure}


% ===========================================================================
% 2. seq-ann --- the pooled feed-forward twin
% ===========================================================================
\begin{figure}[ht]
  \centering
  \begin{tikzpicture}[topo]
    \node[frozen, minimum width=20mm, minimum height=12mm] (tab)
      {frozen item\\latents $Z$\\[1pt]{\scriptsize$19{,}402\times192$}};
    \node[frozen, minimum width=12mm, right=7mm of tab]  (x1) {$x_1$};
    \node[frozen, minimum width=12mm, right=4mm of x1]   (x2) {$x_2$};
    \node[tag, right=2.5mm of x2]                        (xd) {$\cdots$};
    \node[frozen, minimum width=12mm, right=2.5mm of xd] (x3) {$x_t$};
    \begin{scope}[on background layer]
      \node[grp, fit=(x1)(x3)] (pfx) {};
    \end{scope}
    \draw[data] (tab.east) -- (pfx.west);
    % ---- the two parameter-free summaries of the same prefix -------------
    \node[op, right=8mm of pfx, yshift=10mm] (mp)
      {causal mean-pool\\$\frac{1}{t}\sum_{i\le t}x_i$\\[1pt]{\scriptsize\emph{permutation-invariant}}};
    \node[op, right=8mm of pfx, yshift=-10mm] (lt)
      {last item\\$x_t$\\[1pt]{\scriptsize the only recency cue}};
    \draw[flow] (pfx.east) -- (mp.west);
    \draw[flow] (pfx.east) -- (lt.west);
    % ---- the MLP ---------------------------------------------------------
    \node[trained, right=5mm of mp, yshift=-10mm] (mlp)
      {MLP\\$384\to256\to256\to192$\\[1pt]{\scriptsize ReLU $+$ dropout}};
    \draw[flow] (mp.east) -- (mlp.north west);
    \draw[flow] (lt.east) -- (mlp.south west);
    \node[edg] at ([xshift=-4.2mm]mlp.west) {$\|$};
    % ---- retrieval, placed below to keep the picture inside the text ------
    \node[op, below=9mm of mlp] (rk)
      {$\hat z\to$ cosine rank\\{\scriptsize prefix excluded $\to$ top-10}};
    \draw[flow] (mlp) -- (rk);
    \coordinate (c1) at ([yshift=-3mm]tab.south |- current bounding box.south);
    \coordinate (c2) at (rk.south |- c1);
    \draw[data] (tab.south) -- (c1);
    \draw[data] (c1) -- node[tag, above]
      {the same frozen table is the retrieval target} (c2);
    \draw[data] (c2) -- (rk.south);
    \node[note] at ([yshift=-3.5mm]current bounding box.south west)
      {\textbf{The same information, with the order discarded.} Both summaries
       are strictly causal, so the objective, the padded batching and the
       early stopping are reused \emph{verbatim} from
       Figure~\ref{fig:topo-gru} --- but a mean weights the track from twenty
       steps ago exactly like the one before last, and cannot represent ``A
       then B'' differently from ``B then A''. Only the special-cased $x_t$ tap
       carries any notion of recency, and the concatenation ($\|$) is this
       model's entire account of sequence.};
  \end{tikzpicture}
  \caption{\texttt{seq-ann}, the GRU's non-recurrent twin: the recurrence of
    Figure~\ref{fig:topo-gru} is replaced by two parameter-free summaries of
    the same prefix --- a causal mean-pool concatenated with the last item's
    latent --- feeding a plain three-layer ReLU MLP. Everything else is
    literally the same code (\texttt{seq\_ann.py} imports
    \texttt{make\_batches} and \texttt{step\_loss} from
    \texttt{seq\_nexttrack}), which is what makes this pair a clean test of
    the recurrence itself rather than of a training recipe. Visual convention
    as in Figure~\ref{fig:topo-gru}: the two pooling operators are dashed
    because they hold no parameters, and only the MLP is trained. The
    defensible comparison is same-space: on \dsid{seq-20260715-030509}
    (PCA-128) the pooled twin scores Recall@10 \num{0.0901} (\num{129} of
    \num{1431}) against the GRU's \num{0.1146} (\num{164}), paired $\Delta$
    $-\num{0.0245}$ $[-0.039, -0.011]$, a CI \emph{entirely below zero}.
    \texttt{seq-ann} has no run at all on the canonical PCA-192 split, so it
    never belongs in the same column as the \num{0.12299} bar.}
  \label{fig:topo-ann}
\end{figure}


% ===========================================================================
% 3. seq-dualgru --- two towers, a fusion head, and the pre-MLP variant
% ===========================================================================
\begin{figure}[ht]
  \centering
  \begin{tikzpicture}[topo]
    \node[frozen, minimum width=19mm, minimum height=11mm] (pf)
      {prefix\\latents\\$x_1\dots x_t$};
    % ---- the three configurable causal views ----------------------------
    \node[op, right=7mm of pf] (v2) {\texttt{delta}\\$x_t-x_{t-1}$};
    \node[op, above=2.6mm of v2] (v1) {\texttt{latent}\\$x_t$};
    \node[op, below=2.6mm of v2] (v3) {\texttt{cummean}\\$\frac{1}{t}\sum_{i\le t}x_i$};
    \begin{scope}[on background layer]
      \node[grp, fit=(v1)(v3)] (vg) {};
    \end{scope}
    \draw[data] (pf.east) -- (vg.west);
    % ---- two independent towers -----------------------------------------
    \node[trained, right=8mm of vg, yshift=9mm] (ta)
      {tower A on \texttt{view\_a}\\GRU $h_a{=}256$\\[1pt]{\scriptsize own weights}};
    \node[trained, right=8mm of vg, yshift=-9mm] (tb)
      {tower B on \texttt{view\_b}\\GRU $h_b{=}256$\\[1pt]{\scriptsize own weights}};
    \draw[flow] (vg.east) -- (ta.west);
    \draw[flow] (vg.east) -- (tb.west);
    % ---- concat + the two fusion depths ---------------------------------
    \node[op, right=7mm of ta, yshift=-9mm] (cc)
      {concat\\$h_a\|h_b$\\[1pt]{\scriptsize$512$}};
    \draw[flow] (ta.east) -- (cc.north west);
    \draw[flow] (tb.east) -- (cc.south west);
    \node[won, right=7mm of cc, yshift=9mm] (f0)
      {\texttt{fusion\_layers}${}=0$\\$512\to192$, linear};
    \node[refuted, right=7mm of cc, yshift=-9mm] (f1)
      {\texttt{fusion\_layers}${}=1$\\$512\to256\to192$\\[1pt]{\scriptsize$+$ ReLU $+$ dropout}};
    \draw[flow] (cc.north east) -- (f0.west);
    \draw[flow] (cc.south east) -- (f1.west);
    % ---- the refuted per-step pre-MLP variant ---------------------------
    \node[refuted, below=12mm of tb] (pb)
      {\emph{variant} --- per-step pre-MLP\\on a tower:
       Linear$(192{\to}256)$\\$+$ ReLU $+$ Dropout};
    \draw[variant] (vg.south) |- (pb.west);
    \draw[variant] (pb.north) -- (tb.south);
    \node[note] at ([yshift=-3.5mm]current bounding box.south west)
      {\textbf{Both heads emit the same $192$-d $\hat z$ and retrieve exactly
       as in Figure~\ref{fig:topo-gru}.} The towers share nothing --- separate
       modules, separate parameters --- and their per-step hidden states are
       concatenated, so the combiner sees the towers' \emph{representations}
       rather than their scores and gradients flow back into both. There is no
       bidirectional mode by construction: a backward pass would let step $t$
       read item $t+1$ and copy its own teacher-forcing target. All three views
       are strictly causal for the same reason, and so is the pre-MLP, which is
       applied \emph{timestep-wise}; it is drawn solid because it is trained,
       red because it is refuted.};
  \end{tikzpicture}
  \caption{\texttt{seq-dualgru}: two recurrent towers with entirely separate
    weights read the prefix in parallel and a head fuses their hidden states
    --- the strongest available form of ``train the combination end-to-end so
    the parts co-adapt''. Each tower's input is a configurable causal view
    (\texttt{view\_a} / \texttt{view\_b}); the default is
    \texttt{latent}/\texttt{latent}, i.e.\ two independent GRUs over the same
    stream. Two of the boxes are coloured by the record rather than by the
    design. Removing the MLP fusion head (\texttt{fusion\_layers}${}=0$,
    green) \emph{beats} keeping it in three independent measurements, most
    recently by $+\num{0.02655}$ $[+0.01398, +0.03913]$ CI$>$0 --- and with
    some $9$--$10\%$ \emph{fewer} parameters, so the linear readout is the
    family default; and the
    per-step pre-MLP (red) loses the decisive single-tower ablation, Recall@10
    \num{0.10901} (\num{156} of \num{1431}) against the baseline's \num{176},
    paired $\Delta$ $-\num{0.01398}$ $[-0.02657, -0.00140]$ CI$<$0. Visual
    convention as in Figure~\ref{fig:topo-gru}. The best dual arm reaches
    \num{0.11880} and merely \emph{ties} the single-GRU bar, as does a
    $0.098\%$ parameter-matched width control --- so capacity is not the
    explanation.}
  \label{fig:topo-dualgru}
\end{figure}


% ===========================================================================
% 4. seq-embed --- a learned, weight-tied item table (DESIGN ONLY)
% ===========================================================================
\begin{figure}[ht]
  \centering
  \begin{tikzpicture}[topo]
    \node[frozen, minimum width=20mm, minimum height=11mm] (ids)
      {session prefix\\item \emph{ids}\\$i_1\dots i_t$};
    \node[trained, right=10mm of ids, minimum height=11mm] (E)
      {learned item\\table $E$\\[1pt]{\scriptsize$19{,}402\times192$}};
    \node[trained, right=6mm of E] (gru)
      {GRU $h{=}256$\\{\scriptsize$+$ dropout}};
    \node[trained, right=6mm of gru] (hd)
      {linear head\\{\scriptsize$256\to192$}};
    \node[op, right=6mm of hd] (rk)
      {$\cos(\hat z, E_v)$ over all\\$19{,}402$ rows $\to$ top-10};
    \draw[flow] (ids) -- node[edg, above, inner sep=1pt] {gather} (E);
    \draw[flow] (E) -- (gru);
    \draw[flow] (gru) -- (hd);
    \draw[flow] (hd) -- node[edg, above, inner sep=1pt] {$\hat z$} (rk);
    % ---- the three modes -------------------------------------------------
    \node[op, text width=64mm, align=flush left, anchor=south west] (md)
      at ([xshift=4mm,yshift=8mm]E.north east)
      {\texttt{embed\_mode} --- what the table $E$ actually is:\\
       \textbf{\texttt{frozen}} $E=Z$, nothing trained (the reproduction gate)\\
       \textbf{\texttt{residual}} $E=Z+\Delta E$: frozen content vector $+$ a
       trained correction, so a cold item still has a position\\
       \textbf{\texttt{free}} $E$ fully trained, optionally randomly initialised};
    \draw[flow] (md.west) -- (E.north east);
    % ---- weight tying: the same table is the retrieval target -----------
    \coordinate (c1) at ([yshift=-8mm]E.south);
    \coordinate (c2) at (rk.south |- c1);
    \draw[tie] (E.south) -- (c1);
    \draw[tie] (c1) -- node[tag, above, text=topogood]
      {weight tying --- the \emph{same} table is\\the input embedding
       \emph{and} the retrieval target} (c2);
    \draw[tie] (c2) -- (rk.south);
    % ---- the gradient loop that makes the geometry trainable ------------
    \coordinate (t1) at ([yshift=5mm]rk.north |- current bounding box.north);
    \coordinate (t2) at (E.north |- t1);
    \draw[grad] (rk.north) -- (t1);
    \draw[grad] (t1) -- node[tag, above, text=topobad]
      {$\nabla$ --- one InfoNCE step moves the query \emph{and} every target:
       the geometry itself is trainable} (t2);
    \draw[grad] (t2) -- (E.north);
    \node[note] at ([yshift=-3.5mm]current bounding box.south west)
      {\textbf{Why the loop is the whole point.} In
       Figures~\ref{fig:topo-gru}--\ref{fig:topo-dualgru} the retrieval space
       is fixed and only the map into it is learned; here the space is a
       parameter, shaped by next-track adjacency rather than inherited from a
       reconstruction objective. That also creates a collapse hazard absent
       everywhere else: with a learned table a plain cosine objective has a
       trivial optimum at ``send every item to one point'', so only a
       contrastive loss is admissible and \texttt{loss=cosine} is rejected
       outright for the trained modes.};
  \end{tikzpicture}
  \caption{\texttt{seq-embed}, the one axis the record has not tested:
    \textbf{a design, not a result}. No run of this predictor exists on any
    dataset, it is absent from the server's registered predictor list, and
    nothing here is a measured score. The item table $E$ is used twice, as the
    GRU's input embedding and as the retrieval target (weight tying, the
    GRU4Rec/SASRec formulation), so a single gradient step moves the query and
    all the candidates together --- the return arrow along the top is what
    makes the geometry trainable, and it is exactly what every other topology
    in this document lacks. The three \texttt{embed\_mode} settings span the
    space from a reproduction gate to a fully free table. What \emph{is}
    measured is the data budget that bounds them: \num{73632} training steps
    to fit \num{15816} item vectors, a median per-item training frequency of
    \num{1}, only \num{3294} items occurring five or more times, and
    \num{788} of \num{1431} test targets never appearing in any training
    session --- so a \texttt{free} table cannot retrieve the cold
    \SI{55}{\percent} at all and is capped at Recall@10 \num{0.4493}, whereas
    \texttt{residual} keeps the content vector and with it the cold reach.
    Visual convention as in Figure~\ref{fig:topo-gru}.}
  \label{fig:topo-embed}
\end{figure}


% ===========================================================================
% 5. seq-blend --- the champion
% ===========================================================================
\begin{figure}[ht]
  \centering
  \begin{tikzpicture}[topo]
    \node[frozen, minimum width=18mm, minimum height=11mm] (pf)
      {session\\prefix\\$i_1\dots i_t$};
    % ---- the learned metric map: shared by two legs, bypassed by one -----
    \node[trained, right=4mm of pf] (W)
      {learned map $W$\\$192\times192$\\[1pt]{\scriptsize no bias, identity init}\\{\scriptsize
       InfoNCE $\tau{=}0.07$}\\{\scriptsize $36{,}864$ params}};
    % ---- the three legs -------------------------------------------------
    \node[learnspace, right=3mm of W] (C)
      {\Cp\ content-$k$NN in $W$-space\\$\max_{u}\cos(Wz_v,Wz_u)$\\[1pt]{\scriptsize
       no parameters of its own}};
    \node[trained, above=7mm of C] (R)
      {\Rp\ projected GRU\\$h{=}256$, $394{,}944$ params};
    \node[lookup, below=7mm of C] (M)
      {$M$ first-order Markov counts\\$P(v\mid u)$, $u=$ last prefix item\\[1pt]{\scriptsize
       Dirichlet trust $n_u/(n_u{+}8)$,}\\{\scriptsize artist/genre back-off; $0$ params}};
    \draw[flow] (pf.east) -- (W.west);
    % ---- the frozen item table W acts ON: the same Z as Figure~\ref{fig:topo-gru} --
    \node[frozen, anchor=south east, minimum width=18mm] (Z)
      at ([yshift=6mm]W.north west)
      {frozen item latents $Z$\\{\scriptsize$19{,}402\times192$}};
    \draw[data] (Z.south east) -- node[edg, left, inner sep=1.2pt] {$z_i$} (W.north west);
    \draw[proj] (W.north east) -- (R.west);
    \draw[proj] (W.east) -- (C.west);
    \draw[data] (pf.south) |- node[tag, above, pos=0.55]
      {item \emph{ids} only ---\\the $M$ leg bypasses $Z$ and $W$} (M.west);
    % ---- the fixed combiner ---------------------------------------------
    \node[op, right=3mm of C] (cb)
      {z-normalise each leg\\over the candidate set,\\then average\\[1pt]{\scriptsize
       equal thirds ---}\\{\scriptsize\emph{no learned weights}}};
    \draw[flow] (R.east) -- node[edg, above, sloped, pos=0.55] {$s_{\Rp}$} (cb.north west);
    \draw[flow] (C.east) -- node[edg, above, inner sep=1pt] {$s_{\Cp}$} (cb.west);
    \draw[flow] (M.east) -- node[edg, above, sloped, pos=0.30] {$s_M$} (cb.south west);
    \node[op, right=2mm of cb] (rk) {rank\\$\to$ top-10};
    \draw[flow] (cb) -- (rk);
    \node[note] at ([yshift=-3.5mm]current bounding box.south west)
      {\textbf{Which legs are learned.} $W$ and the \Rp\ recurrence are
       trained (\num{431808} parameters between them); \Cp\ has no parameters
       of its own but retrieves inside $W$'s learned geometry, which is why it
       is drawn dashed-blue; $M$ is a table of training-set transition counts,
       drawn dashed-purple. The combiner is a fixed rule: each leg's
       full-vocabulary score vector is z-scored over the admissible candidates
       (every item minus those already in the prefix) and the three are
       averaged with hard-coded equal weights. The z-scoring is the
       load-bearing step --- it puts three incomparable score scales onto one
       per-session scale while \emph{preserving} relative magnitude, so a
       confidently-peaked leg can outvote two diffuse ones. Rank-only fusion,
       which discards that magnitude, loses by $-\num{0.029}$ CI$<$0.};
  \end{tikzpicture}
  \caption{\texttt{seq-blend} with \texttt{projection=true}, the champion:
    three full-vocabulary score vectors, z-normalised over the candidate set
    and averaged in fixed equal thirds. $W$ acts on the \emph{frozen item
    latents} $Z$ of Figure~\ref{fig:topo-gru} --- it warps the retrieval geometry,
    it does not read the prefix ids directly --- and it is shared by two of the
    three legs and bypassed by the third; that is the whole mechanism. $W$ is fit by in-batch InfoNCE on consecutive within-train
    pairs, which rebuilds the retrieval metric from ``reconstruct the item
    matrix'' into ``what follows what''; both legs that retrieve by cosine in
    that space improve --- the parameter-free content-$k$NN leg more than
    doubles, \num{0.0685} to \num{0.1516}, and the recurrent leg rises from
    \num{0.1230} to $\approx$\num{0.172} --- while the item-ID lookup is
    untouched at \num{0.106918}. The two projected-leg standalone figures come
    from the off-server driver of the 2026-07-18 campaign and have no server
    run of their own; only the $M$ leg's number is a canonical-split run.
    Together the three legs reach Recall@10 \num{0.21174} (\num{303} of
    \num{1431}), paired $\Delta$ over the unprojected three-leg blend
    $+\num{0.0259}$ $[+0.012, +0.041]$, a CI entirely above zero and robust
    across three seeds. Visual convention as in Figure~\ref{fig:topo-gru}.}
  \label{fig:topo-blend}
\end{figure}


% ===========================================================================
% 6. seq-stacker --- the learned combiner
% ===========================================================================
\begin{figure}[ht]
  \centering
  \begin{tikzpicture}[topo]
    \node[frozen, minimum width=18mm, minimum height=11mm] (pf)
      {session\\prefix\\$i_1\dots i_t$};
    \node[trained, right=6mm of pf] (R)
      {$R$ GRU $h{=}256$\\[1pt]{\scriptsize trained separately}};
    \node[lookup, above=6mm of R] (M) {$M$ Markov counts};
    \node[trained, below=6mm of R] (A)
      {$A$ pooled MLP\\[1pt]{\scriptsize trained separately}};
    \draw[flow] (pf.east) -- (M.west);
    \draw[flow] (pf.east) -- (R.west);
    \draw[flow] (pf.east) -- (A.west);
    % ---- the per-candidate feature row ----------------------------------
    \node[op, right=6mm of R] (fx)
      {per-candidate feature row\\$[\,s_M,\ s_R,\ s_A,\ \log(1{+}\text{plays}),$\\$\text{recency},\
       \text{artist match},\ \text{genre match}\,]$};
    \draw[flow] (M.east) -- (fx.north west);
    \draw[flow] (R.east) -- (fx.west);
    \draw[flow] (A.east) -- (fx.south west);
    \node[refuted, right=5mm of fx] (xgb)
      {XGBoost meta-learner\\$200$ rounds, depth $6$\\[1pt]{\scriptsize\texttt{binary:logistic}}};
    \draw[flow] (fx) -- (xgb);
    \node[op, below=9mm of xgb] (rk)
      {score all $19{,}402$ items\\$\to$ top-10};
    \draw[flow] (xgb) -- (rk);
    % ---- the two failure mechanisms, marked on the diagram --------------
    \node[refuted, text width=56mm, align=flush left, anchor=north west] (tr)
      at ([yshift=-10mm]fx.south west)
      {\textbf{Training rows}: $1$ positive $+$ $50$ popularity-sampled
       negatives per train session, i.e.\ \num{5723} positives, against the
       \num{73632} teacher-forced steps its own $R$ leg trains on
       ($12.9\times$ less signal) --- and a decision surface calibrated on
       $51$ candidates ($0.26\%$ of the catalogue) is then asked to rank all
       \num{19402}.};
    \draw[grad] (tr.north east) -- (xgb.south west);
    \node[note] at ([yshift=-3.5mm]current bounding box.south west)
      {\textbf{Why the shape matters.} A $200$-round depth-$6$ booster is
       piecewise \emph{constant}, so ranking \num{19402} candidates with it
       produces large tied blocks that the harness breaks by item index. The
       fixed z-blend of Figure~\ref{fig:topo-blend} is instead a strictly
       monotone continuous function of the leg scores: it preserves each leg's
       internal ordering exactly and can only trade the legs off against one
       another.};
  \end{tikzpicture}
  \caption{\texttt{seq-stacker}, the learned counterpart of
    Figure~\ref{fig:topo-blend}'s fixed rule: the same separately-trained base
    legs, but their per-candidate scores become columns of a feature row
    alongside four cheap candidate features, and a gradient-boosted tree
    learns the combination. It earns a diagram because its failure is
    informative --- with the \texttt{markov}, \texttt{gru} and \texttt{ann}
    legs it reaches Recall@10 \num{0.085255} (\num{122} of \num{1431}),
    \emph{below its own GRU base leg} at \num{0.12299}, paired $\Delta$
    $-\num{0.0377}$ $[-0.056, -0.020]$ CI$<$0. The red callout names the two
    mechanisms: supervision starvation, since the meta-learner gets one
    labelled event per session where its base leg gets one per prefix
    position, and a train/serve candidate mismatch. The booster is drawn solid
    because it is trained and red because it is refuted; the feature row and
    the ranking step are dashed because they hold no parameters. Visual
    convention as in Figure~\ref{fig:topo-gru}. This is one of four
    independent learned-combiner nulls, the others being reciprocal-rank
    fusion, oracle weights fitted on the test labels, and the jointly-trained
    dual tower of Figure~\ref{fig:topo-dualgru}.}
  \label{fig:topo-stacker}
\end{figure}


\paragraph{\texttt{seq-embed}: a design, not a result.} One family in scope has
no measurements at all and must be read as an argument rather than a row.
\texttt{predictors/seq\_embed.py} exists (written 2026-07-26) and implements a
\emph{learned} item-embedding table in three modes --- \texttt{frozen},
\texttt{residual} and \texttt{free} --- with InfoNCE forced because cosine has a
trivial collapse optimum on a learned table. It is registered in
\texttt{registry.toml} but does \emph{not} appear in
\texttt{GET /api/predictors}, has no model definition, and has \textbf{zero runs
on any dataset}. Its interest is that the data budget of
Section~\ref{sec:intro} predicts its outcome: a \texttt{free} table is the
GRU4Rec/SASRec formulation, one trainable vector per item, estimated here from a
\emph{median of one observation}, and structurally capped at Recall@10
\num{0.4493} because an unobserved item's vector never receives a gradient. (That
formulation --- one free trainable vector per item, learned end-to-end from
sequences --- is the standard one in the session-recommendation literature, where
GRU4Rec and SASRec are its recurrent and self-attentive canonical instances. It
is the obvious thing to try and it is the one thing this corpus is least able to
support.) This
document quotes the dataset facts (\num{73632} steps, \num{15816} observed
items, median frequency \num{1}, \num{3294} items at frequency $\ge5$) as
measured, and everything about \texttt{seq-embed}'s performance as
\emph{unmeasured}.

{\rightskip=0pt plus 2em \emph{Sources: \texttt{predictors/} sources read directly;
\texttt{PROJECT-FACTS.md}.}\par}

\clearpage
\section{The board: what the trained models actually score}
\label{sec:board}

\emph{Results.} Table~\ref{tab:models} is the in-scope board and
Figure~\ref{fig:mc-model-landscape} is the same information as a landscape, with
the two non-parametric reference bars drawn as \emph{lines} rather than compared
rows. Figure~\ref{fig:nexttrack-phase2-families} shows the family ordering that
first established it. Three reading rules apply.

First, the promoted champion is the \Rp+M+\Cp{} blend at \num{0.21174}
(\num{303} of \num{1431}); the \num{0.22292} (\num{319}) cell above it is a
temperature scan point ($\tau=\num{0.10}$) whose $+\num{0.011}$ nominal edge was
judged \emph{sub-noise} and whose default was \textbf{held} --- it is the
highest Recall@10 ever measured on this split and it is \emph{not} the
champion's number.

Second, the single-model bar is the plain GRU at \num{0.12299}, which has been
reproduced \emph{eleven} times bit-identically
($\num{0.12299091544374564}=\num{176}/\num{1431}$, MRR agreeing to
\num{3.1e-11}) across hosts and across a predictor code edit. Same-configuration
reproduction spread on this deterministic path is therefore exactly zero, which
is why the noise discussion in Section~\ref{sec:power} is about
\emph{between}-configuration jitter, not seeds.

Third, \texttt{seq-ann}'s only run in the entire system is on
\dsid{seq-20260715-030509} (PCA-128), \emph{not} the canonical split. Because
both splits share the same session index its $n_{\text{test}}$ is also
\num{1431}, which makes the number look canonical when it is not; it is
therefore daggered and excluded from the best-cell bolding, and the defensible
GRU-versus-ANN comparison is the same-space one in
Section~\ref{sec:recurrence}. Seven further cells carry a
\textsuperscript{\S}: they are \emph{nominal single-run} comparisons for which no
paired bootstrap CI was ever computed --- the two temperature cells, the
projection-rank cell, the two BPR cells, the anti-eager $\beta$ arm and the
canonical-split LSTM. The reports call some of those ties and one a collapse;
none of those calls is backed by a paired test, and none of the cells is eligible
for a best-cell bold.

\begin{table}[ht]
  \centering
  \caption{The in-scope trained models on the canonical split
    \dsid{seq-20260715-131139} (\num{1431} test sessions), sorted by Recall@10.
    Eval-time MMR arms are held back to Table~\ref{tab:mmr} because they share
    bit-identical weights with their parent. The \emph{family} column abbreviates
    the predictor: blend $=$ \texttt{seq-blend}, next $=$ \texttt{seq-nexttrack},
    dual $=$ \texttt{seq-dualgru}, ann $=$ \texttt{seq-ann}, stack $=$
    \texttt{seq-stacker}, embed $=$ \texttt{seq-embed}, lookup $=$ a
    non-parametric table over the dataset. \textbf{Directions:} R@10, MRR,
    art@10, gen@10 \emph{and} $H$ are all higher-better. $H$ is the
    \emph{generation-3 (live)} composite and is \textbf{not} comparable to
    Table~\ref{tab:hgen}'s gen-1/gen-2 columns (Section~\ref{sec:metrics}); the
    champion's low $H$ (\num{0.008603}, the smallest non-lookup value in the
    column) is the two-board inversion of Section~\ref{sec:metrics}, not a
    defect, and it is not resolvable against the \numrange{0.02127}{0.02487}
    healthy band by anything in this table. For absolute scale, chance Recall@10
    for a random top ten is $10/\num{19402}\approx\num{0.00052}$ and the
    popularity lookup reaches \num{0.00070} (\num{1} hit); the art@10 and gen@10
    chance floors are much higher and have \emph{not} been measured, so read
    those two columns only across rows.
    \textbf{Bolding.} Two rows are bolded, the promoted champion and the prior
    crown it replaced. Best cell per metric is bolded among the
    \emph{on-split, unrefuted, non-nominal} rows only --- which excludes the
    \texttt{seq-ann} row (off-split), the Markov row (a lookup, kept in sorted
    position for scale rather than as a compared model), the $\tau0.10$ scan cell
    and the other \textsuperscript{\S}-marked cells (nominal, no paired CI), and
    every \texttt{seq-dualgru}, pre-encoder, stacker and $\beta$ arm shown to
    \emph{bound} a refuted axis. On that rule the champion holds R@10 and MRR and
    the $h425$ width control holds art@10, gen@10 and $H$.
    ``---'' marks a value not recorded: the nine remote-only runs never emitted
    \texttt{artist\_conc\_at\_k} and so could not be rescored to generation 3
    (their stored $H$ is gen 1 or gen 2 and is therefore withheld rather than
    printed in a gen-3 column), and no parameter count is recorded anywhere for
    \texttt{seq-stacker} or \texttt{seq-ann}.
    \textsuperscript{\dag}~\emph{not} on the canonical split ---
    \dsid{seq-20260715-030509}, PCA-128; do not read it in this column.
    \textsuperscript{\ddag}~scan cell whose nominal edge was judged sub-noise;
    the $\tau=\num{0.07}$ default was held.
    \textsuperscript{\S}~\emph{nominal single-run comparison} --- no paired
    bootstrap CI exists for this cell (Section~\ref{sec:notestablished}).
    \emph{Sources: \texttt{2026-\allowbreak 07-\allowbreak 15c-\allowbreak nexttrack-\allowbreak phase2-\allowbreak model-\allowbreak sweep.md},
    \texttt{2026-\allowbreak 07-\allowbreak 18b-\allowbreak nexttrack-\allowbreak projection-\allowbreak registration-\allowbreak and-\allowbreak scan.md},
    \texttt{2026-\allowbreak 07-\allowbreak 25-\allowbreak nexttrack-\allowbreak dual-\allowbreak tower-\allowbreak fusion-\allowbreak scan.md},
    \texttt{2026-\allowbreak 07-\allowbreak 26-\allowbreak nexttrack-\allowbreak pre-\allowbreak encoder-\allowbreak scan.md},
    and \texttt{2026-\allowbreak 07-\allowbreak 26-\allowbreak nexttrack-\allowbreak mmr-\allowbreak lambda-\allowbreak frontier.md}.}}
  \label{tab:models}
  {\scriptsize\setlength{\tabcolsep}{2.2pt}\begin{tabular}{llrcccccl}
    \toprule
    Model / arm & family & params & R@10 (raw) & MRR & art@10 & gen@10 & $H$ gen 3 & Note \\
    \midrule
    proj-blend $\tau0.10$\textsuperscript{\ddag\S} & blend & 431808 & 0.22292 (319) & 0.1273 & 0.3424 & 0.4032 & --- & held sub-noise \\
    \best{\Rp+M+\Cp{} $\tau0.07$} & blend & 431808 & \best{0.21174 (303)} & \best{0.1220} & 0.3326 & 0.4011 & 0.008603 & promoted, servable \\
    proj-blend rank-64\textsuperscript{\S} & blend & 407232 & 0.20545 (294) & 0.1169 & 0.3326 & 0.3990 & --- & rank is not the gain \\
    proj-blend $\tau0.05$\textsuperscript{\S} & blend & 431808 & 0.19287 (276) & 0.1111 & 0.3354 & 0.3997 & --- & $\tau$ ablation \\
    \best{R+M+C} & blend & 394944 & 0.18588 (266) & 0.1025 & 0.3305 & 0.3941 & --- & prior crown \\
    R+M two-leg & blend & 394944 & 0.17121 (245) & 0.0952 & 0.3291 & 0.3983 & 0.008259 & leg-ablation ref. \\
    proj-blend BPR full\textsuperscript{\S} & blend & 431808 & 0.15793 (226) & 0.0939 & 0.3187 & 0.3913 & --- & InfoNCE $\gg$ BPR \\
    proj-blend BPR rk-64\textsuperscript{\S} & blend & 407232 & 0.15234 (218) & 0.0909 & 0.3194 & 0.3906 & --- & worst proj.\ cell \\
    C2 GRU h454 \emph{ctrl} & next & 969936 & 0.12509 (179) & 0.0539 & 0.3375 & 0.4542 & 0.023988 & width ctrl, tie \\
    GRU h256 $\beta{=}0.05$\textsuperscript{\S} & next & 394944 & 0.12369 (177) & 0.0534 & 0.3368 & 0.4500 & 0.024755 & weak knob \\
    T1-f0 bare/preMLP256 & dual & 888256 & 0.12369 (177) & 0.0547 & 0.3333 & 0.4647 & --- & a repair, not a win \\
    GRU h256 (bar) & next & 394944 & 0.12299 (176) & 0.0534 & 0.3375 & 0.4500 & 0.024538 & single-model bar \\
    C1 GRU h425 \emph{ctrl} & next & 871017 & 0.12089 (173) & 0.0522 & \best{0.3410} & \best{0.4577} & \best{0.024872} & tie, $2.2\times$ params \\
    LSTM h256\textsuperscript{\S} & next & 510144 & 0.11880 (170) & 0.0493 & 0.3277 & 0.4549 & 0.022941 & GRU $\approx$ LSTM \\
    D-LL-f0 latent/latent & dual & 789696 & 0.11880 (170) & 0.0563 & 0.3368 & 0.4563 & 0.023785 & best 07-25 dual: tie \\
    GRU h297 \emph{ctrl} & next & 494697 & 0.11321 (162) & 0.0520 & 0.3284 & 0.4500 & --- & outlier, CI$<$0 \\
    S1 MLP256$\to$GRU256 & next & 493504 & 0.10901 (156) & 0.0446 & 0.3235 & 0.4689 & 0.024102 & pre-encoder harmful \\
    B2 first-order Markov & lookup & 0 & 0.10692 (153) & 0.0694 & 0.2984 & 0.3885 & 0.004409 & reference bar only \\
    D-LL latent/latent, $f1$ & dual & 871872 & 0.10273 (147) & 0.0424 & 0.3089 & 0.4577 & 0.022841 & fusion head harmful \\
    T4 preMLP both towers & dual & 1068992 & 0.09783 (140) & 0.0383 & 0.2991 & 0.4389 & --- & placement irrelevant \\
    pooled-MLP twin\textsuperscript{\dag} & ann & --- & 0.09015 (129) & 0.0394 & 0.2991 & 0.4354 & --- & off-split, PCA-128 \\
    D-DC delta/cummean & dual & 871872 & 0.08665 (124) & 0.0338 & 0.2579 & 0.4242 & --- & beats both parents \\
    XGB stacker (M+G+A) & stack & --- & 0.08526 (122) & 0.0490 & 0.3005 & 0.3753 & --- & below its base leg \\
    D-CC cummean/cummean & dual & 871872 & 0.06639 (95) & 0.0266 & 0.2460 & 0.4060 & --- & degenerate arm \\
    D-DD delta/delta & dual & 871872 & 0.05800 (83) & 0.0282 & 0.2523 & 0.4228 & --- & worst real arm \\
    learned item table & embed & --- & --- & --- & --- & --- & --- & design only, no runs \\
    \bottomrule
  \end{tabular}}
\end{table}

\begin{figure}[ht]
  \centering
  \includegraphics[width=0.8\textwidth]{mc_model_landscape}
  \caption{The whole in-scope landscape, and the things it is easy to misread.
    \num{27} trained arms on the canonical split plus the off-split ANN twin
    (\num{28} rows in all), ordered by Recall@10 with raw hits of \num{1431} and
    coloured by predictor family (\texttt{seq-blend} blue,
    \texttt{seq-nexttrack} green, \texttt{seq-dualgru} gray,
    \texttt{seq-stacker} cyan, \texttt{seq-ann} orange). The champion is the blue
    star. The first-order Markov lookup (\num{0.107}) is drawn only as a
    reference \emph{line}, never as a row, because it is a table over the
    dataset; the single-model GRU bar (\num{0.123}) is drawn \emph{both} as its
    own dotted line and as a row (bold label, green square), because it is
    simultaneously a compared model and the anchor every architecture $\Delta$ is
    measured against. Two markers are hollow for integrity: the ANN row because
    it is PCA-128 off-split, and the $\tau0.10$ cell because its nominal edge was
    held sub-noise --- it is the highest Recall@10 on the split and it is
    \emph{not} the champion. Only the champion carries a recorded marginal 95\%
    CI $[\num{0.191},\num{0.233}]$; the separate $\pm\num{0.0124}$ bar states the
    paired-$\Delta$ resolution of the instrument, which is what actually decides
    the adjacent rows.}
  \label{fig:mc-model-landscape}
\end{figure}

\begin{figure}[ht]
  \centering
  \includegraphics[width=0.85\textwidth]{nexttrack_phase2_families}
  \caption{The family ordering that set the shape of the board, from the
    Phase-2 sweep. Green: the \emph{fixed} z-blends. \textbf{Read the in-figure
    label ``NEW CHAMPION ROW'' historically:} the three-leg R+M+C on PCA-192 at
    \num{0.1859} was the champion at the date of that sweep and is the
    \emph{prior} crown here, superseded by \Rp+M+\Cp{} at \num{0.21174}
    (Section~\ref{sec:projection}), which is not in this figure. Blue: other
    fixed blends. Gray: single models --- the LSTM on PCA-128 (\num{0.1272}) is
    the best single model in that sweep, the pooled ANN twin (\num{0.0901}) the
    worst. Red: the \emph{learned} XGBoost stackers at \numrange{0.0853}{0.0867},
    below every single model they were built on top of. The first-order Markov
    bar (\num{0.107}) is the incumbent reference line. Read left to right the
    figure already contains the document's thesis: fixed fusion of unlike
    signals beats every single model, and learning the fusion loses.
    \emph{Source: \texttt{2026-\allowbreak 07-\allowbreak 15c-\allowbreak nexttrack-\allowbreak phase2-\allowbreak model-\allowbreak sweep.md}.}}
  \label{fig:nexttrack-phase2-families}
\end{figure}

\clearpage
\section{Fusion wins because its legs are correct on almost disjoint sessions}
\label{sec:fusion}

\emph{Rationale.} The champion is a score-level fusion of three legs, and the
naive story --- ``ensembling reduces variance'' --- is wrong here in an
instructive way. What each leg \emph{knows} is structurally different. The
Markov leg is a pure count lookup keyed on the single last prefix item $u$: the
empirical $P(v\mid u)$ over training bigrams. Its counts are shrunk toward a
fallback by a \emph{trust weight} $n_u/(n_u+8)$, where $n_u$ is how many times $u$
was ever followed by anything in training --- so an item seen once is trusted
$1/9$ and an item seen a hundred times almost fully, and the remaining mass goes
to an artist/genre back-off distribution plus a tiny popularity tie-break. (This
is Dirichlet--multinomial smoothing with a pseudo-count of 8; the point is only
that a leg built on counts must have \emph{something} to say when the counts are
thin.) It can only score items it has \emph{seen follow something}. The GRU leg
regresses a point in the item space from the ordered prefix and ranks the whole
catalog by cosine to it; it never needs to have seen the target at all, only to
land near it. The content-kNN leg is the same idea with no training (max cosine
of a candidate to any prefix item). So the fusion is not two estimates of one
quantity, it is a \emph{union of two retrieval mechanisms} --- identity memory
and content geometry --- whose failure modes are close to complementary.

\emph{Results.} The warm/cold decomposition makes the mechanism visible.
Partitioning the \num{1431} test targets into \num{643} warm and \num{788} cold
by the train-seen test and recomputing recall in each half from each run's
\texttt{predictions.json} gives: popularity \num{0.0016} warm / \num{0.0000}
cold; Markov \num{0.2255} (\num{145}/\num{643}) warm / \num{0.0102}
(\num{8}/\num{788}) cold; the GRU \num{0.1198} (\num{77}/\num{643}) warm /
\num{0.1256} (\num{99}/\num{788}) cold; R+M \num{0.2255} / \num{0.1269}; R+M+C
\num{0.2255} / \num{0.1536}; and the champion \num{0.2395}
(\num{154}/\num{643}) / \num{0.1891} (\num{149}/\num{788}).
Figure~\ref{fig:mc-cold-warm} plots this. The asymmetry is extreme: Markov is
$22\times$ better on warm targets than cold ones, while the GRU is \emph{flat}
across the split and in fact very slightly better on cold --- exactly what a
content-space retrieval should be. Hit-set overlap on the identical sessions:
$|M|=\num{153}$, $|R|=\num{176}$, $|M\cap R|=\num{55}$, $|M\cup R|=\num{274}$,
Jaccard \num{0.201}; \num{98} sessions are Markov-only and \num{121} GRU-only.
R+M lands \num{245} hits, \num{211} of them inside that \num{274}-session union
plus \num{34} sessions \emph{neither} leg hit alone. Note the scope carefully:
this overlap argument is about the \textbf{R+M stage}, computed with the
\emph{unprojected} GRU leg $R$. The champion's \num{303} hits do include
\num{109} sessions outside $M\cup R$ entirely, but that number mixes the fusion
effect with the projection upgrade of Section~\ref{sec:projection} --- it is not
evidence about fusion, and the champion's extra reach is measured there, not
here. The fusion win itself is CI-backed. On the AE-64-lineage artifact
\dsid{seq-20260713-014226} where the two-leg blend was first established, paired
$\Delta$ over the Markov bar on that split is $+\num{0.0657}$
$[+\num{0.049},+\num{0.083}]$, three-seed confirmed; on the canonical split the
content third leg is $+\num{0.0147}$ $[+\num{0.003},+\num{0.026}]$ over
same-space R+M, both CIs above zero. The warm/cold and overlap counts are
\emph{single-run observations with no CI}, though they reproduce
independently-reported off-server figures and rest on deterministic runs.

\emph{Interpretation --- and three alternatives ruled out.} (a) ``More models
help.'' No: adding a leg of the \emph{same kind} does nothing.
\emph{All four numbers in this paragraph are on the 2026-07-13/14
AE-64-lineage artifact \dsid{seq-20260713-014226}, not the canonical split, and
must be read only against each other} --- the single GRU on that artifact scores
${\approx}\num{0.089}$, not \num{0.12299}. There, adding the pooled ANN as a
third leg gives R+A+M \num{0.1684} against R+M's \num{0.1726}, i.e.\ slightly
worse, and R+A alone is \num{0.0894}, essentially that split's GRU alone.
Substituting \emph{item2vec} --- a skip-gram co-listening embedding, the
word2vec construction applied to sessions-as-sentences --- for the content leg
gives R+M+I \num{0.1349}, well below R+M. And on the canonical split a fourth
leg of \emph{session-kNN} type (retrieve whole \emph{sessions} similar to the
current one and vote over what followed them; V-SKNN in the literature) drags
R+M+C+V to \num{0.1474}, paired $\Delta$ $-\num{0.038}$ CI below zero --- and it
fails for precisely the mechanism above, being structurally warm-only against a
\SI{55}{\percent} cold test set in a single-user corpus where there is no
cross-user neighbourhood to borrow from. Diversity of signal \emph{kind}, not
model count, is what pays. (b) ``The blend just inherits
Markov's warm hits.'' No: R+M lands \num{100} cold hits where Markov alone lands
\num{8}, and the entire cold half of the champion's performance (\num{149} of
\num{303}) is carried by legs Markov contributes nothing to. (c) ``A smarter
fuser would do better.'' That is Section~\ref{sec:fail-combiners}, and it is
five independent nulls.

{\rightskip=0pt plus 2em \emph{Sources: \texttt{2026-\allowbreak 07-\allowbreak 14-\allowbreak nexttrack-\allowbreak blend-\allowbreak and-\allowbreak ensemble.md},
\texttt{2026-\allowbreak 07-\allowbreak 15c-\allowbreak nexttrack-\allowbreak phase2-\allowbreak model-\allowbreak sweep.md},
and \texttt{2026-\allowbreak 07-\allowbreak 18-\allowbreak nexttrack-\allowbreak literature-\allowbreak fit-\allowbreak campaign.md}.}\par}

\begin{figure}[ht]
  \centering
  \includegraphics[width=0.95\textwidth]{mc_cold_warm}
  \caption{The one decomposition that explains the champion. \emph{Left}:
    Recall@10 split into warm and cold halves for each champion leg and each
    blend stage. The Markov leg collapses to \num{0.010} on cold targets
    (\num{8} of \num{788}) against \num{0.226} warm, whereas the projected
    content leg reads cold \num{0.152} $\approx$ warm \num{0.151} and the base
    GRU is fractionally \emph{better} on cold than warm --- so the fusion lifts
    both halves at once (\num{154}/\num{643} warm, \num{149}/\num{788} cold).
    The \Rp{} and \Cp{} bars are hatched because those standalone figures are
    report provenance from an off-server driver, with no server run and no raw
    hit counts. \emph{Right}: the split's composition --- the warm fraction
    \emph{is} the seen-items-only ceiling, Recall@10 $\le\num{0.4493}$, against
    $\num{788}/\num{1431}=\SI{55.1}{\percent}$ cold. Any model that cannot
    retrieve an unseen item is arguing over the left-hand \SI{45}{\percent}
    only.}
  \label{fig:mc-cold-warm}
\end{figure}

\clearpage
\section{The learned projection wins because it rebuilds the retrieval metric}
\label{sec:projection}

\emph{Rationale.} Before the projection, both the GRU leg and the content-kNN
leg retrieved by cosine in the same unsupervised PCA-192 space. That space is
variance-ordered and completely target-agnostic: it was fit to \emph{reconstruct}
the item feature matrix, so its dominant directions are whatever has the most
raw variance --- on this corpus loudness, era and popularity --- and not ``tracks
that follow each other''. The top ten principal components carry
\SI{65}{\percent} of the variance with essentially zero text loading. The
projection replaces the \emph{metric} rather than the model.

\emph{What is actually learned.} \texttt{seq\_blend.fit\_projection} learns a
single square matrix $W$ (\num{192}$\times$\num{192}, no bias, initialized to
the identity; Adam, lr \num{1e-3}, \num{40} epochs, batch \num{512}) from
\emph{train} consecutive item pairs $(u,v)$ with self-loops dropped. The
objective is the in-batch InfoNCE of Section~\ref{sec:machinery} at
$\tau=\num{0.07}$: for a batch of \num{512} pairs, embed
$a_i=\mathrm{normalize}(Wz_{u_i})$ and $p_j=\mathrm{normalize}(Wz_{v_j})$ and
require each item's true successor to outscore the \num{511} other rows'
successors.

\emph{What $W$ is, stated exactly.} Up to the row normalization, ranking by
$\langle Wz_u, Wz_v\rangle = z_u^{\top}(W^{\top}W)z_v$ means the projection is a
\textbf{learned Mahalanobis metric} on the original \num{192} dimensions, with
positive-semidefinite Gram matrix $W^{\top}W$. Two things follow immediately.
First, by the singular value decomposition $W=U\Sigma V^{\top}$ a general square
$W$ is a rotation, then a per-axis rescaling, then \emph{another} rotation --- so
it can both re-orient the axes along which ``followed by'' varies and rescale
them, which is more than a diagonal reweighting of the existing PCA coordinates
could do. Second, $z_u^{\top}(W^{\top}W)z_v$ is \emph{symmetric in $u$ and $v$ by
construction}: the object being learned is a ``these two go together'' relation,
not a directed ``$u$ is followed by $v$'' one. That is exactly why the directional
two-tower ablation below is the right test of whether the symmetry is a
limitation, and why its losing is informative rather than incidental.
\textbf{The honest gap:} the spectrum of $W^{\top}W$ has never been inspected on
this corpus, so the natural next sentence --- \emph{which} directions were
stretched and which shrunk --- is not something this document can assert. What is
measured is only the consequence: after the fit, cosine in the new space behaves
like transition affinity. Two of the three legs are defined by cosine in that
space, so both improve at once; the Markov leg is item-ID based and is untouched
by construction. And because $W$ is a function of \emph{content
only}, an item that never appeared in training still has a position in the
projected space --- this is the one channel with cold reach, and the cold half is
\SI{55}{\percent} of the test set. The supervision is leak-free: train pairs and
intrinsic latents only.

\emph{Results.} \Rp+M+\Cp{} reaches Recall@10 \num{0.2117}
$[\num{0.191},\num{0.233}]$ at MRR \num{0.1220}, paired $\Delta$ over the R+M+C
crown $+\num{0.0259}$ $[+\num{0.012},+\num{0.041}]$, CI entirely above zero.
Both shared-metric legs improve and the untouched leg does not: content-kNN
$C\to\Cp$ goes \num{0.0685}$\to$\num{0.1516} --- a $2.2\times$ lift on a leg with
\emph{no trainable parameters of its own} (the $k$NN rule is parameter-free; the
\num{36864} parameters of the shared $W$ it now retrieves inside are counted once,
in the champion's total), so the entire change is the geometry ---
the GRU leg $R\to\Rp$ goes \num{0.1230}$\to$\numrange{0.1719}{0.1761} across
seeds, and Markov stays at \num{0.1069}. Unlike the deterministic GRU path this
result has genuine seed variance (random init, shuffled batches): over seeds
1337/7/99 the blend reads \num{0.2117}/\num{0.2110}/\num{0.2145} with paired
$\Delta$ $+\num{0.026}$/$+\num{0.025}$/$+\num{0.029}$, every CI above zero. It
is then hardened on a split disjoint in both index and time
(\dsid{seq-20260718-211238}, the earlier chronological band, byte-identical item
space, cold rate \num{0.64}): \num{0.1705} against the crown's \num{0.1279} and
R+M's \num{0.1174}, paired $\Delta$ $+\num{0.0426}$
$[+\num{0.0280},+\num{0.0573}]$ --- a \emph{wider} margin on the colder split,
exactly what the mechanism predicts, since when more targets are cold the
cold-reaching leg carries proportionally more of the win.
Figure~\ref{fig:nexttrack-projection-hardening} shows both splits.

\emph{The axis ablations, and their status.} Four cells sharpen what is doing
the work, and none of them carries a paired CI --- they are \emph{nominal
single-run} comparisons and are labelled as such here and in
Table~\ref{tab:models}. InfoNCE $\gg$ BPR: replacing the softmax-over-in-batch-
negatives objective with pairwise BPR --- Bayesian Personalized Ranking, the
standard pairwise recommender loss, which for each positive samples \emph{one}
negative and maximizes $\log\sigma(s_{\text{pos}}-s_{\text{neg}})$ --- collapses the
blend to \num{0.1579} and \num{0.1523}, roughly \num{0.055} below the default
and \emph{below even the unprojected crown}, so the specific shape of the
contrastive loss is load-bearing and not merely ``some pairwise supervision''.
Full rank $\approx$ rank 64: constraining $W$ to $64\times192$ gives
\num{0.2055} against \num{0.2117}, a nominal $-\num{0.006}$ and a tie inside the
$\pm\num{0.015}$ band --- \emph{the gain is the supervision, not the capacity}.
Temperature is mildly monotone over the tested range
(\num{0.05}$\to$\num{0.1929}, \num{0.07}$\to$\num{0.2117},
\num{0.10}$\to$\num{0.2229}) and $\tau0.10$'s $+\num{0.0112}$ sits inside the
band, so the decision rule did not re-promote it. A directional two-tower
$W_q/W_k$ variant --- which drops the $W^{\top}W$ symmetry by giving $u$ and $v$
separate maps --- scored \Cp{} \num{0.1349} against the symmetric single $W$'s
\num{0.1516}, so the symmetry is not the binding constraint, and a symmetric
``these two go together'' relation is what this corpus supports.

\emph{Interpretation --- the nuance that is easy to state wrongly.} A separate
campaign asked whether the frozen equal-weighted PCA-192 vector is itself the
bottleneck, building eight alternative \num{192}-dimensional item spaces under
parity gates and evaluating each under two readouts. Under the \emph{frozen} GRU
readout, \emph{six of eight} beat it with paired CI above zero
(\numrange{+0.0252}{+0.0356}), so at the raw level the PCA-192 geometry
genuinely is a handicap. Under the \emph{champion} readout, with the projection
re-fit on each space, \emph{not one} variant beats it: the two nominally
positive spaces straddle zero and two are significantly worse. The correct
reading is that pre-whitening the input and the learned projection are the
\emph{same fix applied in two places}, so they do not stack --- a full-rank
linear map over the same \num{192} dimensions can rescale and rotate away the
loud inert directions, and InfoNCE gives it the supervision to know \emph{which}
ones. It is emphatically not ``the representation does not matter'' (it
demonstrably does, six times over with CI above zero) and not ``whitening helps
the champion'' (it does not). A sub-finding sharpens it: dropping the numeric and
acoustic blocks \emph{helps} the frozen GRU ($+\num{0.0335}$) but \emph{hurts}
the champion ($-\num{0.0168}$), so the lesson is not ``drop the loud blocks'' but
``let a supervised map decide their weight''. The honest caveat is that this
cross-readout conclusion rests on a \emph{null} on the champion readout, which
is weaker evidence than a win. Because those eight spaces are dataset variants
rather than models they are otherwise out of scope here.

{\rightskip=0pt plus 2em \emph{Sources: \texttt{2026-\allowbreak 07-\allowbreak 18-\allowbreak nexttrack-\allowbreak literature-\allowbreak fit-\allowbreak campaign.md},
\texttt{2026-\allowbreak 07-\allowbreak 18b-\allowbreak nexttrack-\allowbreak projection-\allowbreak registration-\allowbreak and-\allowbreak scan.md},
and \texttt{2026-\allowbreak 07-\allowbreak 18c-\allowbreak nexttrack-\allowbreak representation-\allowbreak exploration.md}.}\par}

\begin{figure}[ht]
  \centering
  \includegraphics[width=0.8\textwidth]{nexttrack_projection_hardening}
  \caption{The champion's lead \emph{widens} on the colder split, which is the
    mechanism's own prediction. Two splits $\times$ three blends: the canonical
    split (cold rate \num{0.55}) and a disjoint earlier-holdout split (cold rate
    \num{0.64}), with \Rp+M+\Cp{} in green (\num{0.2117} / \num{0.1705}), the
    unprojected R+M+C crown in blue (\num{0.1859} / \num{0.1279}) and the
    two-leg R+M in gray (\num{0.1712} / \num{0.1174}). The annotated paired
    deltas are $+\num{0.0259}$ $[+\num{0.012},+\num{0.041}]$ and $+\num{0.0426}$
    $[+\num{0.0280},+\num{0.0573}]$, both CIs above zero. Absolute recall is
    lower on the right-hand split by construction; the verdict is the paired
    $\Delta$, not the level. \emph{Sources:
    \texttt{2026-\allowbreak 07-\allowbreak 18-\allowbreak nexttrack-\allowbreak literature-\allowbreak fit-\allowbreak campaign.md},
    and \texttt{2026-\allowbreak 07-\allowbreak 18b-\allowbreak nexttrack-\allowbreak projection-\allowbreak registration-\allowbreak and-\allowbreak scan.md}.}}
  \label{fig:nexttrack-projection-hardening}
\end{figure}

\clearpage
\section{Recurrence beats pooling because the pooled summaries are already inside the state}
\label{sec:recurrence}

\emph{Rationale.} \texttt{seq-ann} is built specifically as the GRU's
non-recurrent twin, and the comparison is unusually clean because everything
except the architecture is reused verbatim: \texttt{seq\_ann.py} imports
\texttt{make\_batches} and \texttt{step\_loss} from \texttt{seq\_nexttrack}, so
objective, batching and early stopping are byte-for-byte identical, and the live
API confirms matched hyperparameters ($h=256$, InfoNCE, \num{40} epochs, lr
\num{1e-3}, dropout \num{0.1}, seed 1337). Where the GRU carries a hidden state
updated by learned gates, the ANN builds a fixed-width bag-of-prefix at each
step --- the concatenation of the \emph{mean} of the prefix latents so far and
the \emph{last} item's latent --- and feeds it through a ReLU MLP. Both are
strictly causal, so neither leaks. The difference is therefore exactly: a mean
plus a special-cased most-recent item, versus a state whose gates decide per
dimension and per step how much of the incoming item to write and how much
context to forget. A mean is permutation-invariant and weights a track from
twenty steps ago as much as the one just played; a gated state can hold a
stylistic thread through an interruption and can represent ``A then B''
differently from ``B then A''.

\emph{Results.} Head-to-head on the same PCA-128 space
\dsid{seq-20260715-030509}: ANN \num{0.09015} (\num{129} of \num{1431}) at MRR
\num{0.03941} against the GRU anchor's \num{0.11461} (\num{164}) at MRR
\num{0.0489}, paired $\Delta$ $-\num{0.0245}$
$[-\num{0.039},-\num{0.011}]$, CI entirely below zero; the LSTM on the same
space reaches \num{0.12718} (\num{182}), so the pooled model is beaten by both
recurrent cells. Where it loses is informative and was recomputed here from the
same predictions pipeline: ANN warm \num{0.1026} (\num{66}/\num{643}) / cold
\num{0.0799} (\num{63}/\num{788}) against the GRU's \num{0.1182} / \num{0.1117},
i.e.\ it gives up \num{10} warm hits but \num{25} cold ones --- the damage
concentrates on the cold half. \emph{Why} it concentrates there is a conjecture
this document cannot settle, and it must not be stated as the earlier
Markov-versus-GRU argument was: neither the ANN nor the GRU has any item-memory
mechanism, so ``you cannot lean on having seen the item'' does not distinguish
them --- both are pure content retrieval on both halves. The plausible reading is
about \emph{density}: warm targets are by definition items the user played
before, so they sit in the well-trodden, densely covered regions of the latent
space where a sloppy predicted point still lands one of ten neighbours, whereas
cold targets sit in sparse regions where the precision of the point is what
decides. That is untested here --- no per-region density measurement exists ---
and it is offered as the likely mechanism, not a result. Hit-set
overlap: \num{92} of the ANN's \num{129} hits are also GRU hits and only
\num{37} are ANN-only, so the pooled model is close to a strict subset --- which
is also why it is redundant as a blend leg (Section~\ref{sec:fusion}). These
warm/cold and overlap figures are single-run, uncertainty-unquantified counts.

\emph{Interpretation --- two triangulations and one reverse test.} First, the
explicit pooled summaries are drastically weaker \emph{even with recurrence over
them}: a two-tower arm reading only the running prefix mean scores \num{0.06639}
(\num{95}) and one reading only step-to-step differences \num{0.05800}
(\num{83}), against \num{0.12299} for a single GRU over raw latents. Both of
those ablation arms carry the harmful \texttt{fusion\_layers}${}=1$ head
(Section~\ref{sec:fail-head}), so part of the ${\sim}\num{0.057}$ gap is the head
rather than the view; the clean within-regime contrast is against the
latent/latent arm at the \emph{same} depth, \num{0.06639} against \num{0.10273},
a still-decisive $-\num{0.036}$. The direction of the finding survives; its
magnitude does not. Second,
handing a raw-latent tower an \emph{explicit} running mean or difference adds
nothing --- latent/cummean versus latent/latent is $-\num{0.00559}$
$[-\num{0.01747},+\num{0.00699}]$ and latent/delta is $-\num{0.01118}$
$[-\num{0.02306},+\num{0.00070}]$, both straddling zero --- because a GRU's
update gate is a leaky accumulator (a running mean is inside its hypothesis
space) and its candidate/gate interaction covers first differences. An
interpretability read on that arm is consistent with the derived views being
lossy functions of what the raw-latent tower already held: the two towers'
representation geometries are only partly distinct (CKA \num{0.458} --- a
\emph{middling} similarity, so there is real new linear structure between them)
yet the sparse-autoencoder decomposition recovers \emph{zero} next-item concepts
in the second tower that the first does not already carry. New geometry, no new
concepts. Third, the reverse test shows the fusion
machinery is not broken: pairing the two \emph{individually hopeless} views with
each other beats both parents with CI above zero ($+\num{0.02865}$ and
$+\num{0.02027}$) at a hopeless absolute \num{0.08665}. Two-view co-adaptation
is a real mechanism; it simply has nothing left to contribute once a raw-latent
tower is present. Figure~\ref{fig:mc-view-redundancy} draws the argument.
Capacity is separately ruled out, with the caveat of
Section~\ref{sec:fail-topology}: on PCA-192 single-GRU width is flat at the
endpoints, $h=425$ at $2.2\times$ the parameters tying $h=256$ ($-\num{0.00210}$
$[-\num{0.01118},+\num{0.00629}]$) and $h=454$ at $2.5\times$ also tying, so the
${\sim}\num{0.025}$ recurrence gap cannot be bought with parameters. And the
GRU's advantage \emph{widens} multi-step: on three-step continuation
($n=\num{1172}$) it beats the first-order bigram on both suffix\_recall@10
(\num{0.0963} versus \num{0.0869}) and cont\_prec@10 (\num{0.0273} versus
\num{0.0236}).

\emph{The honest weakness.} The ANN datapoint is a \emph{single unreplicated run
on PCA-128}. It has never been measured on the canonical PCA-192 split, so it
must not be read in the same column as the \num{0.12299} bar
(Table~\ref{tab:models}, dagger), and the only defensible GRU-versus-ANN
statement is the same-space one above.

{\rightskip=0pt plus 2em \emph{Sources: \texttt{2026-\allowbreak 07-\allowbreak 15c-\allowbreak nexttrack-\allowbreak phase2-\allowbreak model-\allowbreak sweep.md},
\texttt{2026-\allowbreak 07-\allowbreak 25-\allowbreak nexttrack-\allowbreak dual-\allowbreak tower-\allowbreak fusion-\allowbreak scan.md},
and \texttt{2026-\allowbreak 07-\allowbreak 26-\allowbreak nexttrack-\allowbreak pre-\allowbreak encoder-\allowbreak scan.md}.}\par}

\begin{figure}[ht]
  \centering
  \includegraphics[width=0.95\textwidth]{mc_view_redundancy}
  \caption{Why a second tower cannot add information here. \emph{Left}: the
    three tower views --- \texttt{latent} $x_t$, \texttt{delta}
    $x_t-x_{t-1}$, \texttt{cummean} the running prefix centroid --- are all
    deterministic, strictly causal functions of the \emph{same} prefix, so
    $\sigma(\text{tower B})\subseteq\sigma(\text{tower A})$: tower B holds a
    re-compression of what tower A already receives, never new information, and
    the recurrence's own gates already accumulate a running mean and a first
    difference. Measured, the best dual arm \num{0.11880} (\num{170} of
    \num{1431}) ties the \num{0.12299} (\num{176}) baseline, paired $\Delta$
    $-\num{0.00419}$ $[-\num{0.013},+\num{0.004}]$. \emph{Right}: what genuine
    multi-modal fusion looks like, where each stream has a private region the
    other cannot supply --- measured exactly once on this corpus, with
    \texttt{delta}$+$\texttt{cummean} beating \emph{both} parents with CI above
    zero ($+\num{0.02865}$, $+\num{0.02027}$) at a hopeless absolute
    \num{0.08665}. The inner ellipses are drawn strictly inside tower A's
    boundary because the panel asserts a subset relation. The right panel's
    ``CKA \num{0.458}'' is a similarity, not a fraction of newness: it says the
    two towers' representation geometries are only \emph{partly} alike, which
    makes the accompanying ``$+0$ new next-item concepts'' the informative half of
    the pair. All the recall figures quoted on this figure are
    \texttt{fusion\_layers}${}=0$ arms except the parents of the
    \texttt{delta}$+$\texttt{cummean} contrast, which are
    \texttt{fusion\_layers}${}=1$ throughout and therefore internally comparable.}
  \label{fig:mc-view-redundancy}
\end{figure}

\begin{figure}[ht]
  \centering
  \includegraphics[width=0.82\textwidth]{nexttrack_music_scatter}
  \caption{The families differ in \emph{kind}, not only in rank. music@10
    against Recall@10 for all \num{53} next-track runs, one marker per run,
    coloured and shaped by predictor family: \texttt{seq-blend} (blue circles),
    \texttt{seq-nexttrack} (green squares), \texttt{seq-markov} (purple
    diamond), \texttt{seq-ann} (orange triangle), \texttt{seq-stacker} (cyan
    down-triangles), \texttt{seq-popularity} (red $\times$ at the bottom left).
    The two metrics correlate at $r=\num{0.609}$ but crown different winners:
    the content-regressing \texttt{seq-nexttrack} family lands the
    \emph{musically closest misses} (best music@10 \num{0.482}) while the blend
    family wins exact recall. \textbf{One in-figure annotation is stale and must
    be read against Section~\ref{sec:board}}: the point it labels ``Recall@10
    champion (\num{0.223})'' is the $\tau0.10$ scan cell whose nominal edge was
    judged sub-noise and whose default was held --- it is the highest Recall@10
    measured on this split, but the champion is \Rp+M+\Cp{} at \num{0.21174}. Read against Section~\ref{sec:metrics}, this is
    why music@10 makes a usable grounding factor: it is a smooth superset of
    recall, not a substitute for it. \emph{Source:
    \texttt{2026-\allowbreak 07-\allowbreak 20-\allowbreak musical-\allowbreak distance-\allowbreak metric.md}.}}
  \label{fig:nexttrack-music-scatter}
\end{figure}

\clearpage
\section{MMR is the first honest crown lever, and it de-concentrates artists for free}
\label{sec:mmr}

\emph{Rationale.} Every prior attempt to raise holisticness@10 raised it by
predicting worse (Section~\ref{sec:metrics}). MMR re-ranking is different for a
structural reason: it is eval-time only, so the trained weights are
bit-identical across $\lambda$ arms and $\lambda=1$ is an exact identity gate,
and its single greedy penalty moves \emph{three} of the crown's four factors the
same way at once. ild is literally the quantity the penalty maximizes, so it
rises by construction. Eagerness falls as a \emph{side effect} with no notion of
artist anywhere in the rule, because same-artist tracks are sonic near-neighbours
and the max-similarity penalty preferentially evicts them. And music@10
\emph{rises}, because a sonically spread top-ten covers more of the
neighbourhood around the truth, so the best-of-ten cosine to the true next track
goes up. Only exact recall pays.
Figure~\ref{fig:mc-mmr-geometry} draws the geometry;
Figure~\ref{fig:nexttrack-mmr-tradeoff} shows the original offline sweep.

\emph{Results.} Ten arms on the canonical split, same host, serialized, none
failed; every arm's $H$ recomputed from \texttt{predictions.json} under one code
path and validated against storage to $|\Delta|\le\num{2.1e-7}$; all statistics
paired bootstrap with all five facets and hence $H$ recomputed for both arms on
each resampled index set. \textbf{Eight of eight $\lambda<1$ arms lift $H$ with
paired $\Delta H$ CI entirely above zero} (Table~\ref{tab:mmr}). The best
\emph{free} operating point is the single GRU at $\lambda=\num{0.7}$ with
\texttt{mmr\_pool} \num{50}: $H$ \num{0.030355}, $+\SI{22.1}{\percent}$ over the
all-time healthy bar \num{0.024872} and \num{5.4} measured half-widths above it,
at Recall@10 \num{0.120196} (\num{172}) with $\Delta$ recall $-\num{0.002908}$
$[-\num{0.008386},+\num{0.002795}]$ \emph{straddling} zero and $|\Delta|$ far
inside the $\pm\num{0.015}$ band. It also beats the nearest pure-$\lambda$ point
on $H$ by $+\num{0.000889}$ $[+\num{0.000497},+\num{0.001360}]$ at a recall tie,
so its rank as best free arm is resolved rather than nominal.

\emph{Results --- the facets, with the arm named.} The three-factors-together
move is clearest on the GRU's deepest $\lambda$ arm, \textbf{B3
($\lambda=\num{0.7}$, pool \num{200})}, which is \emph{priced} rather than free:
there, as $\lambda$ goes \num{1.0}$\to$\num{0.7}, ild rises
\num{0.4536}$\to$\num{0.5314}, artist\_adj falls
\num{0.3474}$\to$\num{0.3129}, artist\_conc falls
\num{0.3559}$\to$\num{0.2707} and music@10 \emph{rises}
\num{0.4558}$\to$\num{0.4747}, with mood\_coh moving only
\num{0.4788}$\to$\num{0.4585} ($-\SI{4.2}{\percent}$). The \emph{free} arm B4
(same $\lambda$, pool \num{50}) makes the same moves less far --- ild
\num{0.5017}, artist\_adj \num{0.3270}, artist\_conc \num{0.3019}, music@10
\num{0.4670}, mood\_coh \num{0.4671} --- which is precisely the trade the pool
axis buys back. Every quoted facet is a cosine in the content metric $M$
(Section~\ref{sec:machinery}), which is also the space MMR's penalty acts in;
that shared space is why one knob moves four numbers. The instrument transferred
exactly: an offline sweep on the champion checkpoint had projected
\num{0.2103}/\num{0.1936}/\num{0.1642}/\num{0.1097} for $\lambda$
\num{0.9}/\num{0.7}/\num{0.5}/\num{0.3} and the native runs returned
\num{0.21034}/\num{0.19357}/\num{0.16422}/\num{0.10971}.

\begin{table}[ht]
  \centering
  \caption{The eval-time MMR $\lambda$ frontier, ten arms on
    \dsid{seq-20260715-131139} (\num{1431} test sessions), grouped by anchor and
    sorted within group by $\lambda$. $\lambda$ and \texttt{mmr\_pool} are
    \emph{eval-time} parameters: within a group the trained weights are
    bit-identical, so $\lambda=\num{1.0}$ is an exact identity control. Paired
    $\Delta$ is a per-session paired bootstrap against that control (\num{2000}
    resamples, rng 1337); CI bounds are given to three decimals, and a CI
    straddling zero is a \emph{tie}. $H$ is the \emph{generation-3 (live)}
    composite. Tiers are the campaign's own decision rule: TIER-1 \emph{free}
    ($\Delta H$ above zero at a recall tie), TIER-2 \emph{priced} ($\Delta$
    recall CI below zero), TIER-3 \emph{rejected} (frontier mapping only).
    \textbf{Six definitions were registered on 2026-07-26 and \emph{none} was
    promoted}; the recall champion is unchanged. \textbf{Bolding.} The bolded row
    is B4, the best \emph{free} operating point, and the two bolded cells are the
    best $H$ and $\Delta H$ among the free arms. \textbf{Recall is deliberately
    unbolded in this table}: it is not a maximization target here but a
    \emph{constraint} --- each group's $\lambda=\num{1.0}$ control is the highest
    recall in its group by construction, since MMR can only reorder a list the
    control already produced, and a free arm's merit is that its recall $\Delta$
    \emph{straddles} zero rather than that it is large. Recall also must not be
    compared \emph{across} the two groups: the anchors are different models.
    \emph{Source:
    \texttt{2026-\allowbreak 07-\allowbreak 26-\allowbreak nexttrack-\allowbreak mmr-\allowbreak lambda-\allowbreak frontier.md}.}}
  \label{tab:mmr}
  {\scriptsize\setlength{\tabcolsep}{2.2pt}\begin{tabular}{llccccl}
    \toprule
    Arm & $\lambda$ / pool & R@10 (raw) & $\Delta$ recall [95\% CI] & $H$ gen 3 & $\Delta H$ [95\% CI] & Tier \\
    \midrule
    B0 GRU h256 & 1.0 / --- & 0.12299 (176) & --- & 0.024538 & --- & control \\
    B1 GRU h256 & 0.9 / 200 & 0.12089 (173) & $-0.0021$ [$-.006$, $+.002$] & 0.026761 & $+0.0022$ [$+.0016$, $+.0029$] & TIER-1 free \\
    B2 GRU h256 & 0.8 / 200 & 0.11950 (171) & $-0.0036$ [$-.009$, $+.001$] & 0.029458 & $+0.0049$ [$+.0040$, $+.0059$] & TIER-1 free \\
    \best{B4 GRU h256} & 0.7 / 50 & 0.12020 (172) & $-0.0029$ [$-.008$, $+.003$] & \best{0.030355} & \best{$+0.0058$ [$+.0048$, $+.0068$]} & TIER-1 free \\
    B3 GRU h256 & 0.7 / 200 & 0.11461 (164) & $-0.0085$ [$-.016$, $-.001$] & 0.034171 & $+0.0096$ [$+.0083$, $+.0110$] & TIER-2 priced \\
    \midrule
    A0 \Rp+M+\Cp & 1.0 / --- & 0.21174 (303) & --- & 0.008603 & --- & control \\
    A1 \Rp+M+\Cp & 0.9 / 200 & 0.21034 (301) & $-0.0013$ [$-.008$, $+.005$] & 0.010435 & $+0.0018$ [$+.0015$, $+.0022$] & TIER-1 free \\
    A2 \Rp+M+\Cp & 0.7 / 200 & 0.19357 (277) & $-0.0181$ [$-.029$, $-.008$] & 0.015289 & $+0.0067$ [$+.0057$, $+.0077$] & TIER-2 priced \\
    A3 \Rp+M+\Cp & 0.5 / 200 & 0.16422 (235) & $-0.0473$ [$-.061$, $-.033$] & 0.023576 & $+0.0150$ [$+.0132$, $+.0168$] & TIER-3 rejected \\
    A4 \Rp+M+\Cp & 0.3 / 200 & 0.10971 (157) & $-0.1019$ [$-.120$, $-.084$] & 0.031626 & $+0.0230$ [$+.0206$, $+.0256$] & TIER-3 rejected \\
    \bottomrule
  \end{tabular}}
\end{table}

\emph{Interpretation --- the mechanism, resolved.} A pre-registered question
asked whether artist concentration is \emph{responsive} to a sonic penalty or
\emph{sticky}. It resolves unanimously to responsive: in eight of eight arms
$|\Delta\text{artist\_conc}|$ exceeds $|\Delta\text{artist\_adj}|$, by
\numrange{1.59}{3.06}$\times$, where the sticky hypothesis required a ratio
below \num{0.5}. So a rule with no notion of artist in it de-concentrates artists
faster than it reduces seed echo, and the planned explicit artist penalty is
consequently dropped as the next move. The ratio also \emph{shrinks}
monotonically as $\lambda$ falls (\num{3.06}$\to$\num{1.59} on the champion),
i.e.\ the easy de-concentration is bought first and the mechanism saturates.
Because the crown takes $\max(\text{artist\_adj},\text{artist\_conc})$ and
concentration falls ${\sim}2.5\times$ faster, the two must cross: past that
switch the eagerness factor becomes the stiff one and each further recall point
buys progressively less crown. Note that this verdict rests on the \emph{sign
and ratio} of point estimates across eight arms --- its strength is unanimity,
not an interval on the ratio. The $\Delta H$ and $\Delta$ recall figures, and
hence the tiers, are paired-CI-backed.

\emph{Why the plain GRU benefits and the champion does not.} The pre-registered
expectation was the opposite: the champion's eagerness is $1.9\times$ the GRU's,
so its $(1-\text{eager})$ factor enters at \num{0.308} against \num{0.644} and
it has far more headroom. The measurement says no --- at $\lambda=\num{0.7}$ the
GRU reaches $H$ \num{0.034171} against the champion's \num{0.015289}, a
$2.24\times$ gap whose \emph{absolute} size grows as $\lambda$ falls
(\num{0.01633} at $\lambda=\num{0.9}$, \num{0.01888} at $\lambda=\num{0.7}$)
even though the ratio itself narrows slightly, from $2.56\times$ to $2.24\times$,
over the only two $\lambda$ values where both anchors were run
(Figure~\ref{fig:mc-lambda-frontier}). The reason is in the facets. At
\emph{matched} $\lambda=\num{0.7}$ and pool \num{200}, the champion's mood\_coh
falls \num{0.3982}$\to$\num{0.3643} ($-\SI{8.5}{\percent}$) against the GRU's
\num{0.4788}$\to$\num{0.4585} ($-\SI{4.2}{\percent}$) --- a factor of two, not a
collapse. Pushed to $\lambda=\num{0.3}$ the champion's mood\_coh does collapse to
\num{0.2378} ($-\SI{40}{\percent}$), but no GRU arm was run that deep, so the
often-quoted $\SI{40}{\percent}$-versus-$\SI{4}{\percent}$ contrast is
\emph{not} like-for-like and the honest matched figure is the $2\times$ one.
The champion's recall is built substantially on artist
adjacency supplied by its Markov leg, so diversifying away from the seed artist
also walks it out of the prefix's mood neighbourhood --- it pays twice. MMR's
value therefore scales with the \emph{health} of the other three factors, not
with headroom to de-eager. Figure~\ref{fig:mc-candidate-eigen} shows the same
asymmetry from the inside, as the geometry of the candidate set the two models
produce.

\emph{What is not established.} Nothing here is promoted. Six definitions were
\emph{registered} on 2026-07-26; a mandatory second-split robustness confirm and
user approval are still pending; the best free arm's rank over its neighbour is
resolved at only ${\sim}\num{0.9}$ half-widths of its own comparison, on one
split. And the campaign leaves a live calibration question open: the grounding
floor \num{0.34} was chosen to \emph{cancel} recall-destroying $H$ gains, and
MMR is the first lever that moves music@10 the \emph{other} way, so the constant
is currently operating outside the regime it was calibrated for.

{\rightskip=0pt plus 2em \emph{Sources: \texttt{2026-\allowbreak 07-\allowbreak 22-\allowbreak session-\allowbreak holisticness-\allowbreak and-\allowbreak antieager-\allowbreak levers.md},
and \texttt{2026-\allowbreak 07-\allowbreak 26-\allowbreak nexttrack-\allowbreak mmr-\allowbreak lambda-\allowbreak frontier.md}.}\par}

\begin{figure}[ht]
  \centering
  \includegraphics[width=0.95\textwidth]{mc_mmr_geometry}
  \caption{Sonic diversification de-eagers artists for free. \emph{Left}: the
    greedy rule
    $\arg\max_i \lambda\,\mathrm{rel}(i)-(1-\lambda)\max_{j\in\text{selected}}\cos(M_i,M_j)$
    applied at $\lambda=\num{0.7}$ over a candidate pool. Each candidate keeps its
    own artist marker \emph{shape}; a red ring marks the three candidates the
    re-rank drops from the top eight, and all three are the \emph{same} artist as
    an already-picked track sitting close to it. Each pick's green ring is an
    \emph{iso-penalty contour}, not a cutoff: the penalty
    $\max_j\cos(M_i,M_j)$ is continuous, so a sufficiently relevant candidate
    survives inside a contour --- two do, which is the mechanism working correctly
    rather than an inconsistency. \emph{Right}: the two resulting
    lists --- pure relevance returns seven of eight from one artist, MMR returns
    four artists --- with no artist term anywhere in the rule. The pool is
    schematic (real similarity is cosine between \num{517}-dimensional vectors
    and is not drawable in two dimensions, so planar proximity stands in) and the
    list's own Herfindahl arithmetic is the illustration's; the \emph{measured}
    facet moves sit in the separate callout, and they are arm B3 (single GRU,
    \texttt{mmr\_pool} \num{200}, TIER-2 \emph{priced}), not the free arm: at
    $\lambda$ \num{1.0}$\to$\num{0.7}, ild \num{0.4536}$\to$\num{0.5314}, artist\_conc
    \num{0.3559}$\to$\num{0.2707}, artist\_adj \num{0.3474}$\to$\num{0.3129},
    music@10 \num{0.4558}$\to$\num{0.4747}, with
    $|\Delta\text{conc}|>|\Delta\text{adj}|$ in eight of eight arms at
    \numrange{1.59}{3.06}$\times$.}
  \label{fig:mc-mmr-geometry}
\end{figure}

\begin{figure}[ht]
  \centering
  \includegraphics[width=0.9\textwidth]{mc_lambda_frontier}
  \caption{The same lever is free on one model and ruinous on the other. The
    MMR $\lambda$ frontier in (Recall@10, generation-3 $H$) space, one traced
    curve per anchor at fixed \texttt{mmr\_pool} \num{200}: the single GRU's thick
    trace climbs almost \emph{vertically} ($H$ \num{0.0245}$\to$\num{0.0342} over
    $\lambda$ \num{1.0}$\to$\num{0.7}) while the champion's slides
    \emph{sideways}, spending $-\num{0.102}$ recall to buy $+\num{0.023}$ $H$.
    Note that the GRU trace's top end, $H$ \num{0.0342}, is arm B3 --- TIER-2
    \emph{priced}, $\Delta$ recall CI below zero --- so the vertical climb runs
    one step past the free region; the star is the genuinely free best arm B4 and
    it is drawn \emph{off} the trace because it changes \texttt{mmr\_pool} to
    \num{50} and so does not belong to the same fixed-pool sweep. Tier is encoded
    as a ring (free / priced / rejected / $\lambda{=}1$ control), the $\lambda$
    labels carry their pool where it differs, and both anchors' own
    $\pm\num{0.015}$ recall floors are marked (\num{0.108} for the GRU,
    \num{0.197} for the champion). The dashed double arrow is an eye guide between
    the two anchors at matched $\lambda=\num{0.7}$, not a fitted curve.
    \textbf{Nothing plotted here is promoted.}}
  \label{fig:mc-lambda-frontier}
\end{figure}

\begin{figure}[ht]
  \centering
  \includegraphics[width=0.8\textwidth]{nexttrack_mmr_tradeoff}
  \caption{The original offline $\lambda$ sweep on the champion checkpoint,
    which is what motivated running the native arms of Table~\ref{tab:mmr}.
    Left axis: Recall@10 (red) falling
    \num{0.2117}$\to$\num{0.2103}$\to$\num{0.1936}$\to$\num{0.1642}$\to$\num{0.1097}
    across $\lambda$ off/\num{0.9}/\num{0.7}/\num{0.5}/\num{0.3}, with the
    $\pm\num{0.015}$ tie band shaded around the $\lambda$-off point. Right axis:
    ild@10 rising \num{0.406}$\to$\num{0.815}, artist\_adj@10 falling
    \num{0.610}$\to$\num{0.293}, music@10 rising \num{0.454}$\to$\num{0.492} ---
    the three-factors-together move, including the grounding factor, that
    distinguishes this lever from the degenerate ones. \textbf{The right axis
    carries the individual facets, each on its own $0$--$1$ scale, not $H$
    itself} --- the figure's own axis label calls them ``holisticness metrics'',
    which is the facet suite, and $H$ (${\sim}\num{0.01}$--$\num{0.03}$ here)
    is nowhere on this plot. The transfer to the native runs was exact to four
    decimals. \emph{Source:
    \texttt{2026-\allowbreak 07-\allowbreak 22-\allowbreak session-\allowbreak holisticness-\allowbreak and-\allowbreak antieager-\allowbreak levers.md}.}}
  \label{fig:nexttrack-mmr-tradeoff}
\end{figure}

\begin{figure}[ht]
  \centering
  \includegraphics[width=0.9\textwidth]{mc_candidate_eigen}
  \caption{The eigen-shape of a next-track belief, and why the champion is the
    harder model to diversify. \textbf{Provenance first, because it is unlike
    every other figure here: this panel is not in the record.} It is two
    \texttt{extend} invocations run by the author for this document, on one
    five-item seed prefix, single seed, ten \emph{generation steps} --- a
    generation step being one round of ``rank, take the top candidate, append it
    to the prefix, re-rank'', i.e.\ the autoregressive continuation mode rather
    than the single-step protocol of Figure~\ref{fig:protocol}. No server run
    record exists for either invocation and no campaign report contains these
    numbers; they are reproducible only from
    \texttt{docs/figures/make\_concept\_figures.py}. Read it as an illustration
    of a mechanism, never as evidence. \emph{Left}: the participation ratio
    $1/\sum_i s_i^2$ (Section~\ref{sec:machinery}) of the candidate-set variance
    spectrum at each step. The champion blend falls \emph{monotonically} over
    four successive steps (pr \num{2.35}$\to$\num{1.79}$\to$\num{1.23}$\to$\num{1.06},
    mean over all ten \num{3.14}) while the single GRU fluctuates
    non-monotonically without trend over the whole run
    (\numrange{2.07}{8.64}, mean \num{5.87}) --- note that the GRU's own minimum,
    \num{2.07} at step 3, falls \emph{below} the band drawn for it, so the
    contrast is monotone-collapse versus no-trend, not inside-band versus
    outside. \emph{Middle}: three spectra --- the GRU at its step-5 point
    ($\mathrm{pr}=\num{7.11}$, leading share \num{0.256}) against the blend's
    first and fourth steps, the latter putting \num{0.972} of the variance on a
    single axis. \emph{Right}: what that means geometrically --- a spread cloud
    leaves \emph{what} comes next open, a collapsed one leaves only \emph{how far}
    along one line; each cloud's drawn aspect is the measured
    $\sqrt{\lambda_1/\lambda_2}$. The one external check available is that the
    blend's collapsed step reproduces the album-lock signature recorded in
    \texttt{predictors/seq\_extend.py} (pr \num{1.23} versus \num{1.25}, leading
    share \num{0.897} versus \num{0.89}).}
  \label{fig:mc-candidate-eigen}
\end{figure}

\clearpage
\section{The failures, and the measured mechanism in each case}
\label{sec:failures}

Five design axes have been refuted on this corpus. What makes them worth a
section is not the nulls themselves --- Section~\ref{sec:power} shows the
instrument could not have detected a moderate win anyway --- but that in four of
the five the \emph{mechanism} was subsequently measured rather than argued.

\subsection{Dual-tower view fusion: the second tower's input is measurable with respect to the first's}
\label{sec:fail-views}

\emph{Rationale and mechanism.} The failure can be stated before any run. The
three tower views are all deterministic, strictly causal functions of the
\emph{same} prefix, so $\sigma(\text{view}_t)\subseteq\sigma(x_0..x_t)$: tower B
holds a lossier summary of what tower A already receives and can add
optimization-level diversity, never new information. The second half of the
mechanism is about the ensemble story: joint end-to-end training is precisely
what \emph{destroys} the independence that makes combination pay. The fixed
z-blend works because its legs are trained separately and never see each other's
gradients; two towers trained jointly through one loss co-adapt into a single
estimator with two halves. That is why the same corpus supports a \num{0.212}
blend and refuses a \num{0.119} joint dual tower.

\emph{Results.} Zero of the eight arms in the 2026-07-25 dual-tower scan beat the
baseline with paired CI above zero. The best of those eight, latent/latent at
\texttt{fusion\_layers}${}=0$, reaches \num{0.11880} (\num{170}) against
\num{0.12299} (\num{176}), $-\num{0.00419}$
$[-\num{0.01258},+\num{0.00419}]$, a tie. (One \emph{later} dual arm, T1-f0 in
the 2026-07-26 batch, nominally exceeds the bar at \num{0.12369}; it also ties,
$+\num{0.00070}$ $[-\num{0.01188},+\num{0.01328}]$, and
Section~\ref{sec:fail-relu} explains why it is a repair rather than a win. No
dual arm anywhere clears zero.) Capacity is ruled out by the single-GRU width
control at \num{871017} parameters, which ties the bar at \SI{0.098}{\percent}
parity with the \texttt{fusion\_layers}${}=1$ dual arms (\num{871872}) --- note
that the parity is with those, not with the \num{789696}-parameter $f0$ arm just
quoted, which the control outweighs by \SI{10}{\percent} and still only ties.
Figure~\ref{fig:nexttrack-dualtower-matrix} gives the
\texttt{fusion\_layers}${}=1$ view matrix. The view attribution rule makes the
mechanism decidable: mixed cells
containing a raw-latent tower \emph{tie} their latent/latent parent and beat only
their weaker parent, whereas the one pair with no raw-latent tower at all beats
\emph{both} its parents with CI above zero. Complementarity is real and measured;
it has nothing left to contribute once a raw-latent tower is present.

\emph{A citation that could not be made.} Two further diagnostics were proposed
for this section --- a rank comparison of the concatenated two-tower
representation against each tower alone, and a cross-prediction $R^2$ between
identical-view towers --- and neither exists anywhere in this instance: not in
any \texttt{experiments/*.md} report, not in \texttt{PROJECT-FACTS.md}, not in
any run's metrics. Under this project's integrity rules the figures are omitted
and their numbers are not printed here even as illustration. The citable
substitute for the diversity-collapse claim is that two independently initialized
\emph{same-view} towers fused jointly tie both the baseline and the
parameter-matched control, buying none of the implicit-ensemble gain that
independently trained legs demonstrably do buy.

{\rightskip=0pt plus 2em \emph{Source: \texttt{2026-\allowbreak 07-\allowbreak 25-\allowbreak nexttrack-\allowbreak dual-\allowbreak tower-\allowbreak fusion-\allowbreak scan.md}.}\par}

\begin{figure}[ht]
  \centering
  \includegraphics[width=0.72\textwidth]{nexttrack_dualtower_matrix}
  \caption{The raw latent view subsumes the engineered ones. Recall@10 for the
    two-tower view pairing over $\{\texttt{latent},\texttt{delta},
    \texttt{cummean}\}$, \textbf{all six cells at
    \texttt{fusion\_layers}${}=1$}, one run each (upper triangle only; the grey
    lower cells were not run because the pairing is symmetric). Reading the
    diagonal is the argument: latent/latent \num{0.10273}, delta/delta
    \num{0.05800}, cummean/cummean \num{0.06639} --- the views that discard raw
    item identity are catastrophic even with a full recurrence over them.
    \textbf{Every cell sits below the single-GRU baseline \num{0.12299} named in
    the title, and the outlined best cell \emph{loses} to it}, paired $\Delta$
    $-\num{0.02027}$ $[-\num{0.03494},-\num{0.00697}]$, CI below zero. The dual
    arm that merely \emph{ties} the baseline is the \texttt{fusion\_layers}${}=0$
    variant of that same view pair (\num{0.11880},
    $-\num{0.00419}$ straddling zero), and it is \emph{not} in this matrix ---
    which is itself the point of Section~\ref{sec:fail-head}: at fixed views,
    removing the fusion head is worth more than any choice of views.
    \emph{Source:
    \texttt{2026-\allowbreak 07-\allowbreak 25-\allowbreak nexttrack-\allowbreak dual-\allowbreak tower-\allowbreak fusion-\allowbreak scan.md}.}}
  \label{fig:nexttrack-dualtower-matrix}
\end{figure}

\subsection{The MLP pre-encoder: a rectifier deletes half of a zero-centred latent}
\label{sec:fail-relu}

\emph{Rationale.} This is the sharpest mechanism on record because the design
left the ReLU as the only variable. A GRU's own input transform is already
linear: each gate computes $W_i x$. Prepend a linear pre-encoder $V$ of shape
$(\text{pre\_hidden}\times192)$ and the composition is
$W_i(Vx)=(W_iV)x$; when $\text{pre\_hidden}\ge\text{latent\_dim}$, $V$ is
generically full rank, so $W_iV$ ranges over exactly the same space of matrices
as $W_i$. \textbf{A linear pre-encoder is therefore provably
expressivity-neutral} --- it adds zero functions to the hypothesis space. The
only thing the pre-MLP actually changes is the ReLU, and the ReLU is applied to
PCA coordinates, which are zero-centred and near-symmetric \emph{by
construction}: about half of every input vector's coordinates are negative, and
half-wave rectification sets every one of them to exactly zero. Sign is not
recoverable afterwards, so this is a permanent deletion, not a
reparametrization. A ReLU layer is lossless on signed input only if it is
over-complete enough to re-encode the negative half-space --- canonically
$\mathrm{concat}[\mathrm{relu}(x),\mathrm{relu}(-x)]$ at $2\times$ width. At
$\text{pre\_hidden}=256$ on a \num{192}-dimensional input the ratio is only
$1.33\times$, so after rectification the layer carries \emph{fewer live
coordinates than the input has dimensions}. Figure~\ref{fig:mc-relu-sign-loss}
draws exactly this.

\emph{Results.} The decisive single-tower ablation loses: MLP256$\to$GRU256
scores \num{0.10901} (\num{156} of \num{1431}) against the baseline's
\num{0.12299} (\num{176}), paired $\Delta$ $-\num{0.01398}$
$[-\num{0.02657},-\num{0.00140}]$, CI below zero
(Figure~\ref{fig:nexttrack-preencoder-forest}). The mechanism was then
\emph{measured} by a counterfactual tap --- the model's own affine map with every
ReLU deleted. Next-item genre decodability (\num{11184} probe test rows, \num{25}
classes) reads raw latent \num{0.4033} accuracy / \num{0.7609} AUC $\to$
pre-linear \num{0.4084} / \num{0.7667} (the affine map is \emph{neutral},
${\sim}\num{0.8}\sigma$, exactly as $W_i(Vx)=(W_iV)x$ predicts) $\to$
pre-ReLU \num{0.3700} / \num{0.7508} (the ReLU is the \emph{whole} loss,
${\sim}\num{5.9}\sigma$). The zeroing is exact and quantified:
\texttt{frac\_exact\_zero} on the rectified tap is \num{0.5327}, matching
\texttt{frac\_negative\_coords} on the un-rectified tap to four digits, and
$256\times(1-\num{0.5327})=\num{119.6}$ live units to carry \num{192} signed
directions. The damage survives the recurrence: within one model, the tower fed
rectified input reads \num{0.3929}/\num{0.7845} against its identical twin fed
raw input at \num{0.4183}/\num{0.8004} (${\sim}\num{3.9}\sigma$), and the raw
tower reproduces the standalone plain GRU (\num{0.4235}/\num{0.8006}). The
repair signature confirms the direction of causation: adding a \emph{bare} tower
alongside the pre-encoded one recovers $+\num{0.01468}$
$[+\num{0.00278},+\num{0.02655}]$ over the single pre-encoded model yet only
\emph{ties} the plain baseline ($+\num{0.00070}$
$[-\num{0.01188},+\num{0.01328}]$) --- it restores an un-rectified path, it never
exceeds it. Dose is flat and nominally negative throughout
($\text{pre\_hidden}$ \num{0}/\num{128}/\num{256}/\num{384} $\to$
\num{0.10273}/\num{0.09993}/\num{0.09713}/\num{0.09993}, every adjacent pair
straddling zero, no interior peak), and placement is irrelevant --- rectifying
both towers ties rectifying one.

\emph{The honest confound.} The pre-MLP ends in Dropout, so pre-encoder arms
receive GRU-\emph{input} dropout that the bare arms lack (the baseline drops the
recurrent output only). Part of the measured harm may therefore be
over-regularization. A \texttt{pre\_dropout} parameter to disentangle it is
identified but not implemented.

{\rightskip=0pt plus 2em \emph{Source: \texttt{2026-\allowbreak 07-\allowbreak 26-\allowbreak nexttrack-\allowbreak pre-\allowbreak encoder-\allowbreak scan.md}.}\par}

\begin{figure}[ht]
  \centering
  \includegraphics[width=0.95\textwidth]{mc_relu_sign_loss}
  \caption{A rectifier on a signed PCA latent is lossy compression dressed as
    capacity. \emph{Left}: one real item's first \num{28} PCA-192 coordinates
    (blue positive, red negative) with its rectified image overlaid --- every
    negative coordinate is mapped to exactly \num{0} and, ReLU not being
    injective, the sign is unrecoverable. \num{0.5046} of all
    $\num{19402}\times\num{192}$ coordinates in this item space are negative by
    construction, and the measured rectified \emph{tap} inside the trained model
    reads \num{0.5327} (the two are kept distinct: the tap sits after a learned
    affine map). \emph{Right}: the width budget --- at the built width \num{256}
    only \num{119.6} units stay live to carry \num{192} signed directions, a
    plain ReLU at $2\times$ width still averages under \num{192}, and only the
    $\mathrm{concat}[\mathrm{relu}(x),\mathrm{relu}(-x)]$ split is exact, since
    $x=\mathrm{relu}(x)-\mathrm{relu}(-x)$ needs both halves. This panel is a
    budget, not a run measurement. It complements
    Figure~\ref{fig:nexttrack-rectifier}, which measures the same mechanism as
    decodability rather than drawing it.}
  \label{fig:mc-relu-sign-loss}
\end{figure}

\begin{figure}[ht]
  \centering
  \includegraphics[width=0.8\textwidth]{nexttrack_rectifier}
  \caption{The rectifier, not the extra layer, is the whole loss. Next-item
    genre decodability (probe accuracy; AUC printed inside each bar) across six
    representation taps, \num{11184} probe test rows over \num{25} classes,
    paired SE $\le\num{0.0065}$. Input side, left of the divider: raw latent
    (gray) $\to$ the same map with every ReLU \emph{deleted} (green,
    \num{0.4084}, neutral at ${\sim}\num{0.8}\sigma$) $\to$ the map as trained
    (red, \num{0.3700}, ${\sim}\num{5.9}\sigma$ worse). After the recurrence,
    right of the divider: the tower fed raw (green, \num{0.4183}) against its
    identical twin fed rectified (red, \num{0.3929}, ${\sim}\num{3.9}\sigma$),
    with the standalone plain GRU (blue, \num{0.4235}) reproducing the raw
    tower. The affine map is free; the half-wave rectifier is not, and its damage
    is still present after the GRU. \textbf{Two reading cautions.} The $y$ axis is
    truncated to roughly \numrange{0.355}{0.445} with no break drawn, so the bar
    \emph{heights} exaggerate what are \numrange{3}{8}\% relative differences in
    probe accuracy; the claim rests on the annotated $\sigma$ multiples, not on
    the visual drop. And a probe accuracy of \num{0.4033} on \num{25} classes is
    far above chance but far from perfect --- these are decodability
    \emph{differences}, not a statement that genre is fully recoverable from a
    latent. \emph{Source:
    \texttt{2026-\allowbreak 07-\allowbreak 26-\allowbreak nexttrack-\allowbreak pre-\allowbreak encoder-\allowbreak scan.md}.}}
  \label{fig:nexttrack-rectifier}
\end{figure}

\begin{figure}[ht]
  \centering
  \includegraphics[width=0.8\textwidth]{nexttrack_preencoder_forest}
  \caption{The pre-encoder axis at a glance, and the template for reading every
    paired comparison in this document. Paired per-session bootstrap
    $\Delta$Recall@10 against the reproduced single-GRU baseline (\num{2000}
    resamples, rng 1337, $n=\num{1431}$), nine arms sorted by $\Delta$, with the
    $\pm\num{0.015}$ practical tie band shaded and a black line at zero.
    Controls are blue, pre-encoder arms red; the thick row is the decisive
    ablation, MLP256$\to$GRU256 at $-\num{0.01398}$
    $[-\num{0.02657},-\num{0.00140}]$ (\num{156} hits against \num{176}). Note
    what the figure does \emph{not} show: no arm's CI clears zero on the
    positive side, which is the shape of every architecture scan on this corpus
    (Section~\ref{sec:power}). \emph{Source:
    \texttt{2026-\allowbreak 07-\allowbreak 26-\allowbreak nexttrack-\allowbreak pre-\allowbreak encoder-\allowbreak scan.md}.}}
  \label{fig:nexttrack-preencoder-forest}
\end{figure}

\subsection{The fusion head is actively harmful, and the harm grows with each upstream rectifier}
\label{sec:fail-head}

\emph{Rationale.} Why would extra capacity between the encoder and the readout
\emph{cost} accuracy? Because the retrieval geometry is fixed and external to the
network. Training is in-batch InfoNCE over $L2$-normalized vectors and serving
ranks the catalog by cosine to the predicted \num{192}-dimensional vector, so the
readout's only job is to land in a \emph{direction} in a signed linear space it
does not control. With \texttt{fusion\_layers}${}=0$ the map from
$\mathbf{R}^{512}$ to $\mathbf{R}^{192}$ is a single unconstrained linear
operator --- already the most general linear read of the recurrent state and
exactly matched to the geometry it must hit. Inserting
Linear$(512{\to}256)\to$ReLU$\to$Dropout$\to$Linear$(256{\to}192)$ does not
restrict which directions are \emph{reachable}: the final map is unconstrained
and $256>192$, so the image of the non-negative orthant of $\mathbf{R}^{256}$
under a generic linear map generically \emph{positively spans} all of
$\mathbf{R}^{192}$, and no prediction direction is lost. The cost is
parametrization and noise, not reachability, and it comes in two parts. First,
every direction with a negative coefficient must now be produced by
\emph{cancellation} among non-negative activations rather than directly, so the
loss spends gradient undoing a warp it never wanted --- and the warp is the same
half-wave sign deletion as Section~\ref{sec:fail-relu}, now applied to the
recurrent state instead of to the input. Second, and more concretely, the
trailing Dropout injects multiplicative noise directly into the very vector whose
\emph{direction} is being scored at retrieval time; a dropped unit does not
attenuate the score, it rotates the prediction. Neither part is a claim about
expressivity, and the argument's weight rests on the measurement below, which is
CI-backed three times.

\emph{Results.} Three independent paired measurements, all CI above zero in
favour of \emph{removing} the head: $+\num{0.01607}$
$[+\num{0.00210},+\num{0.03005}]$ in the redundant-view regime, $+\num{0.02027}$
$[+\num{0.00629},+\num{0.03354}]$ in the complementary-view regime, and
$+\num{0.02655}$ $[+\num{0.01398},+\num{0.03913}]$ with a pre-MLP on one tower.
The head also \emph{costs} parameters to do this: \num{871872} with it against
\num{789696} without, and \num{970432} against \num{888256} in the pre-encoder
batch --- roughly \SI{10}{\percent} more parameters for a reliably worse result.
\emph{On the ordering, be careful.} It is tempting to read
$+\num{0.016}\to+\num{0.020}\to+\num{0.027}$ as a penalty growing with the number
of upstream rectifiers, but the first two arms have \emph{no} upstream rectifier
at all --- both are bare-tower regimes differing only in which views the towers
read --- so that step cannot be attributed to rectifier count. All three
intervals also overlap heavily and no pairwise test between them was run. What is
defensible is weaker and still useful: the penalty \textbf{reproduces three times
across three different regimes}, and the largest of the three is the one arm that
does carry an upstream rectifier. The ordering itself is nominal.

\emph{Interpretation.} The unifying observation across the whole family is
checkable in the sources: \emph{the winning path contains zero half-wave
rectifiers}. A GRU (sigmoid/tanh gates, sign-preserving and bounded) plus one
\texttt{nn.Linear} to the latent dimension; and the single largest gain on
record, the champion's content projection, is a bias-free \emph{linear} map with
identity initialization trained by InfoNCE. Every arm that inserts a ReLU into
that path loses --- the pooled ANN twin, which is a three-layer ReLU MLP, loses
to the same-space GRU by $-\num{0.0245}$ $[-\num{0.039},-\num{0.011}]$. One
counter-direction must be stated honestly: under the \emph{old} three-factor
crown the \texttt{fusion\_layers}${}=1$ arms scored \emph{higher}, because the
extra warp diffuses the top ten (raising ild, lowering artist\_adj) at the cost
of exact hits. That apparent advantage disappears under the grounded composite,
where no pre-encoder arm reaches $\Delta H$ CI above zero at all.

{\rightskip=0pt plus 2em \emph{Sources: \texttt{2026-\allowbreak 07-\allowbreak 25-\allowbreak nexttrack-\allowbreak dual-\allowbreak tower-\allowbreak fusion-\allowbreak scan.md},
and \texttt{2026-\allowbreak 07-\allowbreak 26-\allowbreak nexttrack-\allowbreak pre-\allowbreak encoder-\allowbreak scan.md}.}\par}

\subsection{Four learned combiners, four nulls, and an oracle that says there is nothing to learn}
\label{sec:fail-combiners}

\emph{Rationale.} The champion's fusion weights are hard-coded equal thirds,
which looks like an obvious place to improve. It has now been attacked from four
directions and all four fail, and the reason is best stated in the order that
makes the weaker argument unnecessary. \textbf{First: there is nothing to learn.}
Fit the three blend weights directly on the \emph{test} labels --- outright
cheating, and therefore an upper bound on what \emph{any} weighting procedure
could achieve --- and the result is statistically indistinguishable from equal
thirds. If the optimum is a tie with the constant, then every parameter a learned
combiner spends is pure variance with no bias left to remove, and no amount of
data or cleverness changes that. \textbf{Second, and only then: even if there were
something to learn, it could not be fitted here.} The count-based legs are built
\emph{from} the train sessions, so on any train-derived validation slice they look
near-perfect --- they are being scored on data they memorized --- and a
weight-fitter therefore learns ``trust the count leg'', which is exactly wrong on
a test distribution where \SI{55}{\percent} of targets have no counts at all.

\emph{Results.} (1) The XGBoost stacker scores below its \emph{own} GRU base leg
in every arm, on both splits, and each $\Delta$ is against the base leg on the
\emph{same} split: on the canonical PCA-192 split \num{0.08526} against
\num{0.12299}, paired $\Delta$ $-\num{0.0377}$
$[-\num{0.056},-\num{0.020}]$; on the PCA-128 split
\dsid{seq-20260715-030509} \num{0.0867} and \num{0.0797} against that split's
GRU anchor of \num{0.1146}, paired $\Delta$ $-\num{0.0280}$
$[-\num{0.045},-\num{0.011}]$ and $-\num{0.0349}$
$[-\num{0.053},-\num{0.017}]$. All three CIs are below zero. A combiner scoring
\emph{below its own input} needs an explanation, and there are three separable
ones. Its
training matrix emits exactly \emph{one} positive per train session plus \num{50}
popularity-sampled negatives: \num{5723} positives against the GRU's
\num{73632} teacher-forced steps over the same sessions, $12.87\times$ less
signal. Its candidate distribution at fit time is one true item against \num{51}
popularity-sampled candidates --- \SI{0.258}{\percent} of the catalog --- and at
eval time it must rank all \num{19402}. And a \num{200}-round depth-6 booster is
piecewise \emph{constant}, so ranking \num{19402} candidates with a step function
produces large tied blocks broken arbitrarily by item index, which destroys
ordering information the GRU leg already had; the z-blend by contrast is a
strictly monotone continuous function of the legs' z-scores and preserves each
leg's internal ordering exactly. (2) Reciprocal-rank fusion, which keeps only
ranks and discards score magnitude, scores \num{0.1572}, $-\num{0.029}$ CI below
zero --- direct evidence that the z-blend's retention of calibrated magnitude is
the working part. (3) The \textbf{oracle}: coordinate-ascent weights fitted
directly on the \emph{test} labels, an upper bound on any achievable weighting,
reach \num{0.1901}, paired $\Delta$ $+\num{0.0042}$
$[-\num{0.004},+\num{0.013}]$, straddling zero. Fitting the same weights on a
train-val slice instead collapses to the degenerate vector $[0,1,0]$, pure
Markov. (4) Representation-level \emph{joint end-to-end} fusion, the strongest
form of the standing rebuttal, is Section~\ref{sec:fail-views}: eight arms, a
\SI{0.098}{\percent} parameter-matched control, best arm a tie, six of eight
losing outright.

\emph{Interpretation.} Put the two halves together and the lane is closed from
both ends. The combination weights would have to be estimated from \num{5723}
sessions with one usable event each, while the things being combined were
estimated from \num{73632} events --- and the oracle says the target of that
estimation is statistically indistinguishable from the constant. So the z-blend
wins because it is a \emph{zero-parameter} estimator of a quantity whose true
value is (indistinguishably) equal thirds. The adopted corollary is that no
further learned combiner should be proposed without a mechanism that is not
capacity, not co-adaptation and not score-level fusion. The residual uncertainty
is scope, not validity --- these nulls are established on one corpus with one
very unusual property (a single user, \SI{55}{\percent} cold) and should not be
generalized to session recommendation at large.

{\rightskip=0pt plus 2em \emph{Sources: \texttt{2026-\allowbreak 07-\allowbreak 14-\allowbreak nexttrack-\allowbreak blend-\allowbreak and-\allowbreak ensemble.md},
\texttt{2026-\allowbreak 07-\allowbreak 15c-\allowbreak nexttrack-\allowbreak phase2-\allowbreak model-\allowbreak sweep.md},
\texttt{2026-\allowbreak 07-\allowbreak 18-\allowbreak nexttrack-\allowbreak literature-\allowbreak fit-\allowbreak campaign.md},
and \texttt{2026-\allowbreak 07-\allowbreak 25-\allowbreak nexttrack-\allowbreak dual-\allowbreak tower-\allowbreak fusion-\allowbreak scan.md}.}\par}

\subsection{Depth, width, direction: a spent lever, a weak instrument, and an objective artifact}
\label{sec:fail-topology}

\emph{Rationale.} Sessions here average \num{13.84} events, so the recurrence's
job is to integrate ${\sim}14$ steps of a signed \num{192}-dimensional latent
into one direction --- a short-horizon integration problem, not one that wants a
deep feature hierarchy. Stacking layers buys hierarchy the task does not need
while inserting more saturating stages between input and readout and adding
inter-layer dropout.

\emph{Results.} Depth is flat then falls off a cliff: 1-layer \num{0.0894},
depth-2 \num{0.0922} (flat, in-band), depth-3 \num{0.0825} --- the only depth arm
to lose to the first-order Markov bar, paired $\Delta$ $-\num{0.0245}$
$[-\num{0.043},-\num{0.006}]$. A residual connection was never \emph{isolated}:
the only residual arm is depth-2 at $h=256$ (\num{0.0908}) and no depth-2 $h=256$
arm was run \emph{without} it, so the honest statement is only that the residual
arm fails to beat the depth-2 $h=128$ arm (\num{0.0922}) at twice the width ---
not that residual connections were tested and found neutral. Width is likewise closed on the rich space
--- $h=425$ at $2.2\times$ the parameters ties $h=256$ --- \emph{but the
instrument that established it is weaker than it looks}, and that caveat matters
more than the result. The four PCA-192 width points read \num{176}/\num{162}/
\num{173}/\num{179} hits for $h=$ \num{256}/\num{297}/\num{425}/\num{454}:
non-monotone, spanning \num{0.0119} recall with no trend, and the $h=297$
control actually \emph{lost} to $h=256$ ($-\num{0.00978}$
$[-\num{0.01747},-\num{0.00280}]$, CI below zero) while $h=454$ tied it. So
between-configuration jitter is ${\sim}\pm\num{0.01}$ --- roughly two-thirds of
the tie band --- while \emph{same}-configuration reproduction is bit-identical.
``Width is flat'' survives as a trend claim resting on the endpoints tying, not
as a per-width measurement, and a single unreplicated param-matched control
cannot carry a verdict. The prescribed replication (two runs on the second
leak-free split; the path is deterministic, so this needs a second \emph{split},
not a second seed) has not been bought.

\emph{The bidirectional arm, and why it is not a directionality result.} The
record's cleanest example of a confound that looks like a finding. A backward
pass would let step $t$ read item $t+1$ and copy its own target, so the
bidirectional arm could not use per-step teacher forcing at all and had to switch
to a leak-free \emph{last-item} objective: encode the whole prefix, predict only
the held-out next item. That reduces supervision from ${\sim}13$ targets per
session to \num{1} --- the same \num{73632} against \num{5723} arithmetic that is
the first mechanism behind the stacker, which is why the two worst designs on
record share a cause. The measured collapse (Recall@10 \num{0.0482}
$[\num{0.037},\num{0.059}]$, MRR \num{0.0225}, paired $\Delta$ versus the Markov
bar $-\num{0.0587}$ $[-\num{0.077},-\num{0.041}]$, the largest single-arm harm on
record) is therefore attributable to the \emph{objective}, not to
directionality, and the report declines to close bidirectionality: a fair test
would need a different leak-free scheme such as masked-item training. The
two-tower family excludes bidirectional mode by construction, with the leakage
argument written into its module docstring.

{\rightskip=0pt plus 2em \emph{Sources: \texttt{2026-\allowbreak 07-\allowbreak 13-\allowbreak nexttrack-\allowbreak topology-\allowbreak scan.md},
\texttt{2026-\allowbreak 07-\allowbreak 25-\allowbreak nexttrack-\allowbreak dual-\allowbreak tower-\allowbreak fusion-\allowbreak scan.md},
and \texttt{2026-\allowbreak 07-\allowbreak 26-\allowbreak nexttrack-\allowbreak pre-\allowbreak encoder-\allowbreak scan.md}.}\par}

\clearpage
\section{Power: at $n=\num{1431}$ the nulls are upper bounds, not proofs}
\label{sec:power}

\emph{Rationale.} Every null in Section~\ref{sec:failures} must be read against
the resolution of the instrument that produced it, or it becomes a false closure.

\emph{Results.} At this split's \num{1431} test sessions the paired-bootstrap CI
half-width is ${\sim}\num{0.0124}$ on Recall@10 (the mean of the nine paired
comparisons in the pre-encoder batch; \num{0.0128} for the nine in the
dual-tower batch) and ${\sim}\num{0.0018}$ on holisticness@10. Against that, an
entire architecture family spans very little: the ten-arm pre-encoder batch
spans \num{0.02796} in recall (\numrange{0.09713}{0.12509}) and \num{0.003147}
in grounded holisticness. Spread over resolution is therefore
$\num{0.02796}/\num{0.0124}\approx\num{2.3}$ on recall and
$\num{0.003147}/\num{0.0018}\approx\num{1.75}$ on the crown. Read that ratio as a
count of \emph{separable tiers}: if the whole family spans $2.3$ half-widths, then
at best about two groups of arms can be told apart, and about one on the crown ---
a heuristic, not a theorem, but one the direct count bears out and then some.
Those ten arms give nine adjacent gaps, and \emph{zero} of the nine adjacent
recall gaps and \emph{zero} of the nine adjacent crown gaps exceed their own
half-width: only the batch \emph{endpoints} separate at all. In the wider
dual-tower batch only two of nine adjacent recall gaps clear it. Resolving a
\num{0.005} recall difference --- a size chosen here as illustrative of a
difference one would want to act on, not as a pre-registered threshold --- would
need $\num{1431}\times(\num{0.0124}/\num{0.005})^2=\num{8801}$ test sessions,
$6.2\times$ the current split, or roughly \num{44000} sessions of listening
history at the split's \SI{20}{\percent} test fraction.
Figure~\ref{fig:nexttrack-power} draws the curve.

\emph{Interpretation --- and exactly which claims the caveat applies to.} The
defensible statement is that \textbf{large harms were detected reliably} --- the
pre-encoder's $-\num{0.0140}$, the fusion-head penalty's $+\num{0.0266}$,
depth-3's $-\num{0.0245}$, the bidirectional arm's $-\num{0.0587}$, the stacker's
\numrange{-0.038}{-0.028}, the degenerate centroid arm's $-\num{0.0566}$, all with
CIs excluding zero --- while \textbf{moderate wins were never detectable at all}.
Those two halves must be kept apart, because the phrase ``the nulls are
power-limited'' is easy to over-apply. Three of the five refuted axes produced
\emph{detections}, not nulls: the pre-encoder's harm, the fusion head's harm and
depth-3's harm all have intervals excluding zero, and nothing about the
instrument's resolution weakens them. What is power-limited is a different and
narrower claim: the \emph{absence of a win} on each axis. Every one of the five
axes ends in ``no arm cleared zero on the positive side'', and that sentence is an
upper bound on the available gain, not a demonstration of zero --- the
pre-encoder axis, for instance, is bounded at a tie ($+\num{0.0007}$ for its best
arm) and no more. The same distinction applies to the \emph{tie} verdicts, which
are the genuinely underpowered ones: the dual-tower family's best arm and the
width closure both rest on straddling intervals that a larger split could
resolve either way. It is \emph{not} ``these architectures do not work''.

\emph{The nuance that must travel with the caveat.} The
${\sim}\num{0.0124}$ figure is a \emph{cross-model} floor: it applies when two
arms have separately trained weights, because then almost every test session can
change its outcome. When two arms share \emph{bit-identical} trained weights and
differ only in an eval-time re-rank, the pairing is far tighter and the
instrument is \numrange{3}{33}$\times$ sharper --- across the eight MMR arms the
\emph{measured} half-widths are $\mathrm{hw}(\Delta H)$
\numrange{0.000382}{0.002476} and $\mathrm{hw}(\Delta\text{recall})$
\numrange{0.004193}{0.018169}. The reason is visible in $n_{\text{disc}}$, the
number of sessions whose hit@10 changed at all: as low as \num{9} at
$\lambda=\num{0.9}$ against \num{188} at $\lambda=\num{0.3}$, because
re-ranking the top of a candidate pool only flips sessions whose true item
crosses the $k=10$ boundary. That sharpness overturned a pre-registered call ---
one arm budgeted as ``likely unresolvable'' at \numrange{0.6}{0.8} conservative
half-widths in fact resolved at \num{4.8} \emph{measured} ones. Two riders:
measured half-widths \emph{grow} with distance from the control
(\num{0.000382} at $\lambda=\num{0.9}$ to \num{0.002476} at
$\lambda=\num{0.3}$), so the instrument is sharpest exactly where the effects
are smallest; and holisticness@10's absolute scale differs ${\sim}3\times$
between models (\num{0.0018} is \SI{21}{\percent} of the champion blend's $H$ but
only \SI{7}{\percent} of the single GRU's), so the crown instrument is
intrinsically ${\sim}3\times$ more precise on the GRU than on the champion. The
practical rule this record now follows is to budget the minimum detectable
effect from the \emph{pairing structure} of the comparison and to report the
measured per-arm half-width rather than a family constant.

{\rightskip=0pt plus 2em \emph{Sources: \texttt{2026-\allowbreak 07-\allowbreak 25-\allowbreak nexttrack-\allowbreak dual-\allowbreak tower-\allowbreak fusion-\allowbreak scan.md},
\texttt{2026-\allowbreak 07-\allowbreak 26-\allowbreak nexttrack-\allowbreak pre-\allowbreak encoder-\allowbreak scan.md},
and \texttt{2026-\allowbreak 07-\allowbreak 26-\allowbreak nexttrack-\allowbreak mmr-\allowbreak lambda-\allowbreak frontier.md}.}\par}

\begin{figure}[ht]
  \centering
  \includegraphics[width=0.8\textwidth]{nexttrack_power}
  \caption{The architecture nulls are power-limited, and this is the figure that
    reframes Section~\ref{sec:failures}. Paired-bootstrap CI half-width on
    Recall@10 against test-set size $n$ on a log axis, falling as
    $1/\sqrt{n}$ and anchored at the measured point $n=\num{1431}$,
    half-width \num{0.0124} (red dot). The dashed gray line at \num{0.028} is
    the \emph{entire} spread of an architecture family; the dotted green line at
    \num{0.005} is a difference one would want to be able to act on --- an
    illustrative choice, not a pre-registered threshold --- and it meets the curve
    only at $n\approx\num{8800}$ sessions, $6.2\times$ this split. The boxed
    callout gives spread over resolution as \num{2.3} on recall, from which the
    ``about two separable tiers'' follows as a heuristic
    ($\text{tiers}\approx\lfloor\text{spread}/\text{half-width}\rfloor$); the
    callout's ``eight adjacent arms'' is loose --- the batch has ten arms and
    therefore nine adjacent gaps, none of which clears its own half-width. Read
    with Section~\ref{sec:power}: this floor applies to \emph{separately trained}
    arms; shared-weight eval-time comparisons such as Table~\ref{tab:mmr}'s are
    \numrange{3}{33}$\times$ sharper. \emph{Sources:
    \texttt{2026-\allowbreak 07-\allowbreak 25-\allowbreak nexttrack-\allowbreak dual-\allowbreak tower-\allowbreak fusion-\allowbreak scan.md},
    and \texttt{2026-\allowbreak 07-\allowbreak 26-\allowbreak nexttrack-\allowbreak pre-\allowbreak encoder-\allowbreak scan.md}.}}
  \label{fig:nexttrack-power}
\end{figure}

\clearpage
\section{Where the gains actually came from}
\label{sec:gains}

Read as one arc, the record says something narrower and more useful than ``the
blend is best''. Four claims survive.

\emph{Gains came from adding a signal of a different kind, never from more
model.} The board moved three times: when a co-occurrence lookup was fused with
a sequence model (paired $\Delta$ $+\num{0.0657}$ over the bar, on the earlier
\dsid{seq-20260713-014226} artifact where that step was taken), when a content
retrieval leg was added to that (${+}\num{0.0147}$, CI above zero), and when the
metric those two content-shaped legs retrieve in was \emph{learned}
(${+}\num{0.0259}$, CI above zero, hardened to ${+}\num{0.0426}$ on a colder
split). It never moved for width, depth, direction, an extra tower, an extra
nonlinearity, or a learned combiner. The mechanism is the cold half: a new
\emph{kind} of signal can reach sessions the existing ones cannot, whereas more
capacity only re-describes what is already reachable.

\emph{The winning path is linear where it touches the retrieval geometry.} The
loss and the serving metric both live in a fixed external space --- cosine
retrieval over signed \num{192}-dimensional latents --- so the readout should be
the most general map \emph{into} that space and nothing more. Every measured
insertion of a half-wave rectifier into that path lost: on the input side
(${-}\num{0.0140}$, with the ReLU isolated as the whole cause at
${\sim}\num{5.9}\sigma$ of decodability), on the output side (three
confirmations across three regimes, at \SI{10}{\percent} more parameters), and in
the pooled ANN twin (${-}\num{0.0245}$, same space, PCA-128). The single largest win, meanwhile, is a bias-free
linear map. ``Add a head'' is a reflex worth resisting whenever the target space
is fixed.

\emph{The combination should not be learned at this data volume, and that is
measurable in one run.} An oracle weighting fitted on the test labels buys
$+\num{0.0042}$ with a CI straddling zero. That single probe bounds the entire
learned-combiner lane, and it should be the \emph{first} experiment run in such a
lane, not the last. The transferable discipline is to count the supervision
events available to the combiner separately from those available to the parts,
and to check whether the candidate distribution at fit time matches the one at
inference.

\emph{The metrics are part of the result.} A composite built from list properties
rather than correctness is maximized by a model that has stopped trying, and the
$\pm\num{0.015}$ recall floor was the only thing standing between this record and
a crown for a model \SI{46}{\percent} worse at its actual job. A grounding factor
helps and is not sufficient. The counterexample that makes the metric salvageable
rather than useless is MMR, which raises three factors \emph{including} the
grounding one, at a recall tie --- same metric, opposite mechanism.

\clearpage
\section{What is not established}
\label{sec:notestablished}

This section exists so that no reader has to reconstruct the caveats from the
prose.

\textbf{Nothing from 2026-07-26 is promoted, and the champion is unchanged.} Six
MMR definitions were \emph{registered}; the auto-generated
\texttt{best-}$\langle$predictor$\rangle$-$\langle$run$\rangle$ records that
exist for those runs are best-models group membership, not deliberate promotion.
Of \num{77} model records only \emph{five} are deliberate promotions, and the
recall champion remains \Rp+M+\Cp{} at \num{0.21174}. The best-models group
currently holds nine of twelve slots, all nine marked \texttt{source: auto}. Eight
of the nine come from the 2026-07-26 batches and the ninth (the two-leg R+M blend)
from 2026-07-22, and two of them --- the TIER-2 \emph{priced} arms B3 and A2 ---
have a $\Delta$ recall CI below zero. The group is nonetheless \emph{curated}, not
an uncurated flood: \texttt{GET /api/best-models} returns per-entry
\texttt{rank}, \texttt{metric}, \texttt{metric\_value}, \texttt{source} and
\texttt{selected\_at}, plus a top-level list of \num{3} \texttt{pinned} model
names and \num{75} \texttt{excluded} identifiers --- and the two TIER-3
\emph{rejected} MMR arms are explicitly among the exclusions. What remains
unverified is only whether the current membership reflects a deliberate curation
pass or the server's automatic recompute; \texttt{PROJECT-FACTS.md} records that
remote-worker completions do not fire that recompute, so the group can be silently
stale after any offloaded batch and its \texttt{updated\_at} must be checked
against the newest run's \texttt{finished\_at} before it is trusted.

\textbf{Several nulls are underpowered rather than refuted.} See
Section~\ref{sec:power}. In particular the width closure rests on endpoints
tying with a jitter of ${\sim}\pm\num{0.01}$ between neighbouring widths, one
control losing outright, and a prescribed two-run replication that has not been
bought; and the pre-encoder result carries an unresolved dropout confound.

\textbf{Several comparisons are nominal, single-run, or report-provenance only.}
No paired CI exists for the temperature cells, the projection-rank ablation, the
BPR cells, the anti-eager $\beta$ arms, or the canonical-split LSTM --- the
reports call the rank-64 and LSTM gaps ties and the BPR gap a collapse, but none
is backed by a paired bootstrap. The champion's \Rp{} and \Cp{} standalone
figures (\num{0.1726}, \num{0.1516}) exist only as off-server-driver numbers in
a campaign report, with no server run and no raw hit counts; only the Markov leg
has a real canonical-split run. The warm/cold decompositions, the hit-set overlap
counts, and the candidate-geometry measurement of
Figure~\ref{fig:mc-candidate-eigen} are single-run observations with no
uncertainty quantified. Two $\beta$ arms executed on a stale worker checkout and
are not code-comparable to any current batch.

\textbf{The crown metric is young and has already been revised twice}, and its
grounding floor \num{0.34} is currently operating \emph{outside} the regime it
was calibrated for: it was set to cancel gains bought by destroying recall, and
MMR is the first lever that moves music@10 in the opposite direction. Stored
holisticness values are comparable only within a generation
(Table~\ref{tab:hgen}), the generation labels used here were re-derived from
stored per-row facets rather than read from the API, and two runs store gen-3
values whose re-derivation is off by ${\sim}\num{1e-3}$ because their
\texttt{artist\_conc\_at\_k} column is missing from the metrics index --- their
$H$ is trustworthy, their facet row is not, and their complete twins were quoted
instead.

\textbf{\texttt{seq-embed} has no numbers at all}, \texttt{seq-ann} has no
canonical-split run, and no parameter count is recorded anywhere for
\texttt{seq-stacker} or \texttt{seq-ann} (the stacker's is not even well defined
--- \num{200} boosted trees at depth 6 plus a base GRU). Three parameter counts
in Table~\ref{tab:models} were derived by instantiating the real classes at
$\text{latent\_dim}=192$ rather than transcribed: the LSTM's \num{510144} and the
champion's \num{431808} ($=\num{394944}+\num{36864}$, the GRU plus the shared
projection $W$).

\clearpage
\section{Conclusions}
\label{sec:conclusions}

Six trained predictor families on one leak-free next-track split converge on a
clear story. (1) The task's defining property is that \num{788} of \num{1431}
test targets are \emph{cold}, capping every seen-items-only mechanism at
Recall@10 \num{0.449336} and making retrieval in a shared content geometry the
only route to the majority of the test set. (2) The champion is a fixed
score-level blend, \Rp+M+\Cp{} at \num{0.21174} (\num{303} of \num{1431}), and it
wins because its legs are correct on almost \emph{disjoint} sessions --- the
Markov leg reads warm \num{0.2255} / cold \num{0.0102}, the GRU leg is flat
across the split, and Jaccard overlap between their hit sets is \num{0.201}. (3)
Fusion of unlike signals pays and fusion of like signals does not: a second
neural leg contributes nothing, a session-kNN leg is actively harmful
($-\num{0.038}$, CI below zero), and adding a leg only helps when it reaches
sessions the others cannot. (4) The single largest win is a learned InfoNCE
projection of the item space, ${+}\num{0.0259}$
$[+\num{0.012},+\num{0.041}]$ over the unprojected crown, seed-robust across
three seeds and \emph{wider} on a colder second split
(${+}\num{0.0426}$ $[+\num{0.0280},+\num{0.0573}]$); it lifts both legs that
share the geometry (content-kNN $2.2\times$, with \emph{zero} trainable
parameters of its own) and leaves the ID-based leg untouched, and full rank ties
rank 64 while BPR collapses \num{0.055} below InfoNCE --- so the gain is the
\emph{supervision}, not the capacity. (5) Recurrence beats pooling
($-\num{0.0245}$ $[-\num{0.039},-\num{0.011}]$ for the pooled twin, same space,
byte-identical objective), and the damage concentrates on the cold half, where
the quality of the predicted point is all there is. (6) Eval-time MMR re-ranking
is the first honest crown lever --- eight of eight arms lift holisticness@10 with
paired CI above zero, \emph{four} of them at a recall tie (three on the plain GRU,
one on the champion), and it de-concentrates artists
\numrange{1.59}{3.06}$\times$ faster than it reduces seed echo with no notion of
artist in the rule --- but it benefits the plain GRU ($H$ \num{0.0342}) far more
than the champion (\num{0.0153}), whose mood coherence falls twice as fast when
diversified at matched $\lambda$ ($-\SI{8.5}{\percent}$ against
$-\SI{4.2}{\percent}$) and collapses $-\SI{40}{\percent}$ when pushed to
$\lambda=\num{0.3}$, and \textbf{nothing from that campaign is promoted}. (7) Five design axes are refuted with measured mechanisms: a second
tower's input is measurable with respect to the first's, so it can add ensemble
diversity but no information --- and joint training destroys exactly that
diversity; a ReLU pre-encoder is provably expressivity-neutral in its linear part
and deletes \num{0.5327} of a zero-centred latent in its nonlinear one, leaving
\num{119.6} live units for \num{192} signed directions; an MLP fusion head is
harmful three times over --- reproduced in three different regimes, though the
apparent ordering among the three is nominal --- at \SI{10}{\percent} more
parameters; four learned combiners land at
or below their own base legs while an oracle weighting fitted on the test labels
buys only $+\num{0.0042}$, CI straddling zero; and depth, width and direction are
spent, with the bidirectional collapse attributable to the
$12.87\times$ supervision cut its objective forced, not to directionality. (8)
The record carries two primary metrics and must not silently pick one: Recall@10
is the comparability currency, holisticness@10 is the server's declared primary,
the two orderings are near-inverted, and the crown has had \emph{three}
incompatible generations whose values differ by up to $6.4\times$ on identical
predictions. (9) Under the live crown the ten separately-trained architecture arms occupy a
band \numrange{0.02127}{0.02487} whose whole width (${\sim}\num{0.0036}$) is
barely twice $\mathrm{hw}(\Delta H)\approx\num{0.0018}$, so their rank order
carries no information --- while the recall champion, at \num{0.008603} far below
that band, and the MMR arms, at \numrange{0.026761}{0.034171} above it, are
\emph{outside} it and separated for real. (10) What is power-limited on all five refuted axes is the
\emph{absence of a win}, not the harms: at $n=\num{1431}$ the paired half-width is
${\sim}\num{0.0124}$ on recall against a whole-family spread of \num{0.028}, zero
of nine adjacent gaps in the last batch exceed their own half-width, and resolving
\num{0.005} would need \num{8801} test sessions --- so large harms were detected
reliably with intervals excluding zero, moderate wins were never detectable, and
each axis's closure bounds an available gain rather than demonstrating zero. The
genuinely underpowered verdicts are the \emph{ties}: the dual-tower family's best
arm and the single-GRU width closure. (11) Shared-weight eval-time comparisons are a different
instrument entirely, \numrange{3}{33}$\times$ sharper because $n_{\text{disc}}$
can be as low as \num{9} of \num{1431}, so the minimum detectable effect must be
budgeted from the pairing structure and not from a family constant. (12) What
actually produced every gain on this board was a new \emph{kind} of signal or a
better shared \emph{geometry}, never more parameters, more towers, more
nonlinearity, or a learned combination --- which is also why
\texttt{seq-embed}'s per-item table is the one untried family the data budget
already argues against, at a median of one observation per observed item.

\clearpage
\section{Source reports}
\label{sec:sources}

This comparison consolidates fifteen next-track campaign reports from
\texttt{experiments/}, plus the living facts file:

\begin{itemize}
  \raggedright
  \item \texttt{2026-\allowbreak 07-\allowbreak 13-\allowbreak nexttrack-\allowbreak sequence-\allowbreak sweep.md} --- establishes the
  session GRU with in-batch InfoNCE as the single-model workhorse and the
  first-order Markov lookup as the reference bar (Section~\ref{sec:board}).
  \item \texttt{2026-\allowbreak 07-\allowbreak 13-\allowbreak nexttrack-\allowbreak topology-\allowbreak scan.md} --- depth, width and
  direction closed on the early item space: depth-3 \textbf{loses} to the bar
  ($-\num{0.0245}$, CI below zero), residual was never isolated at matched width,
  and the bidirectional arm's collapse to \num{0.0482} is an \emph{objective}
  artifact
  (Section~\ref{sec:fail-topology}).
  \item \texttt{2026-\allowbreak 07-\allowbreak 14-\allowbreak nexttrack-\allowbreak blend-\allowbreak and-\allowbreak ensemble.md} --- the fixed
  z-blend of a sequence model with a co-occurrence lookup, $+\num{0.0657}$
  $[+\num{0.049},+\num{0.083}]$ over the bar, three-seed confirmed; a second
  neural leg is redundant (Section~\ref{sec:fusion}).
  \item \texttt{2026-\allowbreak 07-\allowbreak 14-\allowbreak nexttrack-\allowbreak graded-\allowbreak relevance-\allowbreak and-\allowbreak new-\allowbreak legs.md} ---
  introduces artist@$k$ / genre@$k$ / artist\_mrr as first-class graded metrics
  (Section~\ref{sec:metrics}).
  \item \texttt{2026-\allowbreak 07-\allowbreak 15-\allowbreak nexttrack-\allowbreak compression-\allowbreak representation-\allowbreak bakeoff.md}
  --- the compression/representation bake-off.
  \item \texttt{2026-\allowbreak 07-\allowbreak 15b-\allowbreak nexttrack-\allowbreak compression-\allowbreak dim-\allowbreak extension.md} --- with the
  dimension extension, these two fix PCA-192 as the shared item space every
  trained model reads (Section~\ref{sec:intro}).
  \item \texttt{2026-\allowbreak 07-\allowbreak 15c-\allowbreak nexttrack-\allowbreak phase2-\allowbreak model-\allowbreak sweep.md} --- the family
  sweep: the pooled ANN twin loses to the same-space GRU
  ($-\num{0.0245}$, CI below zero) and the XGBoost stacker lands below its own
  base leg (Sections~\ref{sec:recurrence} and~\ref{sec:fail-combiners}).
  \item \texttt{2026-\allowbreak 07-\allowbreak 18-\allowbreak nexttrack-\allowbreak literature-\allowbreak fit-\allowbreak campaign.md} --- the learned
  InfoNCE projection champion, $+\num{0.0259}$ over the prior crown, plus the
  oracle-weight probe that bounds the learned-combiner lane at $+\num{0.0042}$,
  CI straddling zero (Sections~\ref{sec:projection}
  and~\ref{sec:fail-combiners}).
  \item \texttt{2026-\allowbreak 07-\allowbreak 18b-\allowbreak nexttrack-\allowbreak projection-\allowbreak registration-\allowbreak and-\allowbreak scan.md} ---
  registers the champion and scans its axes: InfoNCE $\gg$ BPR, full rank
  $\approx$ rank 64, $\tau$ mildly monotone with the default \textbf{held}
  (Section~\ref{sec:projection}).
  \item \texttt{2026-\allowbreak 07-\allowbreak 18c-\allowbreak nexttrack-\allowbreak representation-\allowbreak exploration.md} --- six of
  eight alternative item spaces lift the \emph{frozen} readout with CI above
  zero and \emph{none} lifts the champion: whitening and the learned projection
  are the same fix in two places (Section~\ref{sec:projection}).
  \item \texttt{2026-\allowbreak 07-\allowbreak 20-\allowbreak musical-\allowbreak distance-\allowbreak metric.md} --- music@10, the graded
  grounding metric that later made the crown salvageable
  (Section~\ref{sec:metrics}).
  \item \texttt{2026-\allowbreak 07-\allowbreak 22-\allowbreak session-\allowbreak holisticness-\allowbreak and-\allowbreak antieager-\allowbreak levers.md} --- the
  holisticness suite, the weak anti-eager training knob, and the offline
  $\lambda$ sweep that transferred to the native arms to four decimals
  (Sections~\ref{sec:metrics} and~\ref{sec:mmr}).
  \item \texttt{2026-\allowbreak 07-\allowbreak 25-\allowbreak nexttrack-\allowbreak dual-\allowbreak tower-\allowbreak fusion-\allowbreak scan.md} --- dual-tower
  fusion \textbf{REFUTED}, the fourth combiner null: zero of eight arms clear
  zero against a single-GRU width control param-matched to \SI{0.098}{\percent} of
  the $f1$ arms, the MLP fusion
  head is harmful twice over, and five $\Delta H$ ``wins'' all fail the recall
  floor (Sections~\ref{sec:fail-views}, \ref{sec:fail-head}
  and~\ref{sec:metrics}).
  \item \texttt{2026-\allowbreak 07-\allowbreak 26-\allowbreak nexttrack-\allowbreak pre-\allowbreak encoder-\allowbreak scan.md} --- the per-step MLP
  pre-encoder \textbf{REFUTED and actively harmful} ($-\num{0.01398}$, CI below
  zero), with the rectifier isolated as the whole mechanism at
  ${\sim}\num{5.9}\sigma$ of genre decodability, and the first batch in the
  project's history to produce \emph{zero} crown candidates
  (Section~\ref{sec:fail-relu}).
  \item \texttt{2026-\allowbreak 07-\allowbreak 26-\allowbreak nexttrack-\allowbreak mmr-\allowbreak lambda-\allowbreak frontier.md} --- the MMR
  $\lambda$ frontier: eight of eight arms lift holisticness@10 with paired CI
  above zero, four free of any recall cost (three on the GRU, one on the
  champion), six definitions
  \textbf{registered and none promoted} (Section~\ref{sec:mmr}).
  \item \texttt{PROJECT-FACTS.md} --- the living facts file: the noise bands,
  the power anchors, the promotion ledger and the metric-generation history
  (Sections~\ref{sec:power} and~\ref{sec:notestablished}).
\end{itemize}

\end{document}
